\documentclass[11pt,oneside]{book}

% BEGIN preamble
\usepackage[T1]{fontenc}
\usepackage{lmodern,amsmath,amssymb,amsthm,mathtools}
\usepackage[a4paper,margin=25mm,headheight=15pt]{geometry}
\usepackage{microtype,enumitem,xcolor,longtable,booktabs,array,needspace}
\usepackage[colorlinks=true,linkcolor=blue!40!black,urlcolor=blue!40!black]{hyperref}
\usepackage[nameinlink,noabbrev]{cleveref}
\usepackage{fancyhdr}
\hypersetup{pdftitle={Ordinary Differential Equations and Dynamical Systems: Unified Notes},pdfauthor={MathGenome preparation series}}
\pagestyle{fancy}\fancyhf{}
\fancyhead[L]{\small MathGenome: ODE and Dynamics}\fancyhead[R]{\small\thepage}
\renewcommand{\headrulewidth}{0.3pt}
\setlength{\parindent}{0pt}\setlength{\parskip}{5pt}
\setlength{\emergencystretch}{3em}\allowdisplaybreaks[2]\raggedbottom
\setcounter{tocdepth}{1}
\newtheorem{theorem}{Theorem}[chapter]
\newtheorem{lemma}[theorem]{Lemma}
\newtheorem{proposition}[theorem]{Proposition}
\newtheorem{corollary}[theorem]{Corollary}
\theoremstyle{definition}
\newtheorem{definition}[theorem]{Definition}
\newtheorem{convention}[theorem]{Convention}
\newtheorem{construction}[theorem]{Construction}
\newtheorem{example}[theorem]{Example}
\theoremstyle{remark}\newtheorem{remark}[theorem]{Remark}
\newcommand{\R}{\mathbb R}\newcommand{\C}{\mathbb C}
\newcommand{\N}{\mathbb N}\newcommand{\Z}{\mathbb Z}\newcommand{\Q}{\mathbb Q}
\newcommand{\dd}{\,d}\newcommand{\id}{\operatorname{id}}
\newcommand{\abs}[1]{\lvert#1\rvert}\newcommand{\norm}[1]{\lVert#1\rVert}
\newcommand{\ip}[2]{\langle#1,#2\rangle}
\DeclareMathOperator{\supp}{supp}\DeclareMathOperator{\dist}{dist}
\DeclareMathOperator{\Lip}{Lip}\DeclareMathOperator{\Span}{span}
\DeclareMathOperator{\Leb}{Leb}\DeclareMathOperator{\diam}{diam}
\newcommand{\source}[1]{\par\begingroup\small\sffamily\color{black!65}Source comparison: #1\par\endgroup}

% END preamble

\begin{document}
\frontmatter

% BEGIN frontmatter
\begin{titlepage}
\centering
\vspace*{25mm}
{\Huge Ordinary Differential Equations\par}
\vspace{4mm}
{\Huge and Dynamical Systems\par}
\vspace{12mm}
{\Large Unified undergraduate and honours notes\par}
\vspace{12mm}
{\large MathGenome preparation series --- Direction 13\par}
\vfill
{\large English synthesis with explicit proofs and foundation bridges\par}
\vspace{6mm}
{\large Prepared 4 November 2026\par}
\end{titlepage}

\chapter*{Purpose, coverage, and conventions}
These notes assemble a common mathematical core from six complementary
references. The structure starts from nineteen curriculum modules and adds
two necessary foundation chapters: local hyperbolic theory and regular planar
separation. Repetition of a theorem across textbooks does not create a second
mathematical obligation. The scope is determined by the precise adopted
statements below, not by a claim to contain every theorem or exercise in all
six reference books.

The classical part covers elementary reductions, existence and uniqueness,
maximal solutions, quantitative dependence, regularity of solution maps,
linear and periodic systems, convergent series, and regular boundary problems.
The dynamical part covers partial flows, stability, planar limit sets, return
maps, discrete dynamics, selected generic bifurcations, symbolic models,
explicit exponents and dimensions, and introductory ergodic theory.
Later chapters supplying a foundation may be read before the chapter that
uses it; the cross references identify those dependencies.

Unless specified otherwise, a vector field is defined on an open subset of a
finite-dimensional real space, and norms are fixed but arbitrary equivalent
norms. A solution includes its interval, regularity, equation and initial data.
An autonomous flow is in general partial; its group law is asserted only
where both sides exist. A symbol such as $C^k$ always applies to the variables
explicitly stated. Continuous, differentiable, smooth and real-analytic
dependence are distinct assertions. Complex analytic statements use
holomorphy on an actual open complex domain.

All adopted theorems, lemmas, propositions, corollaries and required definitions
are part of the mathematical specification. Another valid proof, or an exact
application of a correctly matched library theorem, may replace the printed
proof. Steps used only by one proof route do not become additional independent
acceptance obligations. Conditions, domains, quantifiers, quantitative bounds,
definitions and necessary representation equivalences remain binding.
Examples and counterexamples make boundary cases explicit; their mathematical
claims are not permission to infer stronger general statements.

This volume deliberately bounds graduate extensions. It does not adopt a full
singular-endpoint spectral theory, general coupled center-manifold theorem,
high-codimension global bifurcation classification, KAM theorem, general
multiplicative ergodic theorem, or unrestricted dimension formula. Where a
selected form is proved, its exact scope is stated. Basic analysis and topology
already established in the previous seven MathGenome components may be used
through the following foundation register.

% END frontmatter

\tableofcontents

% BEGIN foundations
\chapter*{Previously established foundations}
\addcontentsline{toc}{chapter}{Previously established foundations}
The following are prior mathematical prerequisites, with their hypotheses
retained. Their Lean representations are checked against the fixed
MathGenome/Mathlib environment in the accompanying interface crosswalk.
These references are not assumptions that replace new ODE theorems.

\begin{description}[style=nextline,leftmargin=11mm]
\item[F1. Finite-dimensional linear algebra.]
Over $\R$ and $\C$: linear maps and matrices, determinants, Jordan reduction
over a splitting field, invariant subspaces, spectral decompositions for
self-adjoint finite-dimensional operators, and equivalence of norms.
These are available from D05, D07 and the fixed Mathlib.
The infinite-dimensional compact spectral argument needed for boundary
problems is supplied in this volume rather than inferred from the
finite-dimensional spectral theorem.

\item[F2. One-variable analysis.]
Completeness of $\R$, uniform convergence, scalar and vector-valued Riemann
integration of continuous functions on compact intervals, the fundamental
theorem of calculus, differentiation under a uniformly controlled integral,
Taylor's theorem with its remainder, and convergent power series.
Use D03 and the compatible scalar/finite-vector bridges in D04 and D11.

\item[F3. Multivariable calculus.]
Fr\'echet derivatives in finite dimensions, the chain rule, higher derivative
symmetry, the local inverse and implicit function theorems, compactly
supported $C^k$ cutoffs in Euclidean open sets, the Jacobian change-of-variables
formula for $C^1$ diffeomorphisms on their actual domains, and Green's formula for an
actual bounded regular planar domain with its oriented boundary.
These are D04 prerequisites. A generic planar separation theorem is
\emph{not} assumed from them; Chapter~\ref{ch:d13-21} supplies the required
regular-curve case.

\item[F4. Metric and function-space topology.]
Compactness and sequential compactness for compact metric spaces,
Heine--Borel in finite dimensions, uniform continuity on compact sets,
completeness of bounded continuous functions with the supremum norm,
the contraction principle on a complete metric space, and Arzel\`a--Ascoli
on compact domains. Uniform limits of continuous functions are continuous.
These are D10 foundations. Weighted spaces in this volume are explicitly
identified with ordinary complete sup-norm spaces.

\item[F5. Measure and integration.]
On complete or explicitly completed finite/probability spaces: monotone,
Fatou and dominated convergence, Tonelli--Fubini with their hypotheses,
H\"older and Minkowski inequalities, completeness of $L^p$ for $1\leq p\leq
\infty$, density of bounded functions in $L^p$ for finite $p$, and the first
Borel--Cantelli lemma when the sum of measures is finite, the
Carath\'eodory extension of a countably additive premeasure with the finite
measure uniqueness theorem on a generating algebra, and regularity of
Lebesgue measure on compact Euclidean domains. On the circle,
continuous functions are dense in $L^2$ of Lebesgue measure. Use D11. Equalities
between $L^p$ elements are almost-everywhere equalities until representatives
and null sets are fixed.

\item[F6. Complex analysis.]
Locally uniform holomorphic limits, the Cauchy integral formula and estimates,
Taylor/Laurent expansions on their actual domains, and winding index
invariance under a homotopy avoiding the evaluation point. Local index
jumps across regular arcs and Green/index bridges retain their path and
boundary hypotheses. Use D12. Real smoothness is not substituted for any
holomorphic hypothesis. The uniform Fej\'er approximation of continuous
periodic functions is a D03 Fourier foundation; together with the preceding
$L^2$ density it justifies the circle-rotation example in Chapter~\ref{ch:d13-19}.

\item[F7. Positive Riesz--Markov representation.]
For a locally compact Hausdorff space $X$, every positive real linear
functional on $C_c(X,\R)$ has a unique positive Radon representing measure
on the Borel sets; conversely integration against such a measure gives that
functional. On compact $X$, $C_c(X)=C(X)$ and its total mass is the functional
at $1$. The actual prior interface is
\texttt{D11.Ch19.PositiveRieszMarkov.real\_representation}, exposed through
\texttt{MathGenome.RealAnalysis.Topics.D11.Ch19}. The finite-dimensional
Riesz representation theorem alone would not justify this invocation.
\end{description}

% END foundations

\mainmatter

% BEGIN chapters/M13.01
\chapter{Equations, solution concepts, and changes of formulation}\label{ch:d13-01}

\section{Objects and domains}
An ordinary differential equation includes a domain, not just a formula.
Throughout the real theory, $U\subset\R\times\R^n$ is open and
$f:U\to\R^n$ is continuous. Coordinates in $U$ are written $(t,x)$.
The norm is any fixed norm on $\R^n$; operator norms always use that norm.
Finite-dimensional equivalence of norms permits a change of norm without
changing existence, continuity, or differentiability assertions.

\begin{definition}[Classical solution and initial value problem]\label{d13:01:solution}
A classical solution of $x'=f(t,x)$ on a nonempty open interval $I$ is a
$C^1$ map $x:I\to\R^n$ with $(t,x(t))\in U$ and
$x'(t)=f(t,x(t))$ for every $t\in I$. It solves the initial value problem
at $(s,\xi)\in U$ if $s\in I$ and $x(s)=\xi$.
On a compact interval we require continuity on the whole interval and the
corresponding one-sided derivative at an endpoint; equivalently, we use the
integral formulation below. A maximal solution is one admitting no extension
to a strictly larger open interval. A global solution is defined for all
$t\in\R$.
\end{definition}

\begin{theorem}[Differential and integral formulations]\label{d13:01:integral}
Let $I$ be an interval containing $s$, and let $x:I\to\R^n$ be continuous
with graph in $U$. Then $x$ solves the initial value problem with value
$\xi$ if and only if
\[
 x(t)=\xi+\int_s^t f(r,x(r))\,\dd r\qquad(t\in I).
\]
The integral is oriented, so the same assertion applies for $t<s$.
\end{theorem}
\begin{proof}
For a classical solution, the fundamental theorem of calculus applied to
each coordinate gives the identity and $x(s)=\xi$. Conversely,
$r\mapsto f(r,x(r))$ is continuous. Its integral has derivative
$f(t,x(t))$ in the interior and the stated one-sided derivatives at any
included endpoints. The identity therefore gives all the required
regularity and the initial value. The argument for $t<s$ uses
$\int_s^t=-\int_t^s$, with no change to the derivative formula.
\end{proof}

\begin{proposition}[Restriction and compatible pasting]\label{d13:01:pasting}
Restriction of a solution to a subinterval is a solution. If two solutions
on intervals $I_1,I_2$ agree on their nonempty intersection, they define a
solution on $I_1\cup I_2$. A left solution and a right solution that are
continuous up to a common time $s$, have the same value $\xi$ there, and
satisfy the equation up to their respective sides can also be pasted.
\end{proposition}
\begin{proof}
Restriction preserves both the graph and derivative identities. For
overlapping open intervals, every point of the union has a neighborhood
where the pasted function is one of the given $C^1$ functions. At a
joining endpoint, continuity of $f$ shows that both one-sided derivatives
converge to $f(s,\xi)$. Alternatively, add the two integral identities at
$s$ and apply Theorem~\ref{d13:01:integral}. Equality on an overlap is a
hypothesis here; it will follow from uniqueness under the conditions of
Chapter~\ref{ch:d13-03}.
\end{proof}

\begin{definition}[Autonomous and parameter-dependent equations]\label{d13:01:autonomous}
An autonomous equation has the form $x'=F(x)$ on an open
$\Omega\subset\R^n$. A parameter-dependent equation is
$x'=f(t,x,p)$, where $p$ belongs to an open subset of $\R^m$ and remains
constant along a solution. A parameter is not an additional freely chosen
function of time. We write $\Phi(t,s,\xi,p)$ for the unique solution value
when uniqueness has been established.
\end{definition}

\section{Equivalent systems}
\begin{theorem}[Reduction of a higher-order equation]\label{d13:01:higher-order}
For $m\geq1$, the equation
\[
 y^{(m)}=G(t,y,y',\ldots,y^{(m-1)})
\]
with continuous $G$ on an open subset of $\R\times(\R^d)^m$ is equivalent
to the first-order system on that same time interval
\[
 Y'=(Y_2,Y_3,\ldots,Y_m,G(t,Y_1,\ldots,Y_m)),
 \qquad Y=(y,y',\ldots,y^{(m-1)}).
\]
Prescribing the $m$ derivatives of $y$ at $s$ is equivalent to prescribing
$Y(s)$.
\end{theorem}
\begin{proof}
A $C^m$ solution $y$ gives a $C^1$ solution $Y$ by differentiating its
coordinates. Conversely, if $Y$ is $C^1$, the equations
$Y_1'=Y_2,\ldots,Y_{m-1}'=Y_m$ imply successively that
$Y_1\in C^m$ and $Y_j=Y_1^{(j-1)}$. The last component equation is the
original equation. Every domain and initial-value condition is preserved,
so this is a bijection of solution sets, not merely a method of constructing
some solutions.
\end{proof}

\begin{proposition}[Autonomization]\label{d13:01:autonomization}
The nonautonomous equation $x'=f(t,x,p)$ is equivalent to the autonomous
system on the open set of allowed $(\tau,x,p)$,
\[
 (\tau,x,p)'=(1,f(\tau,x,p),0).
\]
The solution starting at $(s,\xi,p)$ at elapsed time $0$ has
$\tau(h)=s+h$, constant parameter $p$, and
$x(h)=\Phi(s+h,s,\xi,p)$ whenever either side is defined.
\end{proposition}
\begin{proof}
The first and last equations integrate to the asserted clock and parameter.
Substitution into the middle equation and the change $t=s+h$ give the
original equation. Reversing these operations gives the inverse
correspondence, with exactly the translated solution interval.
\end{proof}

\begin{proposition}[Coordinate and time changes]\label{d13:01:coordinates}
If $H:\Omega\to\widetilde\Omega$ is a $C^1$ diffeomorphism, then
$z=H(x)$ transforms $x'=f(t,x)$ into
\[
 z'=DH(H^{-1}(z))f(t,H^{-1}(z)).
\]
If $\theta:J\to I$ is a $C^1$ diffeomorphism of intervals with nonzero
derivative, $y(r)=x(\theta(r))$ satisfies
$y'(r)=\theta'(r)f(\theta(r),y(r))$. In particular, time reversal changes
$f(t,x)$ into $-f(-t,x)$ after the choice of time origin.
\end{proposition}
\begin{proof}
Both displayed equations are the chain rule. Applying $H^{-1}$ or
$\theta^{-1}$ proves the converse, including the domain and initial data.
The assumption on $\theta'$ is what makes the time change invertible;
an arbitrary reparameterized curve need not have the original timing.
\end{proof}

\section{Initial versus boundary data}
\begin{example}[Boundary values need not determine a solution]\label{d13:01:boundary-example}
Every solution of $y''+y=0$ has the form $A\cos t+B\sin t$, as follows
either by differentiation and initial-value uniqueness or by the
oscillator calculation in Chapter~\ref{ch:d13-02}. The conditions
$y(0)=y(\pi)=0$ leave $B$ arbitrary. The conditions $y(0)=0$,
$y(\pi)=1$ have no solution. Thus existence and uniqueness for initial
values do not settle a two-point boundary problem. Boundary conditions
will be part of the operator domain in Chapter~\ref{ch:d13-10}.
\end{example}

\paragraph{Sources and coverage.}
This chapter adopts the solution and reduction conventions of [T,
Chapter~1 and Section~2.1], with the domain bookkeeping made explicit;
[CL, Chapter~1, Sections~5--6] and [HSD, Chapters~1 and~17] give compatible
system formulations. Implicit equations $F(t,x,x')=0$ are converted to
normal form only where the relevant implicit-function hypothesis holds;
no global normal form is assumed for a singular implicit equation.

% END chapters/M13.01


% BEGIN chapters/M13.02
\chapter{Explicit equations and elementary models}\label{ch:d13-02}

\section{Linear and separable equations}
\begin{theorem}[Scalar variation of constants]\label{d13:02:linear}
Let $a,b:I\to\mathbb K$ be continuous, where $\mathbb K=\R$ or $\C$,
and $s\in I$. Write $A(t)=\int_s^t a(r)\,\dd r$. The unique solution
of $x'=a(t)x+b(t)$, $x(s)=\xi$, on $I$ is
\[
 x(t)=e^{A(t)}\left(\xi+\int_s^t e^{-A(r)}b(r)\,\dd r\right).
\]
\end{theorem}
\begin{proof}
The product rule gives $(e^{-A}x)'=e^{-A}b$ for any solution.
Integration forces the formula and hence uniqueness. Conversely,
differentiating that formula gives the equation and initial value.
The integrating factor never vanishes, so multiplying by it loses no
solutions. The proof works over either stated scalar field.
\end{proof}

\begin{theorem}[Separation on a nonvanishing branch]\label{d13:02:separable}
Let $a:I\to\R$ and $g:J\to\R$ be continuous, where $I,J$ are open
intervals and $g$ is nowhere zero on $J$. Fix $s\in I$ and $\xi\in J$,
and define $H(u)=\int_\xi^u 1/g(v)\,\dd v$. Then $H$ is strictly
monotone with a $C^1$ inverse on $H(J)$. The solution of
$x'=a(t)g(x)$, $x(s)=\xi$, taking values in $J$ is
\[
 x(t)=H^{-1}\left(\int_s^t a(r)\,\dd r\right)
\]
on the connected component containing $s$ of the times for which the
argument lies in $H(J)$. Every solution taking values in $J$ has this
form. If $g(c)=0$ at some point outside this branch, $x\equiv c$ must
be included separately as an equilibrium solution.
\end{theorem}
\begin{proof}
The sign of $g$ is constant on $J$, so $H'=1/g$ has a fixed nonzero
sign. The one-dimensional inverse-function theorem yields the stated
inverse. Along any solution in $J$, the chain rule gives
$(H(x(t)))'=a(t)$. Its initial value is zero, so integration proves
necessity. Differentiating the inverse formula proves sufficiency and
uniqueness within the branch. A zero of $g$ was excluded before division;
substitution into the original equation gives the separately stated
constant solutions.
\end{proof}

\begin{example}[Growth, logistic growth, and actual maximal intervals]\label{d13:02:logistic}
For $r,K>0$, the logistic equation $x'=rx(1-x/K)$ with $x(s)=\xi$
has
\[
 x(t)=\frac{K\xi e^{r(t-s)}}{K+\xi(e^{r(t-s)}-1)}.
\]
For $\xi=0$ and $\xi=K$ this gives the two equilibria. For
$0<\xi<K$, the denominator is positive for all real $t$, the solution
is global, increases, and has limits $0$ and $K$ as $t$ tends to
$-\infty$ and $+\infty$. If $\xi>K$, its maximal interval is
$(s+r^{-1}\log(1-K/\xi),\infty)$. If $\xi<0$, its maximal interval is
$(-\infty,s+r^{-1}\log(1-K/\xi))$. These claims follow by direct
differentiation, the sign of the denominator, and its simple zero at the
stated finite endpoint. Local uniqueness follows because the vector field
is $C^1$; no continuation through a pole is possible in $\R$.
Unrestricted exponential growth is the special linear equation $x'=rx$,
with solution $\xi e^{r(t-s)}$ and no finite-time pole.
\end{example}

\section{Exact equations and useful substitutions}
\begin{proposition}[Exactness on a rectangle]\label{d13:02:exact}
Let $R=I\times J$ be an open rectangle and let $P,Q\in C^1(R)$ with
$P_x=Q_t$. Fix $(s,\xi)\in R$. The function
\[
 F(t,x)=\int_s^t P(r,\xi)\,\dd r+\int_\xi^x Q(t,v)\,\dd v
\]
satisfies $F_t=P$, $F_x=Q$. A $C^1$ graph $x(t)$ satisfies
$P(t,x(t))+Q(t,x(t))x'(t)=0$ exactly when $F(t,x(t))$ is constant.
Where $Q\ne0$, this is a normal-form equation $x'=-P/Q$.
\end{proposition}
\begin{proof}
The fundamental theorem gives $F_x=Q$. Differentiation under the second
integral, valid because $Q_t$ is continuous on every compact subrectangle,
gives
\[
 F_t=P(t,\xi)+\int_\xi^x Q_t(t,v)\,\dd v
     =P(t,\xi)+\int_\xi^x P_x(t,v)\,\dd v=P(t,x).
\]
The graph assertion is the chain rule and the fact that a differentiable
function with zero derivative on an interval is constant. Division by $Q$
is justified only on its nonzero set. No assertion of global exactness on
an arbitrary nonsimply connected domain follows from this calculation.
\end{proof}

\begin{proposition}[Integrating factors]\label{d13:02:integrating-factor}
For $C^1$ functions $P,Q$ on a rectangle, any nowhere-zero $C^1$
function $\mu$ satisfying $(\mu P)_x=(\mu Q)_t$ converts
$P+Qx'=0$ into an exact equation without changing its solution set.
If $Q\ne0$ and $(P_x-Q_t)/Q=h(t)$ depends only on $t$, one may take
$\mu(t)=\exp(\int_s^t h(r)\,\dd r)$. If $P\ne0$ and
$(Q_t-P_x)/P=k(x)$ depends only on $x$, one may take
$\mu(x)=\exp(\int_\xi^x k(v)\,\dd v)$.
\end{proposition}
\begin{proof}
Multiplication by a nowhere-zero scalar is reversible. In the first
one-variable case, $\mu'Q=\mu(P_x-Q_t)$ is exactly the exactness
condition. In the second, $\mu'P=\mu(Q_t-P_x)$ is that same condition.
The exponential formulas solve these scalar linear equations and never
vanish. Apply Proposition~\ref{d13:02:exact}.
\end{proof}

\begin{proposition}[Bernoulli reduction]\label{d13:02:bernoulli}
For continuous real $a,b$, real $p\ne1$, and a positive unknown $x$,
the substitution $u=x^{1-p}$ gives a bijection between positive solutions
of $x'=a(t)x+b(t)x^p$ and positive solutions of
\[
 u'=(1-p)a(t)u+(1-p)b(t).
\]
The inverse is $x=u^{1/(1-p)}$ on the interval where $u>0$.
\end{proposition}
\begin{proof}
Positive real powers are $C^1$, and the chain rule gives the transformed
equation. Applying the chain rule to the inverse power proves the converse.
Positivity is part of both domains. For a general real exponent, extending
the formula through $u=0$ or to negative $u$ is not justified by this
substitution. Possible zero solutions are tested in the original equation
when its real power is defined there.
\end{proof}

\begin{proposition}[Riccati reduction about a known solution]\label{d13:02:riccati}
Let $a,b,c$ be continuous and let $v$ solve $v'=av^2+bv+c$ on $I$.
On any interval on which $x-v$ is nowhere zero, $x$ solves the same
Riccati equation if and only if $u=1/(x-v)$ solves
\[
 u'=-(2av+b)u-a,
\]
with $u$ nowhere zero. The inverse is $x=v+1/u$.
\end{proposition}
\begin{proof}
Substitute $x=v+1/u$ and subtract the equation for $v$. The result is
$-u'/u^2=(2av+b)/u+a/u^2$, which is equivalent to the displayed equation.
Each algebraic operation is reversible on the stated nonzero set. The
known solution $x=v$ is retained separately.
\end{proof}

\section{Oscillators and conservation}
\begin{theorem}[Harmonic oscillator]\label{d13:02:oscillator}
For $\omega>0$, the unique solution of $y''+\omega^2y=0$ with
$y(s)=a$, $y'(s)=b$ is
\[
 y(t)=a\cos(\omega(t-s))+(b/\omega)\sin(\omega(t-s)).
\]
It is global, and $((y')^2+\omega^2y^2)/2$ is constant. If $\omega=0$,
the solution is instead $a+b(t-s)$.
\end{theorem}
\begin{proof}
Differentiation proves existence and the initial conditions. The derivative
of the asserted energy is $y'(y''+\omega^2y)=0$. Applying this identity
to the difference of two solutions with zero initial data gives a sum of
nonnegative squares equal to zero, hence uniqueness when $\omega>0$.
For $\omega=0$, integrate $y''=0$ twice. The formulas are defined on all
of $\R$.
\end{proof}

\begin{proposition}[Mechanical energy]\label{d13:02:energy}
If $V\in C^1(J)$ and $y''=-V'(y)$ takes values in the interval $J$, then
\[
 E(t)=\tfrac12(y'(t))^2+V(y(t))
\]
is constant. On a branch where $E-V(y)>0$ and $y'$ has fixed sign,
the trajectory satisfies $y'=\pm\sqrt{2(E-V(y))}$ and may be recovered
by the separable formula. Turning points must be handled with the original
second-order equation rather than division by the vanishing square root.
\end{proposition}
\begin{proof}
The derivative of $E$ is $y'(y''+V'(y))=0$. Solving the energy identity
for $y'$ is legitimate with the chosen sign and positive radicand.
The chain-rule argument in Theorem~\ref{d13:02:separable} then applies on
that branch. Energy by itself does not choose the sign or the timing at a
turning point, which is why no converse there is asserted.
\end{proof}

\paragraph{Sources and coverage.}
The adopted methods synthesize [T, Sections~1.1--1.6] and [HSD,
Chapter~1]. Exact equations and substitution formulas are verified above
directly from their definitions. The examples supply equilibrium, sign,
denominator, and maximal-interval checks that are essential when an
explicit integration method divides by an unknown function.

% END chapters/M13.02


% BEGIN chapters/M13.03
\chapter{Local existence and uniqueness}\label{ch:d13-03}

\section{The Lipschitz hypothesis and Picard iteration}
\begin{definition}[Local Lipschitz control in the state]\label{d13:03:lipschitz}
A continuous $f:U\to\R^n$ is locally Lipschitz in the state, uniformly
in time, if each $(s,\xi)\in U$ has a product neighborhood on which
\[
 \norm{f(t,x)-f(t,y)}\leq L\norm{x-y}
\]
for a constant $L$ and all permitted $t,x,y$. For a parameter-dependent
field the constant is also uniform for parameters in a sufficiently small
neighborhood. This is a local hypothesis; one constant on all of $U$
is not required.
\end{definition}

\begin{lemma}[A derivative criterion]\label{d13:03:derivative-bound}
On a convex state ball $B$, if $D_xf$ exists, is continuous in all
variables, and has operator norm at most $L$, then the preceding
Lipschitz bound holds on $B$. In particular, a jointly $C^1$ field is
locally Lipschitz in the state, uniformly in its other variables.
\end{lemma}
\begin{proof}
The segment $y+\theta(x-y)$ stays in $B$. Integrating the derivative along
that segment yields
\[
 f(t,x)-f(t,y)=\int_0^1D_xf(t,y+\theta(x-y))(x-y)\,\dd\theta.
\]
The norm estimate follows by the integral triangle inequality. For the
local assertion choose a closed product box inside the open domain;
continuity bounds the derivative on its compact closure.
\end{proof}

\begin{lemma}[Uniform convergence of a contraction iteration]\label{d13:03:contraction}
If $K$ maps a nonempty closed subset $X$ of a Banach space into itself
and $\norm{Ku-Kv}\leq q\norm{u-v}$ for $0\leq q<1$, then $K$ has
exactly one fixed point $u_*$. For $u_{j+1}=Ku_j$,
\[
 \norm{u_j-u_*}\leq\frac{q^j}{1-q}\norm{u_1-u_0}.
\]
\end{lemma}
\begin{proof}
Iteration gives $\norm{u_{j+1}-u_j}\leq q^j\norm{u_1-u_0}$.
Summing from $j$ to $k-1$ bounds $\norm{u_k-u_j}$ by the corresponding
tail of a geometric series. Completeness and closedness give a limit in
$X$. The Lipschitz estimate makes $K$ continuous, so the limit is fixed.
Two fixed points have distance at most $q$ times that distance and hence
are equal. Taking $k\to\infty$ in the geometric tail gives the stated
error estimate.
\end{proof}

\begin{theorem}[Picard--Lindelof on a controlled box]\label{d13:03:picard}
Suppose $f$ is continuous on
$Q=[s-a,s+a]\times\overline B(\xi,b)\subset U$, has norm at most
$M$, and is Lipschitz in $x$ there with constant $L$. Choose $h>0$
with $h\leq a$, $Mh\leq b$, and $Lh<1$. There is a solution with
initial value $\xi$ on $[s-h,s+h]$, staying in the closed state ball.
It is the uniform limit of the Picard iterates
\[
 u_0(t)=\xi,\qquad
 u_{j+1}(t)=\xi+\int_s^t f(r,u_j(r))\,\dd r.
\]
It is the only solution of this initial value problem whose graph stays
in this box on the whole interval. If, additionally, $f$ is locally
Lipschitz in the state uniformly in time throughout $U$, any two local
solutions with these initial data agree on their common interval of definition.
\end{theorem}
\begin{proof}
Let $X$ be the closed subset of $C([s-h,s+h],\R^n)$ whose elements
take values in $\overline B(\xi,b)$, with the supremum norm. The
continuous-function space is complete: a uniformly Cauchy sequence has
coordinatewise limits, the convergence is uniform, and uniform limits
are continuous. Its subset $X$ is closed and contains $u_0$.
The integral operator $K$ in the theorem satisfies
$\norm{Ku(t)-\xi}\leq M|t-s|\leq Mh\leq b$, so it maps $X$ to $X$.
It also satisfies
\[
 \norm{Ku-Kv}_\infty\leq Lh\norm{u-v}_\infty.
\]
Apply Lemma~\ref{d13:03:contraction} and then
Theorem~\ref{d13:01:integral}. This proves existence and box uniqueness.
Under the additional local hypothesis throughout $U$, continuity first puts
both arbitrary local solutions inside one
smaller controlled box about any common point with equal values. Box
uniqueness therefore makes the set of times where the two solutions
agree open in the common interval; it is closed by continuity and
nonempty because it contains $s$. An interval is connected, so this set
is the whole common interval. The same argument works to either side of
$s$.
\end{proof}

\begin{corollary}[Local uniqueness under local hypotheses]\label{d13:03:uniqueness}
If $f$ is continuous and locally Lipschitz in the state uniformly in
time on the open set $U$, every $(s,\xi)\in U$ admits a local solution,
unique on overlaps with that initial value. Two solution trajectories
cannot intersect at a common time and then separate while both remain
defined.
\end{corollary}
\begin{proof}
Choose a sufficiently small compact box about $(s,\xi)$ on which the
local Lipschitz bound holds. Continuity bounds $f$ there. Some positive
$h$ satisfies the three inequalities in
Theorem~\ref{d13:03:picard}, including when $M=0$ or $L=0$, and the
theorem applies. Its overlap argument proves the last assertion with
the intersection time used as the initial time.
\end{proof}

\begin{theorem}[One interval for nearby data]\label{d13:03:uniform-box}
Suppose a field $f(t,x,p)$ is continuous and has uniform bounds $M,L$
on
\[
 [s_0-a,s_0+a]\times\overline B(\xi_0,b)\times\overline B(p_0,c)
\]
inside its open domain, with the bound $L$ applying to state differences.
Choose $h>0$ with $h\leq a$, $Mh\leq b/2$, and $Lh<1$.
For every $s\in(s_0-h/2,s_0+h/2)$,
$\xi\in B(\xi_0,b/2)$ and $p\in B(p_0,c)$, the corresponding solution
exists on the common interval $J=[s_0-h/2,s_0+h/2]$ and stays in
$\overline B(\xi_0,b)$. All three constants are independent of this
choice of nearby initial data.
\end{theorem}
\begin{proof}
On $C(J,\R^n)$ use the closed ball of paths within distance $b$ of
$\xi_0$. The map
$Ku(t)=\xi+\int_s^t f(r,u(r),p)\,\dd r$ takes this ball into itself
because $|t-s|\leq h$ and
$\norm{Ku(t)-\xi_0}\leq b/2+Mh\leq b$.
Its contraction constant is at most $Lh$. The preceding fixed-point
proof applies on this same path space for every permitted $(s,\xi,p)$.
The fixed point has value $\xi$ at $s$ and solves the required equation.
\end{proof}

\section{Existence for a merely continuous field}
\begin{lemma}[Compactness for equicontinuous curves]\label{d13:03:ascoli}
A uniformly bounded, equicontinuous sequence of maps from a compact
interval to $\R^n$ has a uniformly convergent subsequence.
\end{lemma}
\begin{proof}
Choose a countable dense set of times. Boundedness and finite-dimensional
Bolzano--Weierstrass give a subsequence converging at the first time,
a subsequence of it converging at the second, and so on. The diagonal
subsequence converges at every chosen time. Given $\varepsilon>0$,
equicontinuity supplies a single $\delta$ for which each path changes by
less than $\varepsilon/3$ at time distances below $\delta$. Finitely
many $\delta$-balls centered at the dense times cover the interval.
The diagonal sequence is Cauchy at these finitely many centers, so the
triangle inequality makes it uniformly Cauchy everywhere. Completeness
of the continuous-function space, proved in the Picard argument, gives
the uniform limit.
\end{proof}

\begin{theorem}[Peano existence]\label{d13:03:peano}
Suppose $f$ is continuous on
$[s,s+a]\times\overline B(\xi,b)$ and $\norm{f}\leq M$ there.
If $0<h\leq a$ and $Mh\leq b$, there is at least one solution on
$[s,s+h]$ starting at $\xi$ and staying in the closed ball. There is
an analogous solution to the left when the corresponding left box is
available. No uniqueness assertion is made.
\end{theorem}
\begin{proof}
For each $N\geq1$ subdivide $[s,s+h]$ into equal steps
$\Delta=h/N$, with nodes $t_j=s+j\Delta$. Put $u_N(s)=\xi$ and set
\[
 u_N(t)=u_N(t_j)+(t-t_j)f(t_j,u_N(t_j))
 \quad(t_j\leq t\leq t_{j+1}).
\]
Induction gives $\norm{u_N(t)-\xi}\leq M(t-s)\leq b$:
the slope on each new segment has norm at most $M$, so the accumulated
path length gives the claimed bound and validates the next evaluation
of $f$. The curves are all $M$-Lipschitz and uniformly bounded.
Lemma~\ref{d13:03:ascoli} gives a uniformly convergent subsequence
$u_{N_k}\to u$.

Let $\omega(\delta)$ be a modulus of uniform continuity of $f$ on
the compact box, using the product metric $|t-r|+\norm{x-y}$.
On each mesh cell the difference between
$f(t_j,u_N(t_j))$ and $f(t,u_N(t))$ has norm at most
$\omega((1+M)\Delta)$. Integration over the cells, including the last
partial cell, therefore gives
\[
 \sup_{s\leq t\leq s+h}
 \left\|u_N(t)-\xi-\int_s^t f(r,u_N(r))\,\dd r\right\|
 \leq h\,\omega((1+M)h/N)\longrightarrow0.
\]
Uniform convergence of $u_{N_k}$ and uniform continuity of $f$ imply
uniform convergence of $f(r,u_{N_k}(r))$ to $f(r,u(r))$.
The integral identity passes to the limit, so
Theorem~\ref{d13:01:integral} proves the assertion. Reverse time for
the left-hand version. Pasting the two curves at the initial value
gives a two-sided local solution in an open box if desired.
\end{proof}

\begin{example}[Continuity does not give uniqueness]\label{d13:03:nonunique}
The continuous field $f(x)=2\sqrt{|x|}$ has the solutions
\[
 x_a(t)=\begin{cases}0,&t\leq a,\\(t-a)^2,&t\geq a,\end{cases}
 \qquad a\geq0,
\]
all with $x_a(0)=0$, as well as the identically zero solution. The
derivative exists and is zero at the joining time, so these are
classical solutions. The quotient
$|f(x)-f(0)|/|x|=2/\sqrt{x}$ for $x>0$ is unbounded near zero;
the Lipschitz hypothesis fails precisely where the waiting time is
possible.
\end{example}

\section{A useful weaker uniqueness criterion}
\begin{theorem}[Osgood uniqueness]\label{d13:03:osgood}
In a compact local box suppose
\[
 \norm{f(t,x)-f(t,y)}\leq\omega(\norm{x-y}),
\]
where $\omega:[0,\infty)\to[0,\infty)$ is continuous and
nondecreasing, $\omega(0)=0$, $\omega(r)>0$ for $r>0$, and
$\int_0^\delta \dd r/\omega(r)=\infty$ for every $\delta>0$.
Then two solutions with the same initial value agree while they stay
in a sufficiently small common box. Local applications give uniqueness
on their whole common interval.
\end{theorem}
\begin{proof}
For times to the right of $s$ put $d(t)=\norm{x(t)-y(t)}$. The
integral equations give $d(t)\leq\int_s^t\omega(d(r))\,\dd r$.
For $\varepsilon>0$ let
$v_\varepsilon(t)=\varepsilon+\int_s^t\omega(d(r))\,\dd r$.
Then $d\leq v_\varepsilon$ and
$v_\varepsilon'\leq\omega(v_\varepsilon)$ by monotonicity. Since
$v_\varepsilon\geq\varepsilon$, the chain rule and integration yield
\[
 \int_\varepsilon^{v_\varepsilon(t)}
       \frac{\dd r}{\omega(r)}\leq t-s.
\]
If $d(t)>0$, choose $0<\delta<d(t)$ and then $\varepsilon<\delta$.
The left side is at least $\int_\varepsilon^\delta\dd r/\omega(r)$,
which diverges as $\varepsilon\downarrow0$, a contradiction. Thus
$d(t)=0$. Time reversal gives the left side. The open-and-closed overlap
argument from Theorem~\ref{d13:03:picard} extends the conclusion.
\end{proof}

\paragraph{Sources and coverage.}
The fixed-point and Euler-polygonal routes are [T, Sections~2.1--2.3
and~2.7]; the modulus criterion develops [T, the uniqueness exercises
following Section~2.7]. [CL, Chapter~1, Sections~1--4, and Chapter~2,
Sections~2--3] provide complementary existence and uniqueness frameworks.
[HSD, Sections~17.1--17.2] supplies the same classical local theorem.
Only finite-dimensional Peano compactness is adopted here; a direct
replacement by an arbitrary infinite-dimensional state space is not claimed.

% END chapters/M13.03


% BEGIN chapters/M13.04
\chapter{Maximal solutions, continuation, and blow-up}\label{ch:d13-04}

Unless a theorem explicitly says otherwise, $f:U\to\R^n$ is continuous
on the open set $U\subset\R\times\R^n$ and locally Lipschitz in the
state uniformly in time. The uniqueness qualification is retained in
every construction of the solution map.

\section{Maximal intervals}
\begin{theorem}[Unique maximal solution]\label{d13:04:maximal}
For every $(s,\xi)\in U$ there is exactly one maximal solution through
that initial datum. Its domain is an open interval
\[
 I(s,\xi)=(T_-(s,\xi),T_+(s,\xi)),\qquad T_-<s<T_+,
\]
where endpoints may be infinite. Every other solution through the datum
is a restriction of it on its own domain.
\end{theorem}
\begin{proof}
Consider all solutions on open intervals containing $s$. This collection
is nonempty by Corollary~\ref{d13:03:uniqueness}. Any pair agrees on its
intersection by that corollary. The union of all the intervals is open
and is an interval because all of them contain $s$. Assign at each time
the value of any solution defined there. Agreement on overlaps makes
this assignment well defined. Around each time one of the constituent
solutions supplies a $C^1$ representative, so the union is a solution.
Every other solution is contained in it. An extension would itself be
one of the solutions already included in the union, which is impossible
on a strictly larger interval. This proves existence and uniqueness of
the maximal solution.
\end{proof}

\begin{proposition}[Changing the base time]\label{d13:04:rebase}
If $r\in I(s,\xi)$ and $\eta=\Phi(r,s,\xi)$, then
$I(r,\eta)=I(s,\xi)$ and
$\Phi(t,r,\eta)=\Phi(t,s,\xi)$ throughout that interval.
\end{proposition}
\begin{proof}
The maximal solutions share the value $\eta$ at time $r$ and agree on
the intersection. Pasting them gives a solution on the union. Maximality
of each forces the two intervals to be equal, and uniqueness gives the
identity there. This is a nonautonomous composition identity; no time
translation of $f$ is being assumed.
\end{proof}

\section{The actual obstruction to extension}
\begin{theorem}[Finite-endpoint continuation criterion]\label{d13:04:continuation}
Let $x$ solve the equation on $(a,b)$ with $b<\infty$. It extends to a
solution on $(a,b+\varepsilon)$ for some $\varepsilon>0$ if and only if
there is a sequence $t_j\uparrow b$ such that
\[
 (t_j,x(t_j))\longrightarrow(b,\eta)\in U.
\]
In that case $x(t)\to\eta$ as $t\uparrow b$ and the extended value at
$b$ is $\eta$. The time-reversed assertion holds at a finite left endpoint.
\end{theorem}
\begin{proof}
An extension is continuous at $b$, so necessity follows. Conversely,
apply Theorem~\ref{d13:03:uniform-box}, with no parameter, to a compact
box about $(b,\eta)\in U$. It supplies a fixed interval
$J=[b-h/2,b+h/2]$ on which solutions starting at times near $b$ with
values near $\eta$ all exist. For large $j$, the datum
$(t_j,x(t_j))$ is one of these permitted data and $t_j\in J$.
Let $y_j$ be its solution on $J$. Uniqueness gives $y_j=x$ on
$J\cap(a,b)$. Pasting one such $y_j$ to $x$ extends the solution to
the right of $b$. Since $y_j$ is continuous at $b$, $x$ has a limit
there. The given subsequence forces that limit to equal $\eta$.
Restriction to an open interval inside the pasted domain gives the
stated extension. Time reversal proves the left-endpoint assertion.
\end{proof}

\begin{corollary}[Escape from every compact subset]\label{d13:04:escape}
If the maximal right endpoint $T_+$ is finite, the graph of its solution
eventually leaves every compact set $K\subset U$ as $t\uparrow T_+$.
If $U=\R\times\R^n$, then necessarily
$\norm{\Phi(t,s,\xi)}\to\infty$ as $t\uparrow T_+$. Analogous
statements hold at a finite left endpoint.
\end{corollary}
\begin{proof}
Otherwise some compact $K\subset U$ contains graph points at times
$t_j\uparrow T_+$. Compactness gives a convergent subsequence whose
limit has time coordinate $T_+$ and lies in $K\subset U$.
The continuation theorem contradicts maximality. In the full-space
case, failure of the claimed norm limit supplies a sequence with bounded
state values. Those graph points lie in
$[s,T_+]\times\overline B(0,R)\subset U$ for a suitable $R$, a compact
set, again contradicting the first assertion.
\end{proof}

\begin{corollary}[Compact confinement gives global forward existence]\label{d13:04:confinement}
Suppose a forward solution has the following property: for every finite
$B>s$, all its values at times $t\in[s,B)$ at which it is defined lie
in a compact set $K_B\subset\R^n$ with $[s,B]\times K_B\subset U$.
Then its maximal right endpoint is $+\infty$.
\end{corollary}
\begin{proof}
If $T_+<\infty$, choose $B>T_+$. The compact set
$[s,T_+]\times K_B\subset U$ contradicts
Corollary~\ref{d13:04:escape}.
\end{proof}

\begin{theorem}[Global existence under linear growth]\label{d13:04:linear-growth}
Suppose $U=\R\times\R^n$ and, for every finite $B>0$, there are
nonnegative constants $a_B,b_B$ such that
\[
 \norm{f(t,x)}\leq a_B+b_B\norm{x}
 \qquad(|t|\leq B,\ x\in\R^n).
\]
Every maximal solution is global. More quantitatively, on any available
interval between $s$ and $t$ contained in $[-B,B]$, it satisfies
\[
 \norm{x(t)}\leq
 \begin{cases}
 \norm{x(s)}e^{b_B|t-s|}
       +a_B(e^{b_B|t-s|}-1)/b_B,&b_B>0,\\
 \norm{x(s)}+a_B|t-s|,&b_B=0.
 \end{cases}
\]
\end{theorem}
\begin{proof}
For $t\geq s$ the integral equation bounds $\norm{x(t)}$ by
$\norm{x(s)}+\int_s^t(a_B+b_B\norm{x(r)})\,\dd r$.
The elementary integral estimate in
Theorem~\ref{d13:05:gronwall} yields the displayed bound. Its proof
uses only an integrating factor and is independent of continuation.
Time reversal gives the same estimate for $t<s$. If a maximal endpoint
were finite, choose $B$ larger than its absolute value and $|s|$. The
estimate bounds all state values approaching that endpoint, in
contradiction with Corollary~\ref{d13:04:escape}. Thus neither endpoint
is finite.
\end{proof}

\section{Examples that separate the possibilities}
\begin{example}[Quadratic blow-up]\label{d13:04:quadratic-blowup}
The solution of $x'=x^2$, $x(s)=\xi$, is
$x(t)=\xi/(1-\xi(t-s))$. For $\xi>0$ the maximal interval is
$(-\infty,s+1/\xi)$ and $x(t)\to+\infty$ at the right endpoint.
For $\xi<0$ it is $(s+1/\xi,\infty)$, with a pole at the left
endpoint. For $\xi=0$ it is the global zero solution. Direct
differentiation checks the equation. The field is smooth, so uniqueness
identifies the solution on its entire pole-free interval; the pole
prevents any extension through the finite endpoint.
\end{example}

\begin{example}[A bounded state can hit the domain boundary]\label{d13:04:domain-boundary}
On $U=\R\times(-1,1)$, the equation $x'=1$, $x(0)=0$, has the
maximal solution $x(t)=t$ on $(-1,1)$. It is bounded but not global
in the prescribed domain. Its graph escapes every compact subset of
$U$ by approaching $x=1$ or $x=-1$. Extending the formula to all of
$\R$ solves an equation on a different state domain, not the original
problem.
\end{example}

\begin{example}[Local Lipschitz continuity is not linear growth]\label{d13:04:local-not-global}
The field $x^2$ is locally Lipschitz everywhere but has finite-time
blow-up for positive initial data. Conversely, $\sin(x^2)$ is bounded,
so the linear-growth theorem makes all its solutions global, although
its derivative $2x\cos(x^2)$ is unbounded and the field is not globally
Lipschitz. A bounded derivative is a convenient sufficient condition,
not a necessary one for global existence.
\end{example}

\paragraph{Sources and coverage.}
The maximal-union and compact-escape arguments refine [T,
Section~2.6] and [HSD, Section~17.4]. The distinction between leaving
the domain and diverging in norm is retained explicitly. [CL,
Chapter~1, Section~4, and Chapter~2, Section~1] is a complementary
continuation reference. No uniqueness-free maximal solution selection
is used later in this manuscript.

% END chapters/M13.04


% BEGIN chapters/M13.05
\chapter{Comparison estimates and continuous dependence}\label{ch:d13-05}

\section{Integral inequalities}
\begin{theorem}[Gronwall with an inhomogeneous term]\label{d13:05:gronwall}
Let $u,b,c$ be continuous nonnegative functions on $[s,T]$, and let
$A\geq0$. If
\[
 u(t)\leq A+\int_s^t\bigl(b(r)u(r)+c(r)\bigr)\,\dd r,
\]
then
\[
 u(t)\leq A\exp\left(\int_s^t b\right)
       +\int_s^t c(r)\exp\left(\int_r^t b\right)\,\dd r.
\]
For constant $b=L\geq0$ and $c=\delta\geq0$, the right side is
$Ae^{L(t-s)}+\delta E_L(t-s)$, where
\[
 E_L(v)=\begin{cases}(e^{Lv}-1)/L,&L>0,\\v,&L=0.\end{cases}
\]
\end{theorem}
\begin{proof}
Set $v(t)=A+\int_s^t(bu+c)$. Then $0\leq u\leq v$ and
$v'=bu+c\leq bv+c$. Multiplying by
$\exp(-\int_s^t b)$ gives
\[
 \left(v(t)e^{-\int_s^t b}\right)'
       \leq c(t)e^{-\int_s^t b}.
\]
Integrate and multiply by $e^{\int_s^t b}$ to obtain the bound for
$v$ and hence $u$. The constant-coefficient evaluation uses the
elementary integral of the exponential, with its separate value at
$L=0$. The same argument after time reversal proves the corresponding
inequality for an oriented interval to the left of $s$.
\end{proof}

\begin{corollary}[A variable initial bound]\label{d13:05:general-gronwall}
If $u,b,a$ are continuous nonnegative functions on $[s,T]$ and
$u(t)\leq a(t)+\int_s^t b(r)u(r)\,\dd r$, then
\[
 u(t)\leq a(t)+\int_s^t a(r)b(r)
                \exp\left(\int_r^t b\right)\,\dd r.
\]
If $a$ is nondecreasing, this is at most
$a(t)\exp(\int_s^t b)$.
\end{corollary}
\begin{proof}
Put $v(t)=\int_s^t bu$. Then $v'\leq b(a+v)$, $v(s)=0$.
An integrating factor, exactly as in the preceding proof, bounds $v$
by the integral in the display, and $u\leq a+v$. When $a(r)\leq a(t)$,
factor out $a(t)$ and use
$\int_s^t b(r)e^{\int_r^t b}\,\dd r=e^{\int_s^t b}-1$.
\end{proof}

\begin{theorem}[Perturbation and residual estimate]\label{d13:05:perturbation}
Let $x$ solve $x'=f(t,x)$ and let $y$ be a $C^1$ curve on a compact
interval containing $s$. Assume all required graph points lie in the
domain of $f$,
\[
 \norm{f(t,x(t))-f(t,y(t))}\leq L\norm{x(t)-y(t)},\qquad
 \norm{y'(t)-f(t,y(t))}\leq\delta
\]
there, with $L,\delta\geq0$. Then
\[
 \norm{x(t)-y(t)}\leq
 e^{L|t-s|}\norm{x(s)-y(s)}+\delta E_L(|t-s|).
\]
In particular, if $y'=g(t,y)$ and $\norm{f(t,z)-g(t,z)}\leq\delta$
on a set containing the graph of $y$, this is the same bound for two
different vector fields.
\end{theorem}
\begin{proof}
Subtract the integral equations, writing
$y(t)=y(s)+\int_s^t[f(r,y(r))+(y'-f(r,y(r)))]\,\dd r$.
The triangle inequality bounds the difference by its initial norm plus
the integral of $L$ times that difference and $\delta$.
Theorem~\ref{d13:05:gronwall} proves the assertion to the right of $s$;
time reversal proves it to the left. The final assertion follows by
substituting $y'=g(t,y)$ into the residual. Only the field $f$ needs
the displayed Lipschitz bound in this estimate.
\end{proof}

\begin{corollary}[Sharp dependence on initial values]\label{d13:05:initial-bound}
Two solutions of the same $L$-Lipschitz field, on a common interval and
within the set where that bound holds, satisfy
\[
 \norm{\Phi(t,s,\xi)-\Phi(t,s,\eta)}
       \leq e^{L|t-s|}\norm{\xi-\eta}.
\]
The exponential cannot in general be replaced by a uniform-in-time
constant: $x'=Lx$ gives equality for $t\geq s$ in one dimension.
\end{corollary}
\begin{proof}
Take $g=f$ and $\delta=0$ in
Theorem~\ref{d13:05:perturbation}. The explicit scalar solution from
Theorem~\ref{d13:02:linear} verifies the equality example.
\end{proof}

\section{Scalar order comparison}
\begin{theorem}[Subsolutions and supersolutions]\label{d13:05:comparison}
Let $J$ be an open real interval and let $f:[a,b]\times J\to\R$ be
continuous and locally Lipschitz in its second variable uniformly in
time. Suppose $u,v\in C^1([a,b],J)$ satisfy
\[
 u'\leq f(t,u),\qquad v'\geq f(t,v),\qquad u(a)\leq v(a).
\]
Then $u(t)\leq v(t)$ for every $t\in[a,b]$.
\end{theorem}
\begin{proof}
The two ranges are compact subsets of $J$. Choose a compact interval
inside $J$ containing these ranges in its interior. Local Lipschitz
control gives a constant $L$ valid on this whole compact rectangle:
take a finite cover by controlled neighborhoods and subdivide any
state segment into pieces lying in one member of that cover.
For sufficiently small $\varepsilon>0$, the perturbed function
$v_\varepsilon(t)=v(t)+\varepsilon e^{(L+1)(t-a)}$ remains in that
interval. The Lipschitz estimate implies
\[
 v_\varepsilon'(t)\geq f(t,v(t))+
       (L+1)\varepsilon e^{(L+1)(t-a)}
       >f(t,v_\varepsilon(t)).
\]
Initially $u(a)<v_\varepsilon(a)$. If a first contact occurs, the
derivative of $u-v_\varepsilon$ at that contact is nonnegative, since
the function approached zero from below. At the contact the two state
arguments of $f$ coincide, while the differential inequalities make
that derivative strictly negative, a contradiction. Hence
$u<v_\varepsilon$ on the whole interval. Letting $\varepsilon$
decrease to zero proves the result.
\end{proof}

\section{A common domain is part of continuous dependence}
\begin{theorem}[Persistence on a compact solution interval]\label{d13:05:common-interval}
Let $\bar x$ be a solution of a continuous, locally state-Lipschitz
field $f:U\to\R^n$, and let $[a,b]$ be a compact subinterval of its
open solution interval. There are $\rho>0$ and $L\geq0$ such that
the compact tube
\[
 Q_\rho=\{(t,z):a\leq t\leq b,
                 \norm{z-\bar x(t)}\leq\rho\}
\]
lies in $U$ and $f$ has state Lipschitz constant $L$ on each tube
slice. Let $g:U\to\R^n$ be another continuous, locally state-Lipschitz
field with $\sup_{Q_\rho}\norm{g-f}\leq\delta$. For $s'\in[a,b]$
and $\xi'$ put $d=\norm{\xi'-\bar x(s')}$ and $T=b-a$. If
\[
 e^{LT}(d+\delta T)<\rho/2,
\]
the solution of $y'=g(t,y)$, $y(s')=\xi'$, exists on the entire
interval $[a,b]$ and satisfies
\[
 \norm{y(t)-\bar x(t)}
 \leq de^{L|t-s'|}+\delta E_L(|t-s'|)<\rho/2.
\]
\end{theorem}
\begin{proof}
The reference graph is compact in the open domain. Product boxes with
controlled state Lipschitz constants cover it. Choose finitely many
and shrink the time neighborhoods, using continuity of $\bar x$, so
that one fixed-radius ball about $\bar x(t)$ fits in a controlled box
for each $t\in[a,b]$. The minimum of the finitely many radii gives a
positive $\rho$, reduced if necessary so the closed tube lies in $U$;
the maximum of the constants is $L$. Each state slice is convex, so
its full Lipschitz estimate is valid.

The datum for $g$ is inside $Q_\rho$ because $d<\rho/2$. It has a
unique local solution by Corollary~\ref{d13:03:uniqueness}. On any
part of its interval before an exit from the tube, apply
Theorem~\ref{d13:05:perturbation} with the reference field $f$ and
the curve $y$. This gives the displayed bound, since
$E_L(v)\leq ve^{Lv}$. That bound is strictly below $\rho/2$, so an
exit at distance $\rho$ cannot occur by continuity. If a maximal
endpoint occurred inside $[a,b]$, the graph on approach to it would
lie in the compact inner tube $Q_{\rho/2}\subset U$.
Corollary~\ref{d13:04:escape}, applied to $g$, rules out that endpoint.
Thus the solution reaches both ends of $[a,b]$, and the same estimate
holds throughout. The strict margin also permits extension slightly
past either endpoint when it is necessary to interpret derivatives.
\end{proof}

\begin{corollary}[Locally uniform convergence of vector fields]\label{d13:05:field-convergence}
Suppose $f_j,f:U\to\R^n$ are continuous and locally state-Lipschitz,
and $f_j\to f$ uniformly on compact subsets of $U$. If
$(s_j,\xi_j)\to(s,\xi)$ and $[a,b]$ is a compact subinterval of
$I_f(s,\xi)$ with $s$ in its interior, then, for all sufficiently
large $j$, the solutions for $f_j$ starting at $(s_j,\xi_j)$ exist
on $[a,b]$ and converge there uniformly to the solution for $f$.
No uniform bound on the Lipschitz constants of $f_j$ is required.
\end{corollary}
\begin{proof}
Use the tube of Theorem~\ref{d13:05:common-interval} around
$\bar x(t)=\Phi_f(t,s,\xi)$. On that compact tube let
$\delta_j=\sup\norm{f_j-f}\to0$. For large $j$, $s_j\in[a,b]$,
and $d_j=\norm{\xi_j-\bar x(s_j)}\to0$ by continuity of
$\bar x$. The numerical condition of that theorem eventually holds,
and its uniform estimate tends to zero. All comparison constants came
from the reference field $f$; local Lipschitz control for each $f_j$
was used only to give its own uniquely extendible solution.
\end{proof}

\begin{theorem}[Open maximal domain and joint continuity]\label{d13:05:openness}
Let $U\subset\R\times\R^n\times\R^m$ be open and let
$f(t,x,p)$ be continuous and locally state-Lipschitz uniformly in
time and parameters. The maximal solution domain
\[
 \mathcal D=\{(t,s,\xi,p):(s,\xi,p)\in U,
                   \ t\in I(s,\xi,p)\}
\]
is open, and $\Phi:\mathcal D\to\R^n$ is jointly continuous.
\end{theorem}
\begin{proof}
Fix $(t_0,s_0,\xi_0,p_0)\in\mathcal D$ and choose a compact
interval $[a,b]$ in the reference solution interval with both
$s_0,t_0$ in its interior. Its graph at parameter $p_0$ is compact
in $U$. The finite-box argument from
Theorem~\ref{d13:05:common-interval} supplies a state tube and a
closed parameter ball about $p_0$, both of positive radius, whose
product lies in $U$, with a common state Lipschitz constant.
Uniform continuity on this compact set gives
\[
 \sup_{(r,z)\in Q_\rho}
       \norm{f(r,z,p)-f(r,z,p_0)}\longrightarrow0
       \quad(p\to p_0).
\]
For $(s,\xi,p)$ close to $(s_0,\xi_0,p_0)$, the initial discrepancy
$\norm{\xi-\Phi(s,s_0,\xi_0,p_0)}$ is small. Apply
Theorem~\ref{d13:05:common-interval} to the fields at parameters
$p_0$ and $p$. Every such solution exists on $[a,b]$ and converges
uniformly there to the reference solution as the data approach the
reference data. A neighborhood of $t_0$ is also contained in $[a,b]$.
This proves that a product neighborhood of the chosen point lies in
$\mathcal D$. Evaluating the uniform convergence at variable times
$t\to t_0$, and using continuity of the reference curve, proves joint
continuity.
\end{proof}

\begin{corollary}[Semicontinuity of lifespan]\label{d13:05:lifespan}
On its initial-data domain, $T_+(s,\xi,p)$ is lower semicontinuous
and $T_-(s,\xi,p)$ is upper semicontinuous, with infinite endpoints
understood in the extended real sense.
\end{corollary}
\begin{proof}
If $c<T_+(s_0,\xi_0,p_0)$ and $c>s_0$, choose a compact reference
interval extending from slightly before $s_0$ to slightly after $c$.
The preceding proof puts that same interval in all nearby solution
domains, so their right endpoints exceed $c$. If $c\leq s_0$, the
same conclusion follows after making the nearby initial times exceed
$c$, or by choosing an intermediate value between $s_0$ and the right
endpoint. This is lower semicontinuity, including $T_+=+\infty$ by
testing every finite $c$. Reverse time for the left endpoint.
\end{proof}

\begin{example}[Continuous dependence is a finite-time assertion]\label{d13:05:finite-time}
For $x'=x$, solutions with initial values $0$ and $\varepsilon$ at
zero have difference $\varepsilon e^t$. It tends to zero uniformly on
each fixed compact time interval as $\varepsilon\to0$, but not
uniformly on $[0,\infty)$. Similarly, for $x'=x^2$ the common interval
around the zero solution can be chosen arbitrarily long when the
initial value is sufficiently small; a fixed nonzero positive initial
value still has finite forward lifespan. These are compatible with,
not exceptions to, the compact-interval and semicontinuity theorems.
\end{example}

\paragraph{Sources and coverage.}
[T, Section~2.4] supplies the Gronwall and data-perturbation framework;
[HSD, Section~17.3] treats dependence on initial conditions. The compact
tube argument above makes the common-interval and open-domain steps
explicit. [CL, Chapter~1, Section~7, and Chapter~2, Section~4] are
additional references for initial-value and parameter variations.
Continuity alone in a parameter gives continuity, not an unasserted
parameter derivative or a uniform-in-infinite-time stability conclusion.

% END chapters/M13.05


% BEGIN chapters/M13.06
\chapter{Differentiable and analytic dependence}\label{ch:d13-06}

The distinction between continuous, $C^k$, smooth, and analytic dependence
is substantive. This chapter proves the stronger conclusions under their
corresponding stronger hypotheses. All real solution maps are defined on
the open maximal domain $\mathcal D$ of Theorem~\ref{d13:05:openness}.
Differentiability assertions are local on that open domain; no common
infinite time interval is inferred.

\section{A linear integral equation used in the proof}
\begin{lemma}[Linear Volterra equations]\label{d13:06:linear-volterra}
Let $A:[a,b]\to\R^{n\times n}$ and
$B:[a,b]\to\R^{n\times r}$ be continuous, let $s\in[a,b]$, and
let $Y_s\in\R^{n\times r}$. There is a unique solution on $[a,b]$ of
\[
 Y(t)=Y_s+\int_s^t\bigl(A(u)Y(u)+B(u)\bigr)\,\dd u.
\]
If $\norm{A}\leq L$ and $\norm{B}\leq M$ there, then
$\norm{Y(t)}\leq\norm{Y_s}e^{L|t-s|}+ME_L(|t-s|)$.
Continuous variation of $A,B,Y_s$, uniform on this interval, gives
uniform-in-time continuous variation of $Y$; in particular a jointly
continuous family of coefficients and initial values gives a jointly
continuous family of solutions locally in its parameters.
\end{lemma}
\begin{proof}
For the forward interval, start $Y_0(t)=Y_s$ and iterate the displayed
integral operator. Put $C=L\norm{Y_s}+M$. Induction, by integrating
the previous difference, gives
\[
 \norm{Y_{j+1}(t)-Y_j(t)}
       \leq \frac{CL^j(t-s)^{j+1}}{(j+1)!}.
\]
The bound remains valid when $L=0$ by interpreting the zeroth power
as $1$; subsequent differences then vanish. The exponential series
converges uniformly, so the iterates converge uniformly to a continuous
solution and the integral passes to the limit. Reverse time for the
backward interval and paste at $s$. Gronwall applied to the difference
of two solutions proves uniqueness, and applied to the integral norm
inequality gives the asserted bound.

For varying coefficients indexed by $q$, subtract the equations for
$Y_q,Y_{q_0}$ and write the integrand difference as
\[
 A_q(Y_q-Y_{q_0})+(A_q-A_{q_0})Y_{q_0}+(B_q-B_{q_0}).
\]
Locally in $q$ all coefficient norms and the initial norms are uniformly
bounded. Uniform continuity of the data makes the last two terms
uniformly small. The same Gronwall estimate therefore makes
$\norm{Y_q-Y_{q_0}}_\infty$ tend to zero. Continuous dependence on
$t$ together with this uniform convergence proves the joint assertion.
\end{proof}

\section{First derivatives without assuming their existence}
\begin{theorem}[$C^1$ dependence and the first variational equations]\label{d13:06:c1-dependence}
Let $f\in C^1(U,\R^n)$, where
$U\subset\R\times\R^n\times\R^m$ is open. Then
$\Phi(t,s,\xi,p)$ is jointly $C^1$ on $\mathcal D$. With the
initial time, state, and parameter fixed, put
\[
 A(t)=D_xf(t,\Phi(t,s,\xi,p),p),\qquad
 Y(t)=D_\xi\Phi(t,s,\xi,p),\quad
 Z(t)=D_p\Phi(t,s,\xi,p).
\]
These derivatives are the uniquely determined solutions of
\[
 Y'=AY,\quad Y(s)=I_n,\qquad
 Z'=AZ+D_pf(t,\Phi(t,s,\xi,p),p),\quad Z(s)=0.
\]
Also $\partial_t\Phi=f(t,\Phi,p)$.
\end{theorem}
\begin{proof}
First consider an autonomous $C^1$ field $F$ and its flow
$\psi(t,\xi)$, with initial time zero. Fix a solution and a compact
time interval $J$ containing zero inside its maximal interval. By
Theorem~\ref{d13:05:common-interval}, all initial states in a small
neighborhood of $\xi$ have solutions on the same $J$, in a common
compact tube. Shrink that neighborhood and the tube if necessary so
state segments used below stay inside a compact subset of the domain.
On this set let $\norm{DF}\leq L$, and write
$T=\max_{t\in J}|t|$. The initial-value estimate gives
\[
 \norm{\psi(t,\xi+h)-\psi(t,\xi)}\leq e^{LT}\norm h.
\]
Let $Y(t,\xi)$ be the solution of
$Y'=DF(\psi(t,\xi))Y$, $Y(0)=I_n$, furnished by the preceding
lemma, not by assuming a derivative of $\psi$.

Write $\Delta_h(t)=\psi(t,\xi+h)-\psi(t,\xi)$. Uniform
continuity of $DF$ on the compact tube and integration along state
segments imply
\[
 F(\psi(t,\xi)+\Delta_h(t))-F(\psi(t,\xi))
   =DF(\psi(t,\xi))\Delta_h(t)+R_h(t),
\]
where $\norm{R_h(t)}\leq
\omega(e^{LT}\norm h)e^{LT}\norm h$ for a modulus
$\omega(\delta)\to0$ as $\delta\to0$. Subtract the integral equation
for $Yh$ from that for $\Delta_h$. The error
$E_h=\Delta_h-Yh$ has zero initial value and satisfies
\[
 \norm{E_h(t)}\leq
 L\left|\int_0^t\norm{E_h(r)}\,\dd r\right|
 +T e^{LT}\omega(e^{LT}\norm h)\norm h.
\]
Here and below an integral of a nonnegative function over a reversed
interval means its nonnegative absolute value. Applying Gronwall on
either side of zero yields
\[
 \sup_{t\in J}\frac{\norm{E_h(t)}}{\norm h}
       \leq T e^{2LT}\omega(e^{LT}\norm h)\longrightarrow0.
\]
This is the Frechet derivative in the initial state, uniformly on $J$.
The coefficient $DF(\psi(t,\xi))$ is jointly continuous by
Theorem~\ref{d13:05:openness}; the linear Volterra dependence lemma
makes $Y$ jointly continuous. The time derivative is
$F(\psi(t,\xi))$, also jointly continuous. All first partial derivatives
therefore exist continuously, so $\psi$ is jointly $C^1$. The argument
applies around any point of its maximal solution domain because the
reference trajectory between zero and that time is compact.

For the stated nonautonomous, parameter-dependent equation, apply this
autonomous result to
\[
 F(\tau,x,p)=(1,f(\tau,x,p),0)
\]
on the given open set. This field is $C^1$, and its elapsed-time flow
$\Psi$ gives
\[
 \Phi(t,s,\xi,p)=\pi_x\Psi(t-s,(s,\xi,p)).
\]
Composition with these linear coordinate maps proves joint $C^1$
dependence in all the stated variables. Finally differentiate the
integral equation for $\Phi$ with respect to $\xi$ and $p$. On every
compact subinterval the integrands and their first data derivatives
are continuous and uniformly bounded, so difference quotients may be
integrated and the limit taken uniformly. This gives
\[
 Y(t)=I_n+\int_s^t A(r)Y(r)\,\dd r,
 \qquad Z(t)=\int_s^t\bigl(A(r)Z(r)+D_pf(r,\Phi(r),p)\bigr)\,\dd r.
\]
The fundamental theorem of calculus gives their differential equations.
The initial values follow by differentiating $\Phi(s,s,\xi,p)=\xi$.
Uniqueness is the linear Volterra lemma.
\end{proof}

\begin{proposition}[Derivative in the starting time]\label{d13:06:initial-time}
Under the preceding hypotheses,
\[
 \partial_s\Phi(t,s,\xi,p)
    =-D_\xi\Phi(t,s,\xi,p)f(s,\xi,p).
\]
In particular, this is not in general the negative of the terminal
velocity $-f(t,\Phi(t,s,\xi,p),p)$ for a nonautonomous field.
\end{proposition}
\begin{proof}
The rebasing identity from Proposition~\ref{d13:04:rebase} gives, for
small $h$ on a common domain,
\[
 \Phi(t,s+h,\Phi(s+h,s,\xi,p),p)=\Phi(t,s,\xi,p).
\]
Differentiate in $h$ at zero. Joint $C^1$ dependence just proved
justifies the chain rule, and
$\partial_h\Phi(s+h,s,\xi,p)|_{h=0}=f(s,\xi,p)$.
The derivative of the right side is zero, giving the formula. For
example, $x'=t$ has $\Phi(t,s,\xi)=\xi+(t^2-s^2)/2$, whose
starting-time derivative is $-s$, rather than $-t$.
\end{proof}

\section{Higher derivatives and regular perturbation}
\begin{theorem}[$C^k$ and smooth dependence]\label{d13:06:ck-dependence}
If $f\in C^k(U,\R^n)$ for an integer $k\geq1$, then $\Phi$ is
jointly $C^k$ on $\mathcal D$, and its terminal-time derivative
$\partial_t\Phi=f(t,\Phi,p)$ is itself jointly $C^k$. If $f$ is
$C^\infty$, so is $\Phi$.
\end{theorem}
\begin{proof}
It suffices, by autonomization, to prove the autonomous result for $F$.
The case $k=1$ is the preceding theorem. Assume the assertion of
$C^k$ dependence is known for every finite-dimensional $C^k$ vector
field, and let $F\in C^{k+1}$. The first variational system can itself
be written as the autonomous system
\[
 (y,V)'=(F(y),DF(y)V).
\]
Its vector field is $C^k$ on the product of the state domain with
$\R^{n\times n}$. By the inductive assertion its solution starting
at $(\xi,I_n)$ is $C^k$ in elapsed time and initial state. Its
components are $\psi(t,\xi)$ and $D_\xi\psi(t,\xi)$, by the
already proved first variational equation and uniqueness. Thus the
state partial derivatives of $\psi$ are $C^k$. Also $\psi$ is
$C^k$ by the inductive assertion applied to $F$, so
$\partial_t\psi=F\circ\psi$ is $C^k$. A map whose first partial
derivatives are all $C^k$ is $C^{k+1}$, by the finite-dimensional
calculus criterion. This completes the induction. The variational
solutions exist on every compact part of the original solution
interval, since their linear Volterra bounds preclude a new finite
endpoint there.

Returning to $f$, autonomization includes both the clock and constant
parameters as state coordinates, so the induction proves the full
joint statement, not just derivatives with respect to $\xi$.
The identity $\partial_t\Phi=f(t,\Phi,p)$ and composition of
$C^k$ maps give the additional regularity of that derivative.
Applying the conclusion for every finite $k$ proves the smooth case.
\end{proof}

\begin{theorem}[Higher variational equations]\label{d13:06:variations}
Let $F\in C^k(\Omega,\R^n)$ and
$H_j(t,\xi)=D_\xi^j\psi(t,\xi)$ for $1\leq j\leq k$.
For vectors $v_1,\ldots,v_j$, $H_1(0)=I_n$,
$H_j(0)=0$ for $j\geq2$, and
\[
 \frac{\dd}{\dd t}H_j(t)[v_1,\ldots,v_j]
 =\sum_{\pi}
 D^{|\pi|}F(\psi(t))
 \left[H_{|B|}(t)[v_i:i\in B]:B\in\pi\right],
\]
where the sum is over all partitions $\pi$ of $\{1,\ldots,j\}$.
The partition with one block contributes $DF(\psi)H_j$; all other
terms use variations of order less than $j$. In particular,
\[
 \begin{split}
 (H_2[v,w])'&=DF(\psi)H_2[v,w]
             +D^2F(\psi)[H_1v,H_1w],\\
 H_2(0)[v,w]&=0.
 \end{split}
\]
Thus each new-order variation solves a linear inhomogeneous equation
once the lower-order variations are known.
\end{theorem}
\begin{proof}
Differentiate the integral identity
$\psi(t,\xi)=\xi+\int_0^t F(\psi(r,\xi))\,\dd r$ a total of
$j$ times in the initial state. Theorem~\ref{d13:06:ck-dependence}
makes every required state derivative continuous. On a smaller compact
data neighborhood, these derivatives are uniformly bounded; successive
difference quotients and the fundamental theorem for differentiation
under a continuous integral justify each interchange. The higher chain
rule follows inductively by placing a newly differentiated direction
either into one existing block or into its own block; this enumerates
each partition exactly once and yields the displayed sum. The initial
map is $\xi\mapsto\xi$, giving the stated initial derivatives.
The resulting integral equation has a continuous integrand, so its
time derivative exists and equals that sum. In a partition with at
least two blocks every block is smaller than $j$, which proves the
linear dependence on the highest-order variation. The formula for
$H_2$ is the two-partition case. Symmetry of these multilinear
derivatives and equality of their mixed components are supplied by
the $C^j$ calculus theorem, not by an unproved formal differentiation.
\end{proof}

\begin{corollary}[First-order regular perturbation]\label{d13:06:perturbation-expansion}
Let $f(t,x,\varepsilon)$ be $C^1$ near the graph of a reference
solution $x_0$ on a compact interval $J$ containing $s$. Suppose
$\xi(\varepsilon)$ is $C^1$ near zero, with $\xi(0)=\xi_0$, and
let $x_\varepsilon(s)=\xi(\varepsilon)$. For all sufficiently small
$\varepsilon$, $x_\varepsilon$ exists on the same $J$ and
\[
 x_\varepsilon(t)=x_0(t)+\varepsilon z(t)+o(\varepsilon)
\]
uniformly for $t\in J$, where
\[
 z'=D_xf(t,x_0,0)z+\partial_\varepsilon f(t,x_0,0),
 \qquad z(s)=\xi'(0).
\]
If the field and initial curve are $C^2$, the remainder is
$O(\varepsilon^2)$ uniformly on $J$.
\end{corollary}
\begin{proof}
Common existence follows from Theorem~\ref{d13:05:common-interval}.
The $C^1$ dependence theorem and the chain rule give the derivative
$z(t)$ and its inhomogeneous variational equation, including its initial
value. On $J$ times a small closed parameter interval that derivative
is uniformly continuous. The integral form of the one-variable Taylor
formula therefore gives a remainder divided by $|\varepsilon|$ that
tends uniformly to zero. In the $C^2$ case the second parameter
derivative is uniformly bounded on this compact set; integrating its
bound twice gives the asserted quadratic remainder.
\end{proof}

\begin{corollary}[Invertibility of the solution map]\label{d13:06:local-diffeomorphism}
For fixed $s,t,p$, let
$\Omega_{s,t,p}=\{\xi:(t,s,\xi,p)\in\mathcal D\}$. If $f$ is
$C^k$, $k\geq1$, the map
$\xi\mapsto\Phi(t,s,\xi,p)$ is a $C^k$ diffeomorphism from the
open set $\Omega_{s,t,p}$ onto $\Omega_{t,s,p}$, with inverse
$\eta\mapsto\Phi(s,t,\eta,p)$. In particular its state derivative
is invertible at every allowed datum. These sets need not be all of
$\R^n$.
\end{corollary}
\begin{proof}
The sets are open by Theorem~\ref{d13:05:openness}. Rebasing the
maximal solution at time $t$ gives both inverse identities and shows
that each image lies in the stated reverse-time domain. Both maps are
$C^k$ by Theorem~\ref{d13:06:ck-dependence}. Differentiating the two
inverse identities makes their derivatives inverse linear maps.
\end{proof}

\section{Analytic dependence requires analytic hypotheses}
\begin{lemma}[Uniform limits on complex polydiscs]\label{d13:06:polydisc-limit}
A locally uniform limit of holomorphic maps on an open subset of
$\C^d$ is holomorphic. Here holomorphic means locally represented by
an absolutely convergent power series in the $d$ complex variables.
The assertion applies componentwise to maps with values in $\C^n$.
\end{lemma}
\begin{proof}
Fix a closed polydisc of radii $r_1,\ldots,r_d$ strictly inside the
domain, with center $a$. Iterating the one-variable Cauchy formula
expresses each member $g_j$ of the sequence on its interior as
\[
 g_j(z)=\frac{1}{(2\pi i)^d}
 \int_{|\zeta_1-a_1|=r_1}\cdots
 \int_{|\zeta_d-a_d|=r_d}
 \frac{g_j(\zeta)}{\prod_{\nu=1}^d(\zeta_\nu-z_\nu)}
 \,\dd\zeta_d\cdots\dd\zeta_1.
\]
Local uniform convergence includes uniform convergence on the compact
integration torus, so the formula passes to the limit $g$. On a
smaller concentric polydisc of radii $r'_\nu<r_\nu$, expand each
denominator in its geometric series. The resulting product series
converges absolutely and uniformly, bounded by a constant times
$\prod_\nu(1-r'_\nu/r_\nu)^{-1}$. Its integrals may therefore be
taken term by term and give a convergent power series for $g$ there.
All Cauchy formulas used are the one-variable formulas from the prior
complex-analysis foundation. The displayed argument supplies the
several-variable uniform-limit bridge needed here.
\end{proof}

\begin{theorem}[Holomorphic local dependence]\label{d13:06:analytic-dependence}
Let $f(z,w,\lambda)$ be holomorphic on an open neighborhood of
$(z_0,w_0,\lambda_0)\in\C\times\C^n\times\C^m$. The complex
initial value problem
\[
 \partial_z w=f(z,w,\lambda),\qquad w(z_0)=\xi
\]
has, for $(z,\xi,\lambda)$ in sufficiently small product
neighborhoods of $(z_0,w_0,\lambda_0)$, a solution jointly
holomorphic in these variables and unique as a local solution germ.
The initial complex time may also vary near $z_0$, with joint
holomorphic dependence on that time, the terminal time, initial state,
and parameters.
\end{theorem}
\begin{proof}
Translate $z_0$ to zero. Choose a closed time disc, a closed convex
complex state ball $\overline B(w_0,b)$, and a closed parameter
polydisc whose product lies inside the domain, with a slightly larger
product still inside it. Holomorphic differentiation and compactness
give bounds $M$ for $\norm f$ and $L$ for $\norm{D_w f}$ on that
product. Integrating $D_wf$ along straight state segments gives the
complex state Lipschitz bound $L$. Restrict initial states to
$B(w_0,b/4)$ and choose $h>0$ inside the time-disc radius with
$hM\leq b/4$ and $hL<1/2$.

Start $u_0(z,\xi,\lambda)=\xi$ and iterate
\[
 u_{j+1}(z,\xi,\lambda)=\xi+
       \int_0^z f(\zeta,u_j(\zeta,\xi,\lambda),\lambda)
                   \,\dd\zeta.
\]
The integrals are on the straight segment; the integrands are
holomorphic in $\zeta$ on a disc, so the one-variable Cauchy theorem
also makes them independent of the path inside that disc. Each
iterate has distance at most $b/2$ from the fixed state center $w_0$,
since its initial-state distance is at most $b/4$ and its integral
has norm at most $hM\leq b/4$. This also verifies inductively that
every evaluation of $f$ stays inside the product where $M,L$ apply.

The supremum norm over this whole product of variables gives a
contraction factor at most $hL<1/2$ between successive iterates.
The geometric-series argument of Lemma~\ref{d13:03:contraction}
therefore gives uniform convergence to a continuous map $w$.
Every iterate is jointly holomorphic: composition preserves local
power series, and writing its integral as
$z\int_0^1 f(\theta z,u_j(\theta z,\xi,\lambda),\lambda)\,\dd\theta$
shows holomorphy by locally uniform Riemann sums. Those sums are
holomorphic and their convergence is uniform on each smaller compact
product; Lemma~\ref{d13:06:polydisc-limit} applies.
The same lemma makes the uniform limit $w$ jointly holomorphic.
Uniform convergence and uniform state Lipschitz control pass the
integral equation to the limit. Differentiating it in $z$ gives the
complex differential equation and the initial value.

Any other holomorphic solution germ stays in the same state ball on
a sufficiently small disc, by continuity at its initial value. The
same contraction estimate on that smaller disc makes the two
solutions equal there, which is precisely uniqueness of germs.
For variable initial time, autonomize with the holomorphic vector
field $(1,f(\tau,w,\lambda),0)$ and evaluate its elapsed-time solution
at $z-\eta$, starting from $(\eta,\xi,\lambda)$. The preceding
argument applies to this larger finite-dimensional state. This
proves joint holomorphy in the varying initial time $\eta$ as well.
\end{proof}

\begin{corollary}[Real-analytic dependence]\label{d13:06:real-analytic}
If a real field $f(t,x,p)$ is real analytic on its open domain, its
real maximal solution map is real analytic on $\mathcal D$.
\end{corollary}
\begin{proof}
An absolutely convergent real power series extends, after shrinking
its polydisc, to a holomorphic complex power series. Apply
Theorem~\ref{d13:06:analytic-dependence} to such a local
complexification. Its restriction to real variables is real valued:
complex conjugation gives another solution with the same real data,
and uniqueness identifies them. It is therefore the real solution,
which is locally real analytic for small elapsed time.

For any allowed longer real elapsed time, the reference trajectory
between its endpoints is compact in the open domain. Cover it by
finitely many of these local analytic solution neighborhoods and
subdivide the elapsed-time interval so that each consecutive step
uses one such neighborhood. Continuity and the uniform local
existence neighborhoods keep nearby data in the same finite chain.
The resulting solution is the finite composition of the local
analytic maps, by uniqueness. All but the final time increment may
be fixed; varying the final increment in a small neighborhood gives
joint analyticity also in the terminal time. Autonomization already
included the starting time and parameters among the initial-state
coordinates. Thus the solution map is real analytic around every
point of $\mathcal D$. This proof uses no global assertion about
single-valued continuation around complex singularities.
\end{proof}

\begin{example}[A smooth equation with a nonanalytic solution]\label{d13:06:smooth-not-analytic}
Set $g(t)=e^{-1/t^2}$ for $t>0$ and $g(t)=0$ for $t\leq0$.
Each positive-side derivative is a polynomial in $1/t$ times
$e^{-1/t^2}$ and tends to zero at zero, because the exponential
decays faster than every power. Induction therefore shows that the
extension is $C^\infty$ and every derivative at zero vanishes.
The equation $x'=g(t)$, $x(0)=0$, has the smooth solution
$x(t)=\int_0^t g(r)\,\dd r$. All its derivatives at zero vanish,
but it is positive for every $t>0$. Its Taylor series at zero is
identically zero and does not represent it on any neighborhood.
Thus the smooth-dependence theorem cannot be used as an
analytic-dependence theorem.
\end{example}

\paragraph{Sources and coverage.}
The real variational arguments develop [T, Sections~2.4--2.5,
especially Theorems~2.10--2.12] and [HSD, Section~17.6].
[CL, Chapter~1, Sections~7--8, and Chapter~2, Section~4] is a
complementary source for parameter variations and complex systems.
The holomorphic Picard proof expands the uniformity step in [T,
Section~4.1, especially Theorem~4.2] and treats vector states and
parameters explicitly. No assertion about analytic dependence is
derived merely from $C^\infty$ regularity.

% END chapters/M13.06


% BEGIN chapters/M13.07
\chapter{Linear systems and matrix exponentials}\label{ch:d13-07}

Throughout this chapter \(\mathbb K=\R\) or \(\C\), \(I\) is a real interval,
and matrix and vector norms are compatible. Time is real, even for complex
unknowns. The distinction matters when complex eigenvalues describe a real
system.

\section{The exponential and constant coefficients}

\begin{definition}[Matrix exponential]\label{d13:07:exp-definition}
For \(A\in\mathbb K^{n\times n}\) put
\[
 e^{tA}=\sum_{j=0}^{\infty}\frac{t^j A^j}{j!},\qquad t\in\R .
\]
\end{definition}

\begin{theorem}[Exponential calculus]\label{d13:07:exp}
The series is locally uniformly convergent, is infinitely differentiable,
and satisfies
\[
 \frac{\dd}{\dd t}e^{tA}=Ae^{tA}=e^{tA}A,\quad e^{0A}=I,\quad
 \norm{e^{tA}}\le e^{|t|\norm A}.
\]
Moreover \(e^{(s+t)A}=e^{sA}e^{tA}\), so \((e^{tA})^{-1}=e^{-tA}\).
If \(AB=BA\), then \(e^{A+B}=e^Ae^B\); no such conclusion is asserted for
noncommuting matrices.
\end{theorem}
\begin{proof}
On \(|t|\le R\) the \(j\)-th term has norm at most
\((R\norm A)^j/j!\), and its differentiated series is bounded by the
corresponding exponential series times \(\norm A\). Repeated derivatives
are treated in the same way. Uniform convergence therefore justifies
termwise differentiation, and summing the norm majorant gives the bound.
The absolutely convergent product of the series for \(sA,tA\) can be
regrouped by total degree. Its coefficient of \(A^k\) is
\(\sum_{j=0}^k s^j t^{k-j}/(j!(k-j)!)=(s+t)^k/k!\).
For commuting \(A,B\), the same calculation uses the binomial identity for
commuting matrices. Setting \(s=-t\) proves invertibility.
\end{proof}

\begin{proposition}[Jordan blocks and real solutions]\label{d13:07:jordan}
Over \(\C\), a block \(J=\lambda I+N\), \(N^m=0\), has
\[
 e^{tJ}=e^{\lambda t}\sum_{j=0}^{m-1}\frac{t^jN^j}{j!}.
\]
Consequently every component of a constant-coefficient solution is a finite
sum of polynomials times exponentials. For a real matrix, a conjugate pair
\(\alpha\pm i\beta\) gives real combinations of
\(t^j e^{\alpha t}\cos(\beta t)\) and
\(t^j e^{\alpha t}\sin(\beta t)\).
If every eigenvalue has real part less than \(-a<0\), then for each
\(0<b<a\) some \(C_b<\infty\) satisfies
\(\norm{e^{tA}}\le C_b e^{-bt}\) for \(t\ge0\).
\end{proposition}
\begin{proof}
The scalar part commutes with \(N\), so the preceding theorem and nilpotence
give the formula. Similarity transports it to \(A\), using the Jordan
decomposition supplied by linear algebra. Taking real and imaginary parts
produces the real solutions. Finally each polynomial \(t^j\) times
\(e^{-(a-b)t}\) is bounded on \([0,\infty)\), and there are finitely many
blocks. The similarity matrix only changes the constant.
\end{proof}

\section{Fundamental matrices on a whole interval}

\begin{theorem}[Evolution matrix]\label{d13:07:fundamental}
Let \(A:I\to\mathbb K^{n\times n}\) be continuous and \(s\in I\).
There is a unique matrix \(U(t,s)\), defined for all \(t\in I\), with
\[
 \partial_tU(t,s)=A(t)U(t,s),\qquad U(s,s)=I .
\]
It is invertible and, whenever the indicated times belong to \(I\),
\[
 U(t,r)U(r,s)=U(t,s),\quad U(t,s)^{-1}=U(s,t),\quad
 \partial_sU(t,s)=-U(t,s)A(s).
\]
In particular
\(\norm{U(t,s)}\le\exp(\int_{\min(s,t)}^{\max(s,t)}\norm{A(r)}\,\dd r)\).
Every fundamental matrix is \(X(t)=U(t,s)C\) for a fixed invertible \(C\).
\end{theorem}
\begin{proof}
Apply local existence and uniqueness to each column, with initial data the
coordinate basis. For complex coefficients, identify \(\C^n\) with
\(\R^{2n}\) and replace \(A\) by its real block representation; time remains
real and the same uniqueness, integral and norm estimates apply.
On each compact subinterval, Gronwall gives the displayed
bound (reverse time for \(t<s\)). A finite endpoint inside \(I\) would then
allow continuation, so the solution covers \(I\).
If \(U(t,s)v=0\), the solution through zero at \(t\), traced back to \(s\),
is identically zero by uniqueness; hence \(v=0\). Finite dimensionality
gives invertibility. Both sides of the composition formula solve the same
initial-value problem at \(r\). It also gives
\(U(t,s)=U(t,s_0)U(s,s_0)^{-1}\); differentiating the inverse proves the
\(s\)-derivative formula. If \(X\) is fundamental, uniqueness applied to its
columns yields \(X(t)=U(t,s)X(s)\), which proves the last assertion.
\end{proof}

\begin{theorem}[Variation of constants]\label{d13:07:variation}
For continuous \(b:I\to\mathbb K^n\), the unique solution of
\(x'=A(t)x+b(t)\), \(x(s)=x_0\), is
\[
 x(t)=U(t,s)x_0+\int_s^t U(t,r)b(r)\,\dd r
     =X(t)\left(X(s)^{-1}x_0+\int_s^t X(r)^{-1}b(r)\,\dd r\right).
\]
\end{theorem}
\begin{proof}
The product rule gives
\((X^{-1}x)'=X^{-1}(x'-Ax)=X^{-1}b\). Integration proves necessity.
Conversely differentiating the displayed expression, including the endpoint
term in the integral, proves the equation and the initial value. The same
calculation applied to a difference proves uniqueness.
\end{proof}

\begin{theorem}[Liouville's formula]\label{d13:07:liouville}
For every fundamental matrix,
\[
 \det X(t)=\det X(s)\,
   \exp\left(\int_s^t\operatorname{tr} A(r)\,\dd r\right).
\]
For a real system this determinant never changes sign.
\end{theorem}
\begin{proof}
Differentiating a determinant by its columns and using \(X'=AX\) gives
\((\det X)'=(\operatorname{tr}A)\det X\). Equivalently one can apply the
adjugate identity, which is polynomial in the entries and does not require
an a priori division by the determinant. Solve this scalar linear equation
using its initial value. The exponential is positive for real \(A\).
\end{proof}

\begin{example}[Polynomial growth is not exponential growth]\label{d13:07:nilpotent-example}
For \(A=\left(\begin{smallmatrix}0&1\\0&0\end{smallmatrix}\right)\),
\(e^{tA}=\left(\begin{smallmatrix}1&t\\0&1\end{smallmatrix}\right)\).
The only eigenvalue is zero, but the origin is not Lyapunov stable:
initial data \((0,\varepsilon)\) reach first component \(1\) at
\(t=1/\varepsilon\). Thus a nonpositive real-part condition alone does not
guarantee bounded solutions.
\end{example}

\paragraph{Source and coverage.}
Teschl, Chapter 3, supplies the principal comparison; the linear-system
chapters of Coddington--Levinson and Perko provide the classical and geometric
viewpoints. Jordan decomposition is an existing linear-algebra foundation.
The results above apply to arbitrary finite dimension, not only planar
examples.

% END chapters/M13.07


% BEGIN chapters/M13.08
\chapter{Periodic linear equations and Floquet theory}\label{ch:d13-08}

Let \(A:\R\to\mathbb K^{n\times n}\) be continuous and \(T\)-periodic,
where \(T>0\). Normalize the principal matrix by \(X(0)=I\).

\begin{definition}[Monodromy and multipliers]\label{d13:08:monodromy}
The monodromy matrix is \(M=X(T)\). Its complex eigenvalues, counted with
algebraic multiplicity, are the Floquet multipliers. With an initial time
\(s\), the monodromy is \(U(s+T,s)\); these matrices are similar, since
\(U(s+T,s)=X(s)MX(s)^{-1}\). Hence the multipliers do not depend on \(s\).
\end{definition}

\begin{lemma}[Translation identity]\label{d13:08:translation}
For all real \(t\), \(X(t+T)=X(t)M\). Thus
\(X(kT+r)=X(r)M^k\) for every integer \(k\) and real \(r\).
\end{lemma}
\begin{proof}
\(X(t+T)\) and \(X(t)M\) solve \(Y'=A(t)Y\) and agree at \(t=0\).
Uniqueness proves the first identity, and induction, using invertibility
for negative integers, proves the second.
\end{proof}

\begin{lemma}[Complex matrix logarithm]\label{d13:08:matrix-log}
Every invertible complex matrix \(M\) is \(e^L\) for some complex matrix \(L\).
\end{lemma}
\begin{proof}
Put \(M\) in Jordan form. On a block \(J=\lambda I+N\) of size \(m\),
\(\lambda\ne0\), choose any scalar logarithm \(\ell\) of \(\lambda\), and set
\[
 L_J=\ell I+\sum_{j=1}^{m-1}\frac{(-1)^{j+1}}{j}(N/\lambda)^j .
\]
The formal scalar identity \(\exp(\log(1+z))=1+z\), reduced modulo
\(z^m\), is an identity of finite polynomials after all terms of degree
\(\ge m\) are removed. Substitution of the nilpotent \(N/\lambda\)
is therefore legitimate and yields \(e^{L_J}=J\).
Take the direct sum and conjugate back. This proof neither chooses one
global scalar logarithm branch on \(\C\setminus\{0\}\) nor asserts reality.
\end{proof}

\begin{theorem}[Complex Floquet representation]\label{d13:08:floquet}
Even if \(A\) is real, regard it temporarily as complex. There are a constant
complex matrix \(B\) and a complex \(C^1\), invertible, \(T\)-periodic matrix
\(P(t)\) such that
\[
 X(t)=P(t)e^{tB},\qquad P(0)=I .
\]
The change of variables \(x=P(t)y\) converts \(x'=A(t)x\) into \(y'=By\).
\end{theorem}
\begin{proof}
Choose \(L\) with \(e^L=M\) and put \(B=L/T\),
\(P(t)=X(t)e^{-tB}\). Since \(M=e^{TB}\) commutes with every \(e^{tB}\),
\[
 P(t+T)=X(t)M e^{-tB}e^{-TB}=X(t)e^{-tB}=P(t).
\]
Differentiation gives \(P'=AP-PB\), exactly the claimed transformed
equation. Invertibility and regularity follow from those of both factors.
\end{proof}

\begin{example}[A real \(T\)-periodic representation can fail]\label{d13:08:real-log-counterexample}
Take \(T=1\), let \(R(t)\) denote rotation through angle \(\pi t\), and set
\[
 X(t)=R(t)\begin{pmatrix}1&0\\0&2^t\end{pmatrix},\qquad
 A(t)=\pi\begin{pmatrix}0&-1\\1&0\end{pmatrix}
       +R(t)\begin{pmatrix}0&0\\0&\log2\end{pmatrix}R(t)^{-1}.
\]
Then \(A(t+1)=A(t)\), \(X'=AX\), and \(M=\operatorname{diag}(-1,-2)\).
If \(e^B=M\) with real \(B\), then \(B\) commutes with \(M\), so its
off-diagonal entries vanish. Its exponential would consequently have two
positive diagonal entries, a contradiction. A real, period-one Floquet
factorization would imply precisely such a \(B\).
The complex theorem therefore cannot be silently read as a real theorem
with the same period.
\end{example}

\begin{theorem}[Multiplier criteria]\label{d13:08:stability}
The zero solution is exponentially asymptotically stable if and only if all
multipliers satisfy \(|\mu|<1\). It is Lyapunov stable if and only if all
\(|\mu|\le1\) and each Jordan block belonging to \(|\mu|=1\) has size one.
If some \(|\mu|>1\), or a unit-modulus multiplier has a nontrivial Jordan
block, it is unstable.
\end{theorem}
\begin{proof}
On the compact interval \([0,T]\), \(X\) and \(X^{-1}\) are bounded.
The translation identity reduces all forward-time bounds to bounds on
\(M^k\), \(k\ge0\). Jordan's formula gives exponentially decreasing powers
exactly when the spectral radius is less than one, and bounded powers
exactly under the second condition. A multiplier of modulus one prevents
exponential convergence along an eigenvector. A larger modulus or a
nontrivial unit-modulus Jordan block gives one vector whose iterates are
unbounded. Multiplying that vector by an arbitrarily small scalar proves
failure of Lyapunov stability. For real matrices take real or imaginary
parts; at least one has the required unbounded orbit. For a linear system,
boundedness of the evolution matrix is equivalent to Lyapunov stability,
directly from \(x(t)=X(t)x_0\).
\end{proof}

\begin{proposition}[Periodic solutions and forcing]\label{d13:08:periodic-forcing}
The homogeneous system has a nonzero \(T\)-periodic solution exactly when
\(1\) is an eigenvalue of \(M\). If \(b\) is continuous and \(T\)-periodic,
a \(T\)-periodic solution of \(x'=Ax+b\) exists exactly when
\[
 (I-M)x_0=M\int_0^T X(r)^{-1}b(r)\,\dd r
\]
is solvable. In particular it exists uniquely if \(1\) is not a multiplier.
\end{proposition}
\begin{proof}
Variation of constants computes \(x(T)\); the displayed condition is
exactly \(x(T)=x(0)\). When it holds, \(x(t+T)\) and \(x(t)\) solve the same
inhomogeneous equation with the same initial value, so they agree for all
time. The homogeneous assertion is the special case \(b=0\).
\end{proof}

\paragraph{Source and coverage.}
The principal source is Teschl, Section 3.6. The matrix-logarithm step is
included explicitly. We adopt the complex representation and the real
counterexample, rather than leaving a real-logarithm or doubled-period
convention implicit.

% END chapters/M13.08


% BEGIN chapters/M13.09
\chapter{Second-order equations and convergent series}\label{ch:d13-09}

\section{Wronskians and reduction of order}

\begin{proposition}[Reduction of order]\label{d13:09:reduction}
On a real interval, let \(p,q\) be continuous and let \(y_1\) be a nowhere
zero solution of \(y''+py'+qy=0\). A linearly independent second solution is
\[
 y_2(t)=y_1(t)\int_{t_0}^t
          \frac{\exp(-\int_{t_0}^s p(r)\,\dd r)}{y_1(s)^2}\,\dd s .
\]
The same formula holds locally on a simply connected complex domain on
which the coefficients and \(y_1\) are holomorphic and \(y_1\ne0\),
using holomorphic primitives.
\end{proposition}
\begin{proof}
Write \(y=y_1v\). Substitution and cancellation of the \(v\)-term give
\(y_1v''+(2y_1'+py_1)v'=0\). Thus
\(v'=\exp(-\int p)/y_1^2\) solves the remaining first-order equation.
Its Wronskian with \(y_1\) equals
\(y_1^2v'=\exp(-\int p)\ne0\), proving independence.
This argument also explains why the formula is used only where \(y_1\)
does not vanish.
\end{proof}

\begin{theorem}[Series at an ordinary point]\label{d13:09:ordinary-series}
Suppose \(p(z)=\sum_{k\ge0}p_kz^k\) and
\(q(z)=\sum_{k\ge0}q_kz^k\) are holomorphic on \(|z|<R\).
Every solution of \(y''+py'+qy=0\) with prescribed \(y(0),y'(0)\)
is holomorphic there, and its coefficients satisfy
\[
 (n+2)(n+1)a_{n+2}
 =-\sum_{k=0}^n\big((n-k+1)p_ka_{n-k+1}+q_ka_{n-k}\big).
\]
Conversely these coefficients determine that solution.
\end{theorem}
\begin{proof}
Turn the equation into \(Y'=A(z)Y\) with holomorphic \(A\).
For completeness its solution on the whole disc can be constructed
without a path-independence assumption: set \(Y_0(z)=Y(0)\) and
\(Y_{j+1}(z)=\int_0^z A(\zeta)Y_j(\zeta)\,\dd\zeta\).
Holomorphic primitives on a disc exist. On \(|z|\le\rho<R\),
the radial integral and induction give
\(\norm{Y_j(z)}\le\norm{Y(0)}M^j|z|^j/j!\), where
\(M=\max_{|z|\le\rho}\norm{A(z)}\).
Thus \(\sum_jY_j\) converges locally uniformly to a holomorphic function.
Its integrated identity, or locally uniform convergence of
\(\sum_{j\ge0} A Y_j\), gives \(Y'=AY\), with the correct initial value.
Local uniqueness identifies it with every initial-value solution.
Holomorphic Taylor expansion is therefore available throughout the
disc. Multiplication of absolutely convergent series gives the recurrence;
\(a_0,a_1\) and the nonzero denominators determine all remaining
coefficients. This proves convergence, not merely a formal recurrence.
\end{proof}

\section{Regular singular points}

\begin{definition}[Indicial polynomial]\label{d13:09:indicial}
For
\[
 z^2y''+zp(z)y'+q(z)y=0,\qquad p,q\text{ holomorphic on }|z|<R,
\]
the indicial polynomial is \(I(r)=r(r-1)+p_0r+q_0\).
Choose a branch of \(\log z\) on a slit punctured disc and write
\(z^r=e^{r\log z}\). Claims involving these powers use that chosen domain.
\end{definition}

\begin{theorem}[Nonresonant Frobenius solutions]\label{d13:09:frobenius}
If \(I(r)=0\) and \(I(r+n)\ne0\) for every integer \(n\ge1\), there is a
solution
\[
 y=z^r\sum_{n=0}^\infty a_nz^n,\qquad a_0=1,
\]
whose power series converges on \(|z|<R\), with
\[
 I(r+n)a_n
 =-\sum_{k=1}^n\big((r+n-k)p_k+q_k\big)a_{n-k}.
\]
If the two indicial roots differ by a noninteger, this construction gives
two independent solutions.
\end{theorem}
\begin{proof}
Substitution gives the recurrence by matching powers; the stated
nonvanishing makes it uniquely solvable. We check convergence.
Fix \(0<\rho<R_0<R\). Cauchy estimates give
\(|p_k|,|q_k|\le M R_0^{-k}\) after increasing \(M\).
Since \(I(r+n)\) is monic quadratic in \(n\), it has modulus at least
\(c n^2\) for all sufficiently large \(n\). Put \(b_n=|a_n|\rho^n\)
and \(B_n=\max_{j\le n}b_j\). The recurrence implies
\[
 b_n\le \frac{M}{c n^2}
       \sum_{k=1}^n(n+|r|+1)(\rho/R_0)^k B_{n-1}
       \le \frac{C}{n}B_{n-1}.
\]
For \(n\) large the coefficient is at most one, so all \(b_n\) are bounded.
The series then converges absolutely on every smaller disc
\(|z|<\rho\), as do its differentiated series. As \(\rho<R\) was
arbitrary, the convergence radius is at least \(R\), and the formal
calculation becomes an actual solution identity.
For roots \(r_1,r_2\), the Wronskian of their constructed solutions has
leading term \((r_2-r_1)z^{r_1+r_2-1}\), which is nonzero near zero.
\end{proof}

\begin{theorem}[Repeated and resonant roots]\label{d13:09:resonance}
Suppose the roots are ordered so that \(r_1-r_2=m\) is a nonnegative
integer. There is a solution \(y_1=z^{r_1}a(z)\), \(a(0)=1\).
On a sufficiently small slit disc a second independent solution has form
\[
 \begin{cases}
 z^{r_2}b(z)+c\,y_1(z)\log z,\quad b(0)=1, & m\ge1,\\
 y_1(z)\log z+z^{r_1}b(z), &m=0,
 \end{cases}
\]
with \(a,b\) holomorphic. In the first case \(c\) may be zero; it must
not be assumed zero merely because the coefficients are analytic away
from the singular point.
\end{theorem}
\begin{proof}
The larger root satisfies \(I(r_1+n)\ne0\) for \(n\ge1\), so the preceding
theorem constructs \(y_1\). Shrink the disc so that \(a\) has no zeros.
Reduction of order for \(y''+(p(z)/z)y'+(q(z)/z^2)y=0\) gives
\[
 y_2=y_1\int z^{-m-1}h(z)\,\dd z,\qquad
 h(z)=\frac{\exp(-\int_0^z(p(\zeta)-p_0)/\zeta\,\dd\zeta)}{a(z)^2},
 \quad h(0)=1.
\]
Here \(r_1+r_2=1-p_0\) accounts for the exponent. Expand
\(h(z)=\sum h_nz^n\). Uniform convergence on smaller discs permits
termwise integration:
\[
 \int z^{-m-1}h(z)\,\dd z
 =h_m\log z+\sum_{n\ne m}\frac{h_n}{n-m}z^{n-m}.
\]
For \(m\ge1\) the leading term is \(-z^{-m}/m\); multiply the second
solution by \(-m\) to normalize \(b(0)=1\).
For \(m=0\), \(h_0=1\), giving the second displayed form.
The integrated series converges locally, and its derivative is the
integrand, so the Wronskian is nonzero by reduction of order.
\end{proof}

\begin{example}[Euler and Bessel equations]\label{d13:09:singular-examples}
For \(z^2y''+azy'+by=0\), distinct indicial roots give \(z^{r_1},z^{r_2}\);
a repeated root gives \(z^r,z^r\log z\), directly verified by
differentiation. For Bessel's equation
\(z^2y''+zy'+(z^2-\nu^2)y=0\), \(\nu\ge0\), the larger-root solution is
\[
 z^\nu\sum_{j=0}^\infty
 \frac{(-1)^j z^{2j}}{4^j j!\prod_{k=1}^j(\nu+k)} .
\]
The empty product is one. Its ratio test proves convergence for all
\(z\), and substitution verifies the equation. This normalization omits
the optional gamma-factor in the customary notation \(J_\nu\).
\end{example}

\begin{proposition}[Polynomial Legendre solutions]\label{d13:09:legendre}
For every integer \(m\ge0\), the equation
\[
 ((1-x^2)y')'+m(m+1)y=0
\]
has a polynomial solution \(P_m\) of degree \(m\), unique up to a nonzero
scalar. Solutions of different degrees are orthogonal in \(L^2(-1,1)\).
\end{proposition}
\begin{proof}
The coefficient relation at an ordinary point is
\((n+2)(n+1)a_{n+2}=(n(n+1)-m(m+1))a_n\).
Starting from the parity of \(m\), choose the nonzero coefficient so that
the recurrence terminates at \(n=m\); starting from the other parity
cannot yield a polynomial. Equivalently solve downward from the leading
coefficient, which also proves uniqueness of a polynomial solution.
For \(m\ne l\), multiply the two equations by \(P_l,P_m\), subtract, and
integrate on \([-1,1]\). The boundary term
\((1-x^2)(P_lP_m'-P_mP_l')\) vanishes at both endpoints, so
\((m(m+1)-l(l+1))\int_{-1}^1P_mP_l=0\).
This is a polynomial identity argument; it does not misclassify the
vanishing endpoint coefficient as a regular Sturm--Liouville problem.
\end{proof}

\paragraph{Source and coverage.}
Teschl, Sections 3.5 and 4.1--4.2, supplies the ordinary and singular
equations. Coddington--Levinson is a classical comparison source.
The resonance and convergence arguments are included rather than
presenting two formal power series without justification.

% END chapters/M13.09


% BEGIN chapters/M13.10
\chapter{Regular Sturm--Liouville problems}\label{ch:d13-10}

Fix a nondegenerate compact interval \([a,b]\). Unless stated otherwise,
\(p\in C^1[a,b]\), \(q,w\in C[a,b]\) are real,
\(p>0\) and \(w>0\). Set
\[
 Ly=-(py')'+qy,\qquad H=L^2([a,b],w(x)\,\dd x).
\]
The inner product is \(\langle u,v\rangle_w=\int_a^bu\overline v\,w\).
For classical solutions the equation is \(Ly=\lambda wy\).
Weighted \(L^2\) classes, their Hilbert completeness and integration rules
are measure-theoretic foundations; the spectral argument needed here is
proved below.

\section{Boundary forms, uniqueness, and comparison}

\begin{definition}[Separated real boundary conditions]\label{d13:10:boundary}
A separated condition is
\[
 \alpha_a y(a)+\beta_a p(a)y'(a)=0,\qquad
 \alpha_b y(b)+\beta_b p(b)y'(b)=0 ,
\]
where each coefficient pair is real and nonzero as a pair.
Dirichlet conditions mean \(y(a)=y(b)=0\). In this chapter the full
eigenfunction-completeness theorem is adopted for Dirichlet conditions;
the boundary-form and orthogonality results also cover the displayed
general separated conditions.
\end{definition}

\begin{proposition}[Lagrange identity]\label{d13:10:lagrange}
For \(C^2\) functions \(u,v\),
\[
 \int_a^b\big((Lu)\overline v-u\overline{Lv}\big)\,\dd x
   =\big[p(u\overline{v'}-u'\overline v)\big]_a^b .
\]
The boundary term vanishes for functions satisfying the same separated
real conditions. Every corresponding eigenvalue is real, and eigenfunctions
for distinct eigenvalues are orthogonal in \(H\).
\end{proposition}
\begin{proof}
Differentiate the expression in square brackets and integrate. At an
endpoint the vectors \((u,pu')\) and \((v,pv')\) lie in the same
one-dimensional complexification of a real line; their displayed
determinant with conjugation therefore vanishes. With \(u=v\) an
eigenfunction, the identity yields
\((\lambda-\overline\lambda)\norm u_H^2=0\). With distinct real eigenvalues
it yields \((\lambda-\mu)\langle u,v\rangle_w=0\).
\end{proof}

\begin{proposition}[Simple zeros and separation]\label{d13:10:separation}
A nonzero real solution has no simultaneous zero of \(y,py'\); its zeros
are simple and isolated. Between consecutive zeros of one solution, every
independent solution of the same equation has exactly one zero.
For \(\mu>\lambda\), between consecutive zeros of a solution at \(\lambda\)
every nonzero real solution at \(\mu\) has a zero.
\end{proposition}
\begin{proof}
Initial-value uniqueness for \((y,py')'=(py'/p,(q-\lambda w)y)\)
prohibits a simultaneous zero. Accumulating zeros in the interval would,
by Rolle's theorem and continuity, give such a zero.
For equal spectral parameters the modified Wronskian
\(p(u'v-uv')\) is constant and nonzero for independent solutions.
On an interval where \(u\ne0\), \((v/u)'\) consequently has a fixed
nonzero sign. At the two successive simple zeros of \(u\), its one-sided
limits have opposite infinite signs, as the same Wronskian formula at
the endpoints shows. Hence \(v/u\) vanishes exactly once.

For the last claim, suppose \(v\) has no zero between successive zeros
\(c,d\) of \(u\), and choose its sign positive there. Direct differentiation
gives Picone's identity
\[
 \big[puu'-p(v'/v)u^2\big]'
 =p\big(u'-(v'/v)u\big)^2+(\mu-\lambda)wu^2 .
\]
The endpoint expression tends to zero at both \(c,d\).
If \(v\) also vanishes at an endpoint, its zero is simple and
\(u^2v'/v=O(|x-c|)\), or its right-endpoint counterpart, still tends to
zero. Integration gives zero equal to a strictly positive quantity,
a contradiction.
\end{proof}

\section{An explicit compact inverse}

\begin{lemma}[Green operator]\label{d13:10:green}
Choose \(c>0\) so that \(Q=q+cw\ge1\) on \([a,b]\).
There is a real continuous symmetric kernel \(G(x,t)\), strictly positive
when \(a<x,t<b\), such that
\[
 (Kf)(x)=\int_a^bG(x,t)f(t)w(t)\,\dd t
\]
is the inverse Dirichlet solution of
\(-(pu')'+Qu=wf\). As an operator \(H\to H\), \(K\) is compact,
self-adjoint, injective, and strictly positive:
\(\langle Kf,f\rangle_w>0\) for \(f\ne0\).
\end{lemma}
\begin{proof}
Let \(u(a)=0,p(a)u'(a)=1\) and \(v(b)=0,p(b)v'(b)=-1\) solve the
homogeneous equation. On the interval to the right of \(a\), \(u\) and
\(pu'\) remain positive: as long as \(u>0\), \((pu')'=Qu>0\),
so a first return to zero is impossible. The reversed-time argument gives
\(v>0\) to the left of \(b\). The constant
\(D=p(u'v-uv')\) is positive, for example by evaluating at \(a\).
Define
\[
 G(x,t)=D^{-1}u(\min(x,t))v(\max(x,t)).
\]
It is continuous, symmetric and positive in the interior, vanishes at
the endpoints, and has derivative jump
\(p(t)(G_x(t+,t)-G_x(t-,t))=-1\).
For continuous \(f\), splitting the integral at \(t=x\), differentiating
the two pieces and using this jump gives the differential equation.
The same split formula for \(f\in H\) gives \(Kf\in C^1\),
\(p(Kf)'\) absolutely continuous, and the equation almost everywhere.
All these integrals are valid because \(w\) is bounded above and below,
so \(f\in L^1[a,b]\). Integration by parts gives
\[
 \langle Kf,f\rangle_w
 =\int_a^b\big(p|(Kf)'|^2+Q|Kf|^2\big)\,\dd x .
\]
This is positive unless \(Kf=0\), in which case the differential equation
gives \(f=0\) almost everywhere. It proves injectivity and positivity.
Symmetry of \(G\) and Fubini prove self-adjointness.
Finally Cauchy--Schwarz bounds \(Kf\) uniformly on the unit ball of \(H\).
Uniform continuity of \(G\) on the square gives, uniformly on that ball,
\[
 |Kf(x)-Kf(y)|
 \le\left(\int_a^b|G(x,t)-G(y,t)|^2w(t)\,\dd t\right)^{1/2}\longrightarrow0
\]
as \(y\to x\). Arzela--Ascoli gives compactness into \(C[a,b]\);
uniform convergence implies convergence in \(H\). Uniqueness of the
inverse follows from the displayed positive energy with \(f=0\).
\end{proof}

\begin{lemma}[Positive compact spectral theorem]\label{d13:10:compact-spectral}
Let \(K\) be a compact self-adjoint positive operator on a Hilbert space.
The closed span of its positive-eigenvalue eigenspaces is
\((\ker K)^\perp\). Every nonzero eigenvalue has finite multiplicity.
There are at most countably many of them; if infinitely many are listed
with multiplicity, they tend to zero.
\end{lemma}
\begin{proof}
For \(m=\sup_{\norm x=1}\langle Kx,x\rangle\), positivity gives the
Cauchy--Schwarz inequality for the positive semidefinite form
\(\langle Kx,y\rangle\): minimize the quadratic
\(\langle K(x+ty),x+ty\rangle\) in a real or complex scalar \(t\).
Consequently
\(\norm{Kx}^2\le m\langle Kx,x\rangle\), by taking the supremum over
unit \(y\). Thus \(\norm K=m\).
If \(m>0\), choose unit \(x_j\) with
\(\langle Kx_j,x_j\rangle\to m\). Then
\[
 \norm{Kx_j-mx_j}^2
 \le m^2-m\langle Kx_j,x_j\rangle\longrightarrow0.
\]
Compactness gives a subsequence \(Kx_j\to z\), and hence
\(x_j\to z/m\). Its limit is a unit eigenvector at \(m\).
The orthogonal complement of an eigenvector is invariant, since \(K\)
is self-adjoint. Repeat this construction there. An infinite orthonormal
family of eigenvectors with eigenvalues bounded below by
\(\varepsilon>0\) would have images at pairwise distance at least
\(\sqrt2\varepsilon\), contradicting compactness. This proves finite
multiplicity and convergence of the successively extracted eigenvalues
to zero; it also gives countability by grouping those above \(1/r\).
On the orthogonal complement of all extracted eigenspaces the restricted
positive compact operator must be zero. If the extraction terminated
after finitely many steps this follows from its stopping rule.
If it did not terminate, a nonzero residual restriction would, by the
maximizing argument, have an eigenvalue \(a>0\). Its eigenvector belongs
to every earlier restricted space, so every greedily extracted maximum
would be at least \(a\), contradicting their convergence to zero.
That complement is therefore exactly \(\ker K\).
\end{proof}

\begin{theorem}[Dirichlet eigenfunction expansion]\label{d13:10:spectrum}
The regular Dirichlet problem has a sequence of simple real eigenvalues
\[
 \lambda_1<\lambda_2<\cdots,\qquad \lambda_j\longrightarrow+\infty,
\]
and real eigenfunctions \(\phi_j\) forming an orthonormal basis of \(H\).
Thus every \(f\in H\) has
\[
 f=\sum_{j=1}^{\infty}\langle f,\phi_j\rangle_w\phi_j
 \quad\hbox{in }H,\qquad
 \norm f_H^2=\sum_{j=1}^{\infty}|\langle f,\phi_j\rangle_w|^2 .
\]
This is an \(L^2\) assertion; pointwise or uniformly convergent expansions
require additional hypotheses.
\end{theorem}
\begin{proof}
Apply the spectral lemma to the injective Green operator. An eigenvector
\(K\phi=\mu\phi\), \(\mu>0\), is continuous and then solves the classical
equation \(L\phi=(\mu^{-1}-c)w\phi\), with the two Dirichlet conditions;
the regularity follows successively from the split Green formula and
the differential equation. Conversely apply the inverse Green operator
to \((L+cw)\phi=(\lambda+c)w\phi\), obtaining
\(\phi=(\lambda+c)K\phi\). Since \(\phi\ne0\), this first excludes
\(\lambda+c=0\). Taking its inner product with \(\phi\) and using strict
positivity of \(K\) then gives \(\lambda+c>0\). Only now divide to obtain
\(K\phi=(\lambda+c)^{-1}\phi\).
The left Dirichlet condition and one initial derivative determine a
solution, so each eigenspace is one-dimensional. Infinitely many are
needed because \(H\) is infinite-dimensional whereas \(K\) is injective.
The eigenvalues of \(K\) decrease to zero, giving the asserted sequence
for \(L\). The orthonormal expansion and Parseval follow by the Hilbert
orthogonal-projection identity and density proved in the spectral lemma.
As the coefficients are real, real and imaginary parts show that real
normalized eigenfunctions can be selected.
\end{proof}

\begin{proposition}[Resolvent alternative]\label{d13:10:resolvent}
For \(f\in H\) and \(\lambda\notin\{\lambda_j\}\), the equation
\((L-\lambda w)y=wf\) with Dirichlet conditions has a unique solution
in the following specified sense: \(y\in C^1\), \(py'\) is absolutely
continuous, and the equation holds almost everywhere. It is given in \(H\) by
\[
 y=\sum_j\frac{\langle f,\phi_j\rangle_w}{\lambda_j-\lambda}\phi_j.
\]
At \(\lambda=\lambda_k\), it is solvable exactly when
\(\langle f,\phi_k\rangle_w=0\), and then the solutions differ by multiples
of \(\phi_k\).
\end{proposition}
\begin{proof}
The equation is equivalent to
\((I-(\lambda+c)K)y=Kf\). In the orthonormal basis just constructed this
is the scalar equation
\((1-(\lambda+c)\mu_j)y_j=\mu_j f_j\).
Away from a spectral value the factors on the left are bounded away from
zero: they tend to one and only finitely many are exceptional.
Thus the basis calculation defines a bounded solution operator.
The resulting identity \(y=K(f+(\lambda+c)y)\) supplies the asserted
regularity and the differential equation. At an eigenvalue exactly one
factor vanishes, proving both the compatibility condition and the kernel
description.
\end{proof}

\section{Oscillation and the order of eigenvalues}

\begin{theorem}[Dirichlet oscillation count]\label{d13:10:oscillation}
The eigenfunction \(\phi_j\) has exactly \(j-1\) zeros in \((a,b)\).
\end{theorem}
\begin{proof}
Let \(y(x,\lambda)\) solve \(y(a,\lambda)=0\), \(p(a)y'(a,\lambda)=1\).
Write
\[
 y=r\sin\theta,\qquad py'=r\cos\theta,\qquad \theta(a,\lambda)=0,
\]
with continuous argument and \(r>0\). Direct differentiation gives
\[
 \theta_x=p^{-1}\cos^2\theta+(\lambda w-q)\sin^2\theta .
\]
At a zero of \(y\), \(\theta_x=1/p>0\), so each crossing of an integer
multiple of \(\pi\) is upward. Differentiable parameter dependence gives,
with \(v=\partial_\lambda y\),
\[
 \partial_\lambda\theta(b,\lambda)
 =\frac{p(y'v-yv')(b)}{r(b,\lambda)^2}
 =\frac{\int_a^bw(x)y(x,\lambda)^2\,\dd x}{r(b,\lambda)^2}>0.
\]
The integral identity follows by differentiating the equation and its
modified Wronskian; its value at \(a\) is zero.

The first eigenfunction can be chosen positive in the interior.
Indeed for a normalized real top eigenvector \(v\) of \(K\),
the strictly positive kernel gives
\(\langle K|v|,|v|\rangle\ge\langle Kv,v\rangle=\norm K\).
The maximizing argument in the spectral lemma shows that \(|v|\) is
also a top eigenvector. Its Green representation makes it strictly
positive inside. Hence \(\theta(b,\lambda_1)=\pi\).
At each other eigenvalue \(\theta(b,\lambda_j)\) is a positive integer
multiple of \(\pi\). These integers strictly increase. None can be skipped:
otherwise continuity and strict monotonicity in \(\lambda\) would give a
Dirichlet eigenvalue between two successive listed eigenvalues. Thus
\(\theta(b,\lambda_j)=j\pi\). Its upward crossings between the endpoints
are precisely \(\pi,\ldots,(j-1)\pi\), proving the count.
\end{proof}

\begin{example}[Constant coefficients]\label{d13:10:dirichlet-example}
For \(p=w=1,q=0,[a,b]=[0,\ell]\),
\(\lambda_j=(j\pi/\ell)^2\) and
\(\phi_j=\sqrt{2/\ell}\sin(j\pi x/\ell)\).
Substitution proves the equation and boundary conditions, and the preceding
theorem identifies their completeness and nodal counts.
By contrast the initial-value problem for \(y''+\lambda y=0\) is unique for
every \(\lambda\), while a homogeneous boundary problem has nonzero
solutions at the discrete eigenvalues.
\end{example}

\paragraph{Source and coverage.}
Teschl, Chapter 5, is the main source for the compact spectral and
Sturm--Liouville route, with Coddington--Levinson as classical comparison.
All endpoints here are regular; singular endpoints and their domain
classification are separate extensions. The adopted complete-basis result
is the general regular Dirichlet case, while general separated conditions
are included in the boundary-form and orthogonality statements.

% END chapters/M13.10


% BEGIN chapters/M13.11
\chapter{Local flows, invariant sets, and first integrals}\label{ch:d13-11}

Let \(\Omega\subset\R^n\) be open and \(f\in C^k(\Omega,\R^n)\),
\(k\ge1\). Write \(I_x\) for the maximal interval through time zero and
\(\Phi(t,x)\) for the solution with initial point \(x\).

\begin{theorem}[A flow has a domain]\label{d13:11:flow}
The set \(\mathcal D=\{(t,x):x\in\Omega,\ t\in I_x\}\) is open,
\(\Phi:\mathcal D\to\Omega\) is \(C^k\), and
\[
 \Phi(0,x)=x,\qquad
 I_{\Phi(s,x)}=I_x-s,\qquad
 \Phi(t,\Phi(s,x))=\Phi(t+s,x)
\]
whenever \(s\in I_x\) and \(t+s\in I_x\).
For each fixed \(t\), \(x\mapsto\Phi(t,x)\) is a \(C^k\) diffeomorphism
from its open domain to the domain of \(\Phi(-t,\cdot)\).
Furthermore
\[
 D_x\Phi(t,x)f(x)=f(\Phi(t,x)).
\]
\end{theorem}
\begin{proof}
Openness and \(C^k\) dependence are the common-domain results of
Chapters~\ref{ch:d13-05}--\ref{ch:d13-06}. Shifting the time of a solution
by \(s\) gives the same initial-value problem at \(\Phi(s,x)\).
Uniqueness gives equality on the common interval. If either shifted
maximal interval were longer, the equality would extend the other solution,
contrary to maximality. This proves the exact domain identity as well.
The group identities give mutually inverse maps at times \(t,-t\).
Differentiate \(\Phi(t,\Phi(s,x))=\Phi(t+s,x)\) at \(s=0\) to obtain the
last formula. In particular no composition is used outside its domain.
\end{proof}

\begin{proposition}[Orbit alternatives]\label{d13:11:orbits}
An equilibrium \(p\) is characterized by \(f(p)=0\), and its solution is
constant. Two maximal orbits are disjoint or identical as sets. A
nonconstant orbit is either injectively parametrized by time or periodic.
In the periodic case its solution is complete and has a least positive
period \(T\); all its periods are \(T\Z\).
\end{proposition}
\begin{proof}
The equilibrium assertion and orbit coincidence follow from uniqueness.
If \(\Phi(t_1,x)=\Phi(t_2,x)\), \(t_1<t_2\), uniqueness repeats the
segment between these times indefinitely in both directions, producing
a complete periodic solution. Its periods form a closed subgroup \(G\)
of \(\R\). If \(G\) had arbitrarily small positive elements, division
with remainder would approximate every real time by elements of \(G\);
continuity would make every time a period and the solution constant.
Thus \(T=\inf(G\cap(0,\infty))>0\). Closedness gives \(T\in G\);
division with remainder then shows \(G=T\Z\).
\end{proof}

\begin{definition}[Invariant sets]\label{d13:11:invariant}
A set is positively invariant if every solution starting in it stays in
it for all nonnegative times in its maximal interval. It is invariant
if this holds for both signs of time. A compact positively invariant
subset of \(\Omega\) supports complete forward solutions, by the
continuation theorem. An invariant compact subset supports complete
two-sided solutions.
\end{definition}

\begin{theorem}[First integrals and regular constraints]\label{d13:11:first-integral}
A \(C^1\) function \(H:\Omega\to\R\) is constant on each solution exactly
when \(DH(x)f(x)=0\) for every \(x\).
More generally let \(h:\Omega\to\R^m\) be \(C^1\), have rank \(m\)
on \(S=h^{-1}(c)\), and satisfy \(Dh(x)f(x)=0\) for \(x\in S\).
Then \(S\) is invariant.
\end{theorem}
\begin{proof}
For the first assertion the chain rule proves sufficiency; differentiating
at time zero proves necessity. For the constraint assertion, a regular
level-set chart parametrizes \(S\) near \(p\in S\) by a \(C^1\) map
\(\psi:V\subset\R^{n-m}\to S\) whose derivative is an isomorphism onto
\(\ker Dh\). Tangency defines a continuous vector field
\(g(u)=(D\psi(u))^{-1}f(\psi(u))\), where the inverse means the
coordinate inverse on that tangent space. Peano existence gives a local
solution of \(u'=g(u)\). Its image under \(\psi\) solves the original
equation and remains in \(S\), so uniqueness for \(f\) forces the original
solution through \(p\) to remain in \(S\) locally. Repeat at each time;
closedness of \(S\) relative to \(\Omega\) and local continuation prove
the assertion on the entire maximal interval, in both directions.
Notice that differentiating \(h\) along an arbitrary nearby curve would
not suffice when tangency is known only on the level set.
\end{proof}

\begin{theorem}[Straightening a nonzero field]\label{d13:11:flowbox}
If \(f(p)\ne0\), a neighborhood of \(p\) has \(C^1\) coordinates
\((s,z)\in(-\varepsilon,\varepsilon)\times V\) in which the field is
\(\partial/\partial s\). For a \(C^k\) field the coordinate map is \(C^k\).
\end{theorem}
\begin{proof}
Choose a linear embedding \(E:\R^{n-1}\to\R^n\) complementary to \(f(p)\).
The map \(\Psi(s,z)=\Phi(s,p+Ez)\) has derivative
\([f(p),E]\) at \((0,0)\), which is invertible.
The inverse function theorem gives the chart. Shrink to a product
neighborhood whose closure remains in the chart's domain.
The flow law gives
\(\Phi(t,\Psi(s,z))=\Psi(s+t,z)\) whenever both chart points are present;
this is precisely the asserted constant field. Its regularity follows
from the regularity of \(\Phi\) and the inverse function theorem.
\end{proof}

\begin{theorem}[Compact omega-limit sets]\label{d13:11:omega}
Suppose a positive orbit has compact closure \(K\subset\Omega\).
Its omega-limit set
\[
 \omega(x)=\bigcap_{s\ge0}\overline{\{\Phi(t,x):t\ge s\}}
\]
is nonempty, compact, connected and invariant. Its points are exactly
limits \(\Phi(t_j,x)\) with \(t_j\to\infty\), and
\(\operatorname{dist}(\Phi(t,x),\omega(x))\to0\).
\end{theorem}
\begin{proof}
The forward solution is complete by continuation inside \(K\).
The closed tails are nested nonempty compact connected sets:
each is the closure of a continuous image of a connected interval.
Compactness gives nonempty intersection. If the intersection were
disconnected, disjoint open neighborhoods of its two compact parts
would eventually contain a closed tail, contradicting connectedness;
otherwise points outside these neighborhoods in arbitrarily late
tails would have a limit in the intersection. This proves connectedness.
The sequence characterization follows from choosing time at least \(j\)
within distance \(1/j\) of a point in the intersection; its converse
is immediate.

For \(y=\lim\Phi(t_j,x)\), the continuous-dependence theorem gives
\(\Phi(s,y)=\lim\Phi(t_j+s,x)\) on a uniform small time interval for
both signs of \(s\). The latter points stay in \(K\), hence the shifted
limits belong to \(\omega(x)\). Local intervals can be continued for
all time because the orbit remains in the compact set \(K\).
This proves invariance and completeness on \(\omega(x)\).
Finally, failure of the distance assertion would give a sequence of
late orbit points a fixed distance from \(\omega(x)\); a convergent
subsequence in \(K\) contradicts the sequence characterization.
\end{proof}

\begin{proposition}[Hamiltonian energy and volume]\label{d13:11:hamiltonian}
For \(H\in C^2(\Omega)\), \(\Omega\subset\R^{2m}\), the system
\(q'=\partial_pH,\ p'=-\partial_qH\) preserves \(H\).
Its local time maps have Jacobian determinant one, and therefore
preserve Lebesgue measure on their open domains and measurable subsets
thereof.
\end{proposition}
\begin{proof}
The derivative of \(H\) is
\(\partial_qH\cdot\partial_pH-\partial_pH\cdot\partial_qH=0\).
Equality of mixed derivatives gives \(\operatorname{div}f=0\).
The variational matrix \(D_x\Phi(t,x)\) solves
\(Y'=Df(\Phi(t,x))Y,\ Y(0)=I\); Liouville's formula yields determinant
one. Apply the change-of-variables theorem to the local diffeomorphism
from Theorem~\ref{d13:11:flow}. This makes no assertion that an incomplete
flow is defined globally.
\end{proof}

\paragraph{Source and coverage.}
The flow, orbit, transverse-coordinate and invariant-set results follow
the material of Teschl, Sections 6.1--6.4 and 8.1; first integrals also
connect with its Hamiltonian discussion. The proofs expose the domain
and tangency issues needed to compose interfaces in later chapters.

% END chapters/M13.11


% BEGIN chapters/M13.12
\chapter{Equilibria and Lyapunov stability}\label{ch:d13-12}

Translate an equilibrium of \(x'=f(x)\), \(f\in C^1(\Omega,\R^n)\), to
the origin. All neighborhoods below have closure in \(\Omega\).

\begin{definition}[Stability notions]\label{d13:12:definitions}
The origin is Lyapunov stable if for every \(\varepsilon>0\) some
\(\delta>0\) ensures \(\norm{x_0}<\delta\) implies
\(\norm{\Phi(t,x_0)}<\varepsilon\) for every \(t\ge0\).
It is locally asymptotically stable if it is stable and all sufficiently
nearby solutions tend to zero. It is locally exponentially stable if
some \(C,\alpha,\delta>0\) satisfy
\(\norm{\Phi(t,x_0)}\le Ce^{-\alpha t}\norm{x_0}\) for \(t\ge0\)
and \(\norm{x_0}<\delta\). These assertions include forward existence.
\end{definition}

\begin{theorem}[Lyapunov's direct method]\label{d13:12:lyapunov}
Let \(V\in C^1\) near zero, \(V(0)=0\), \(V(x)>0\) for \(x\ne0\),
and \(\dot V(x)=DV(x)f(x)\le0\). Then the origin is stable.
If \(\dot V(x)<0\) for \(x\ne0\), it is asymptotically stable.
If, on a neighborhood,
\[
 c_1\norm x^2\le V(x)\le c_2\norm x^2,\qquad
 \dot V(x)\le-c_3\norm x^2,\qquad c_i>0,
\]
it is exponentially stable with estimate
\[
 \norm{\Phi(t,x_0)}
 \le\sqrt{c_2/c_1}\,e^{-c_3t/(2c_2)}\norm{x_0}
\]
for sufficiently small \(x_0\).
\end{theorem}
\begin{proof}
Given a small sphere of radius \(r\), the minimum \(m_r\) of \(V\) on
that sphere is positive. Continuity at zero supplies \(\delta\) with
\(V(x_0)<m_r/2\) when \(\norm{x_0}<\delta\).
While the orbit stays in the ball, \(V\) cannot increase; a first exit
through the sphere is consequently impossible. The closed sublevel in
the ball is compact and positively invariant, so continuation gives
existence for all positive time.
For strict \(\dot V<0\), \(V(\Phi(t,x_0))\) decreases to a limit.
On the omega-limit set it has that constant value, and invariance makes
its derivative zero there. Hence the omega-limit set is \(\{0\}\);
Theorem~\ref{d13:11:omega} gives convergence.
For the last assertion,
\(\dot V\le-(c_3/c_2)V\). Integrating this scalar differential inequality
and using the two quadratic bounds gives the displayed estimate.
The initial point is chosen small enough that the preceding invariant
sublevel lies in the neighborhood where all inequalities hold.
\end{proof}

\begin{lemma}[The Lyapunov matrix equation]\label{d13:12:matrix-lyapunov}
A real matrix \(A\) has all eigenvalues in the open left half-plane
if and only if, for any prescribed symmetric positive definite \(Q\),
there is a symmetric positive definite \(P\) with
\[
 A^{\mathsf T}P+PA=-Q .
\]
When it exists, it is unique and is
\(P=\int_0^\infty e^{tA^{\mathsf T}}Qe^{tA}\,\dd t\).
\end{lemma}
\begin{proof}
The exponential estimate of Chapter~\ref{ch:d13-07} proves convergence
of the integral. It is symmetric and positive definite: for \(x\ne0\)
the continuous integrand \(x^{\mathsf T}e^{tA^{\mathsf T}}Qe^{tA}x\)
is positive at \(t=0\). Integrate its matrix derivative to get the
equation; the boundary at infinity is zero.
If \(A^{\mathsf T}P+PA=-Q\) and \(Av=\lambda v\) with complex \(v\ne0\),
then \(2\Re\lambda\,v^*Pv=-v^*Qv<0\), so every real part is negative.
For uniqueness, a difference \(D\) solves \(A^{\mathsf T}D+DA=0\).
The matrix \(e^{tA^{\mathsf T}}De^{tA}\) is constant but tends to zero,
so \(D=0\).
\end{proof}

\begin{theorem}[Linearized stability and instability]\label{d13:12:linearized-stability}
Let \(A=Df(0)\).
If every eigenvalue of \(A\) has negative real part, the equilibrium is
locally exponentially stable. If some eigenvalue has positive real part,
the equilibrium is unstable. Eigenvalues on the imaginary axis alone
do not decide either conclusion.
\end{theorem}
\begin{proof}
Write \(f(x)=Ax+g(x)\) with \(g(x)=o(\norm x)\).
In the stable case take \(Q=I\) and the preceding \(P\), and put
\(V=x^{\mathsf T}Px\). Its quadratic bounds hold by positive
definiteness, and
\(\dot V=-\norm x^2+2x^{\mathsf T}Pg(x)\le-\norm x^2/2\)
on a small ball. The direct method gives the result.

For instability use the real invariant spectral splitting
\(\R^n=E_u\oplus E_{cs}\), where the eigenvalues on \(E_u\) have
positive real parts and those on \(E_{cs}\) have nonpositive real parts.
Choose \(0<\eta<\alpha\) smaller than every positive real part.
The integral construction above, applied to \(-A_u+\alpha I\) and
\(A_{cs}-\eta I\), gives positive quadratic forms \(U(u),W(v)\) with
\[
 \dot U\ge2\alpha U,\qquad \dot W\le2\eta W
\]
for the linear system. The missing space \(E_{cs}=\{0\}\) is allowed,
with \(W=0\). For the nonlinear system, equivalence of quadratic norms
and \(g=o(\norm x)\) gives an error at most
\(\varepsilon C(U+W)\) in the derivative of \(V=U-W\), on a small
fixed ball. In the cone \(U>W\),
\[
 \dot V\ge2\alpha U-2\eta W-\varepsilon C(U+W)
          \ge c(U+W)\ge cV>0
\]
for some \(c>0\), after choosing \(\varepsilon\) sufficiently small.
Initial data \((u,0)\), \(u\ne0\), may be arbitrarily small.
While inside the ball their \(V\) stays positive, so they cannot leave
the cone through \(V=0\), and \(V(t)\ge e^{ct}V(0)\).
They cannot remain in the ball forever, since \(V\) is bounded there;
nor can a finite maximal time occur without leaving its compact closure.
Thus arbitrarily small initial points leave a fixed neighborhood,
which proves instability.
\end{proof}

\begin{example}[The inconclusive zero eigenvalue]\label{d13:12:zero-example}
All three scalar equations \(x'=-x^3,\ x'=x^3,\ x'=0\) have derivative
zero at the equilibrium. The first has
\(x(t)=x_0(1+2t x_0^2)^{-1/2}\), hence asymptotic stability; the second
has finite-time escape for each positive nonzero initial value and is
unstable; the last is stable but not attractive. Linearization must not
be used to classify these nonhyperbolic cases.
\end{example}

\begin{theorem}[LaSalle's invariance principle on a compact set]\label{d13:12:lasalle}
Let \(K\subset\Omega\) be compact and positively invariant.
Suppose \(V\in C^1\) on a neighborhood of \(K\) and \(\dot V\le0\) on \(K\).
Let \(E=\{x\in K:\dot V(x)=0\}\), and let \(M\) be the largest invariant
subset of \(E\), meaning the union of complete orbits contained in \(E\).
For each \(x\in K\),
\(\operatorname{dist}(\Phi(t,x),M)\to0\).
If \(M\) is a singleton equilibrium, every such orbit converges to it.
\end{theorem}
\begin{proof}
Along the orbit, \(V\) is decreasing and bounded below by its minimum
on \(K\), hence has limit \(\ell\). At each point of \(\omega(x)\)
continuity and the sequence characterization give \(V=\ell\).
The entire orbit of such a point remains in \(\omega(x)\), so its
derivative of \(V\) is zero. Thus \(\omega(x)\subset M\).
Since \(E\) is compact, the set \(M\) is also compact: a limit of points
whose complete orbits remain in \(E\) has, by continuous dependence
and compact continuation, its entire orbit in \(E\).
Distance to \(\omega(x)\) tends to zero, hence so does distance to \(M\).
The singleton assertion follows immediately.
\end{proof}

\begin{example}[A damped mechanical equilibrium]\label{d13:12:damped-pendulum}
For \(x'=y,\ y'=-\sin x-cy\), \(c>0\), the energy
\(V=y^2/2+1-\cos x\) satisfies \(\dot V=-cy^2\).
In a small compact sublevel around \((0,0)\), the only complete orbit
contained in \(\{y=0\}\) is the equilibrium: staying there also requires
\(\sin x=0\). The direct method and LaSalle therefore give local
asymptotic stability. In contrast the undamped equation conserves \(V\),
so its nearby nonzero-energy solutions cannot converge to \((0,0)\).
\end{example}

\paragraph{Source and coverage.}
Teschl, Sections 6.5--6.7 and 9.1, and the local-theory part of Perko
provide the source framework. The compactness assumptions in the
invariance argument and the distinction between stability, attraction
and exponential decay are retained. Hyperbolic invariant manifolds
and topological linearization are developed separately in
Chapter~\ref{ch:d13-20}; they are not prerequisites for the direct
stability proofs above.

% END chapters/M13.12


% BEGIN chapters/M13.13
\chapter{Planar dynamics and limit sets}\label{ch:d13-13}

In this chapter \(\Omega\subset\R^2\) is open and \(f\in C^1(\Omega,\R^2)\).
The dimension and regularity are part of every limit-set statement.
A transverse segment means a sufficiently short straight segment on which
the normal component of \(f\) never vanishes. Its points are ordered by a
fixed affine coordinate.

\section{Nullclines and phase regions}
\begin{definition}[Coordinate nullclines]\label{d13:13:nullclines}
For $x'=F(x,y)$, $y'=G(x,y)$, the coordinate nullclines are
$N_x=\{(x,y):F(x,y)=0\}$ and $N_y=\{(x,y):G(x,y)=0\}$.
Equilibria are exactly $N_x\cap N_y$. Away from both nullclines, the signs
of $F,G$ specify the quadrant of the tangent vector. At a point of
$N_x\setminus N_y$ the tangent is vertical and nonzero; at a point of
$N_y\setminus N_x$ it is horizontal and nonzero. A nullcline need not
itself be invariant. These are coordinate-dependent descriptions of the
vector field, not additional conservation laws.
\end{definition}

\begin{example}[A complete sign chart in the positive quadrant]
\label{d13:13:nullcline-example}
Consider $x'=x(1-x-y)$, $y'=y(x-a)$ with $0<a<1$.
The nullclines are $x=0$ or $y=1-x$ for the first component, and
$y=0$ or $x=a$ for the second. In $x,y>0$, the first component is
positive exactly below $y=1-x$; the second is positive exactly to the
right of $x=a$. Combining the two comparisons gives all four sign
regions. The equilibria in the closed positive quadrant are
$(0,0)$, $(1,0)$ and $(a,1-a)$, obtained by intersecting the listed
zero sets. In particular most points on either nullcline are not
equilibria. Positive coordinates stay positive while a solution exists:
along that solution,
\[
 x(t)=x(0)\exp\!\left(\int_0^t(1-x(s)-y(s))\,\dd s\right),\quad
 y(t)=y(0)\exp\!\left(\int_0^t(x(s)-a)\,\dd s\right).
\]
The formulas follow by scalar integrating factors and also show the
invariance of either axis when its initial coordinate is zero.
They justify the phase-region interpretation rather than relying on a
numerically drawn picture.
\end{example}

\section{The planar crossing argument}

We use the integer winding index of a closed piecewise \(C^1\) curve
about a point off its trace, with its local constancy on the complement.
These are complex-analysis foundations. The following local calculation
records the precise crossing property needed below, without assuming
that all components of a curve's complement have already been classified.

\begin{lemma}[Index jump across a straight edge]\label{d13:13:index-jump}
Suppose a closed piecewise \(C^1\) curve meets a small disc in one
oriented straight arc and has no other points in that disc. The winding
index immediately to the left of the oriented arc is one greater than
the index immediately to its right.
\end{lemma}
\begin{proof}
Translate and rotate an interior point of the arc to zero and orient
the arc from \(-\varepsilon\) to \(\varepsilon\) on the real axis.
Its contribution to the index about \(iy\) is
\[
 \frac1{2\pi}\int_{-\varepsilon}^{\varepsilon}
             \frac{y}{s^2+y^2}\,\dd s .
\]
The limits as \(y\downarrow0\) and \(y\uparrow0\) are respectively
\(1/2\) and \(-1/2\). The rest of the curve stays away from zero, so its
contribution is continuous in the reference point and has the same
two limits. On each sufficiently small half-disc the index is constant
and integer-valued. Their difference is therefore exactly one, as the
limiting calculation shows. Reversal of orientation reverses the jump.
\end{proof}

\begin{lemma}[Monotone intersections]\label{d13:13:monotone-crossings}
Successive intersections of a forward orbit with a transverse segment
are either all the same point, making the orbit periodic, or are strictly
monotone in the segment coordinate. In particular an omega-limit set
intersects a sufficiently short transverse segment in at most one point.
\end{lemma}
\begin{proof}
Make an affine coordinate change so that the segment is horizontal
and \(f\) crosses it upwards. Let \(p=(a,0)\), \(q=(b,0)\) be consecutive
distinct intersections, with \(a<b\). Join the orbit arc from \(p\) to
\(q\) to the segment from \(q\) back to \(p\), forming a closed curve \(J\).
The orbit arc is injective: a self-return would make the solution
periodic and would force a return to \(p\) before the purported next
intersection \(q\). Its only intersections with the segment are its
endpoints. Near either endpoint a flowbox makes the two incident arcs
the two sides of a right angle.

Let \(c\) be the index just above the open segment between \(p,q\).
The index just below it is \(c+1\), by the preceding lemma, since the
closing segment is oriented to the left. These values are consistent
along the segment: every compact subsegment has a thin strip disjoint
from the orbit arc, and overlapping strips have the same index.
At \(p\), the orbit leaves upwards and the closing segment approaches
from the right; a path around the left side of this local corner joins
the segment to the left of \(p\) to the lower strip. Hence the left
part of the transverse segment has index \(c+1\).
At \(q\), the orbit arrives from below and the closing segment leaves
to the left. A path around the upper right corner joins the segment
to the right of \(q\) to the upper strip, so that part has index \(c\).
These statements use only the two local arcs in a flowbox; compactness
and injectivity let the neighborhoods avoid the rest of \(J\).

Immediately after \(q\), the orbit has index \(c\). It cannot cross the
open closing segment: an upward crossing would have to approach from
index \(c+1\). Nor can it first meet the interior of the old orbit arc,
since uniqueness would make an interval immediately before that meeting
part of the same arc already. A first meeting with \(q\) has the same
backward-uniqueness obstruction, and a first meeting with \(p\) would
approach from its lower side, of index \(c+1\). Thus until another
transverse intersection it avoids \(J\) and retains index \(c\).
It cannot meet the segment left of \(p\), whose index is \(c+1\), or
between \(p,q\); its next intersection must lie to the right of \(q\).
Induction proves strict monotonicity. For \(b<a\) reverse the segment
coordinate. A repeated intersection point gives periodicity by
uniqueness.

If \(z\) is an omega-limit point on the interior of the segment,
nearby late orbit points can be moved a time tending to zero onto the
segment, using its flowbox. Their new times still tend to infinity,
and the resulting intersections tend to \(z\).
An infinite sequence of intersections cannot accumulate at a finite
time, since a flowbox at the limiting regular point permits only one
crossing in a uniform small time neighborhood. The chronologically
ordered intersections are therefore the monotone sequence just proved,
and have at most one limit in the segment. Two different omega-limit
points there are impossible.
\end{proof}

\begin{theorem}[Poincare--Bendixson, regular compact case]\label{d13:13:poincare-bendixson}
Suppose a positive orbit has compact closure in \(\Omega\), and its
omega-limit set contains no equilibrium. That omega-limit set is one
nonconstant periodic orbit.
\end{theorem}
\begin{proof}
Let \(K=\omega(x)\); it is nonempty compact connected and invariant.
Choose \(y\in K\), and then \(z\in\omega(y)\subset K\).
The latter set is nonempty because the orbit of \(y\) remains in \(K\).
Since \(f(z)\ne0\), choose a transverse segment through \(z\).
Moving late orbit points onto it as in the preceding proof gives
intersections of the orbit of \(y\) converging to \(z\).
Every such intersection belongs to \(K\), and the monotone-intersection
lemma says \(K\) meets this segment only at \(z\). Thus the orbit of
\(y\) passes through \(z\) at two different times and is periodic.
Call this compact orbit \(\Gamma\subset K\).

It remains to show equality, not merely the existence of a periodic
orbit in the limit set. For \(p\in\Gamma\), use a flowbox with a transverse
segment through \(p\). Any point of \(K\) sufficiently close to \(p\)
can be moved onto that segment by a small time. Invariance keeps the
image in \(K\), so it equals \(p\) by the one-intersection property.
The original point consequently belongs to \(\Gamma\).
Thus \(\Gamma\) is relatively open in \(K\); compactness makes it closed.
Connectedness and nonemptiness of \(K\) give \(K=\Gamma\).
\end{proof}

\begin{corollary}[A trapping set without equilibria]\label{d13:13:trapping}
A nonempty compact positively invariant subset \(K\subset\Omega\)
containing no equilibrium contains a periodic orbit.
\end{corollary}
\begin{proof}
Start at any point of \(K\) and apply the theorem to its omega-limit set.
\end{proof}

\section{Exclusion and representative portraits}

\begin{theorem}[Bendixson--Dulac exclusion]\label{d13:13:dulac}
Let \(D\subset\Omega\) be open and simply connected and let \(B\in C^1(D)\).
If \(\operatorname{div}(Bf)\) is everywhere strictly positive on \(D\),
or everywhere strictly negative, no nonconstant periodic orbit is
contained in \(D\). The choice \(B=1\) is Bendixson's criterion.
\end{theorem}
\begin{proof}
A hypothetical nonconstant periodic orbit has nonzero tangent \(f\)
and is a simple regular closed curve when parametrized by its least
period. The planar separation results
Theorems~\ref{d13:21:smooth-jordan} and~\ref{d13:21:domain-interior}
give a bounded region \(V\) with that boundary and compact closure
contained in \(D\). Green's theorem gives
\[
 \iint_V\operatorname{div}(Bf)\,\dd x\,\dd y
       =\int_{\partial V} Bf\cdot n\,\dd s=0 ,
\]
because \(f\) is tangent to the boundary. The area integral has a strict
nonzero sign, a contradiction. Thus the topological condition on \(D\)
and the interior containment used by Green's theorem are explicit.
\end{proof}

\begin{example}[Why the domain matters]\label{d13:13:annulus}
For
\[
 x'=(1-x^2-y^2)x-y,\qquad
 y'=x+(1-x^2-y^2)y ,
\]
polar coordinates on \(r>0\) give \(r'=r(1-r^2)\), \(\theta'=1\).
The annulus \(1/2\le r\le2\) is positively invariant: on its inner
boundary the radial derivative is positive and on its outer boundary
negative, precluding a first exit. It contains no equilibria and contains
the periodic orbit \(r=1\), of period \(2\pi\).
Every nonzero solution tends to that circle, since the scalar radial
equation is monotone toward its only positive equilibrium.
\end{example}

\begin{example}[The pendulum and an equilibrium in the limit set]\label{d13:13:pendulum}
For \(x'=y,\ y'=-\sin x\), the energy
\(E=y^2/2+1-\cos x\) is constant.
At \((2k\pi,0)\) its Hessian is positive definite, giving Lyapunov-stable
centers. At \(((2k+1)\pi,0)\) the linearization has eigenvalues \(1,-1\),
so those equilibria are unstable. For \(0<E<2\), the compact energy oval
surrounding a center is regular and invariant; the preceding theorem
shows it is a periodic orbit. At energy \(2\), a branch can approach
the saddle instead. The hypothesis excluding equilibria from the
omega-limit set cannot be deleted.
\end{example}

\paragraph{Source and coverage.}
The principal source is Teschl, Section 7.3, with its mechanical and
planar examples in Sections 6.7 and 7.1--7.2; Perko's global theory is
a comparison source. The monotone-crossing proof above records the
winding-index barrier explicitly. The generalized classification of
limit sets containing arbitrary equilibrium connections is not
silently asserted by the regular compact theorem.

% END chapters/M13.13


% BEGIN chapters/M13.14
\chapter{Periodic orbits and return maps}\label{ch:d13-14}

Let \(f\in C^k(\Omega,\R^n)\), \(k\ge1\), and let
\(p\) lie on a nonconstant periodic orbit \(\Gamma\) of least period \(T>0\).
Then \(f(p)\ne0\). Choose a linear functional \(\ell\) with
\(\ell(f(p))\ne0\) and a small local section
\(\Sigma=\{z:\ell(z-p)=0\}\) around \(p\).

\begin{theorem}[Local return time and map]\label{d13:14:return-map}
On an open neighborhood of \(p\) in \(\Sigma\) there is a unique \(C^k\)
function \(\tau(z)\) close to \(T\) such that
\(\tau(p)=T\) and
\[
 P(z)=\Phi(\tau(z),z)\in\Sigma .
\]
After shrinking the neighborhoods, \(P\) is a \(C^k\) local
diffeomorphism fixing \(p\). The section can be chosen so that this is
the first positive return for the points under consideration.
Writing \(M=D_x\Phi(T,p)\), for \(v\in\ker\ell\),
\[
 D\tau(p)v=-\frac{\ell(Mv)}{\ell(f(p))},\qquad
 DP(p)v=Mv-f(p)\frac{\ell(Mv)}{\ell(f(p))}.
\]
\end{theorem}
\begin{proof}
Apply the implicit function theorem to
\(F(t,z)=\ell(\Phi(t,z)-p)\) at \((T,p)\):
\(\partial_tF(T,p)=\ell(f(p))\ne0\).
The exact common-domain flow theorem supplies the required open
neighborhood of \((T,p)\). Its derivative formulas give the displayed
identities, and the second one maps \(\ker\ell\) to itself.
The analogous implicit return time near \(-T\) constructs an inverse:
for \(z\) near \(p\), time \(-\tau(z)\), close to \(-T\), takes
\(P(z)\) back to \(z\); uniqueness of the implicit branch proves
the inverse identity on smaller neighborhoods.

For the first-return statement, a flowbox at \(p\) excludes a second
intersection within a uniform short initial or terminal passage.
For any fixed small \(\varepsilon>0\), the compact arc
\(\{\Phi(t,p):\varepsilon\le t\le T-\varepsilon\}\) does not contain
\(p\), because \(T\) is the least period. Choose the section inside a
ball disjoint from that arc. Uniform continuous dependence on this
compact time interval keeps nearby orbit arcs disjoint from that
section as well. The implicit return is therefore the first one.
These are local conclusions; arbitrary points of a large hyperplane
need not return at all.
\end{proof}

\begin{theorem}[The neutral multiplier and transverse multipliers]\label{d13:14:multipliers}
The variational equation along \(\Gamma\) has \(T\)-periodic coefficients
\(A(t)=Df(\Phi(t,p))\). Its monodromy \(M\) satisfies \(Mf(p)=f(p)\).
The operator induced by \(M\) on
\(\R^n/\operatorname{span}\{f(p)\}\) is conjugate to \(DP(p)\).
Consequently the characteristic polynomial factors as
\[
 \det(\lambda I-M)=(\lambda-1)\det(\lambda I-DP(p)),
 \qquad
 \det DP(p)=\exp\left(\int_0^T\operatorname{div}f(\Phi(t,p))\,\dd t\right).
\]
For a planar orbit its single transverse multiplier is therefore
strictly positive.
\end{theorem}
\begin{proof}
Periodicity of \(A\) follows from periodicity of the orbit.
Differentiating the flow law in its time-shift variable gives
\(D_x\Phi(T,p)f(p)=f(\Phi(T,p))=f(p)\).
The map \(\ker\ell\to\R^n/\operatorname{span}\{f(p)\}\),
\(v\mapsto[v]\), is an isomorphism.
The derivative formula for \(P\) subtracts only a multiple of \(f(p)\),
so under this isomorphism it induces exactly \([v]\mapsto[Mv]\).
In a basis starting with \(f(p)\), the matrix of \(M\) is block
triangular with diagonal blocks \(1,DP(p)\).
Take characteristic polynomials and determinants, and then apply
Liouville's formula to the variational equation.
\end{proof}

\begin{definition}[Orbital stability]\label{d13:14:orbital-definition}
Orbital stability means that solutions starting sufficiently close to
the compact set \(\Gamma\) remain close to that set for all future time.
Orbital asymptotic stability additionally requires their distance to
\(\Gamma\) to tend to zero. No claim is made that they tend to one
fixed point of the moving periodic orbit at the same time phase.
\end{definition}

\begin{theorem}[Transverse contraction implies orbital attraction]\label{d13:14:orbital-stability}
If every eigenvalue of \(DP(p)\) has modulus less than one, then
\(\Gamma\) is locally exponentially orbitally stable:
some \(C,\alpha,\delta>0\) satisfy
\[
 \operatorname{dist}(\Phi(t,x),\Gamma)
 \le C e^{-\alpha t}\operatorname{dist}(x,\Gamma),\qquad t\ge0,
\]
whenever \(\operatorname{dist}(x,\Gamma)<\delta\).
\end{theorem}
\begin{proof}
Put \(B=DP(p)\), choose \(0<r<\beta<1\) with
\(\norm{B^j}\le C_0r^j\) (Jordan's formula), and define on the section
\[
 \norm v_*=\sum_{j=0}^{\infty}\beta^{-j}\norm{B^jv}.
\]
This converges and is equivalent to the original norm, with
\(\norm{Bv}_*\le\beta\norm v_*\).
Continuity of \(DP\) therefore makes \(P\) a contraction with constant
\(\theta<1\) on a sufficiently small convex section ball centered at
\(p\), mapped into itself. Hence its iterates \(z_j\) obey
\(\norm{z_j-p}_*\le\theta^j\norm{z_0-p}_*\).
Shrink again so that \(T/2<\tau(z)<3T/2\).
Uniform differentiable dependence of the flow on the compact time
interval \([0,3T/2]\) gives
\[
 \operatorname{dist}(\Phi(s,z_j),\Gamma)
 \le L\norm{z_j-p},\quad 0\le s\le\tau(z_j).
\]
At cumulative return times \(t_j=\sum_{i<j}\tau(z_i)\), concatenate
these solution segments. They remain in a fixed compact neighborhood
of \(\Gamma\) inside \(\Omega\), so continuation is legitimate.
Since \(t_j\le3jT/2\), the geometric bound becomes an exponential
in actual time, with a constant independent of \(j\).

Finally this section estimate applies to all nearby initial points.
For \(q=\Phi(s,p)\), \(0\le s\le T\), the orbit of \(q\) returns to
\(p\) at time \(2T-s\in[T,2T]\). The same implicit-function argument
gives nearby points a nearby positive crossing time and a projection
onto \(\Sigma\). Compactness of the parameter interval gives finitely
many such charts with uniformly bounded derivatives, entry times and
flow derivatives before entry. If \(q\in\Gamma\) is nearest to \(x\),
the section projection \(z\) satisfies
\(\norm{z-p}\le L_1\norm{x-q}
=L_1\operatorname{dist}(x,\Gamma)\), while the pre-entry orbit stays
within \(L_2\operatorname{dist}(x,\Gamma)\) of \(\Gamma\).
The bounded entry time can be absorbed into the exponential constant.
This proves the stated distance estimate and stability.
\end{proof}

\begin{proposition}[An unstable transverse multiplier]\label{d13:14:orbital-instability}
If \(DP(p)\) has an eigenvalue of modulus greater than one,
\(\Gamma\) is orbitally unstable.
For a planar orbit, the criteria reduce to attraction when
\(\int_0^T\operatorname{div}f\,\dd t<0\) and instability when it is positive.
Equality leaves the nonlinear question undecided.
\end{proposition}
\begin{proof}
The discrete linearized-instability theorem
(\ref{d13:15:linearized-instability}) supplies initial section points
arbitrarily close to \(p\) whose return-map iterates leave a fixed small
section ball. Choose that ball with its closure in the return-map
domain and with \(\Sigma\cap\Gamma=\{p\}\).
By continuity of \(P\), the first exiting iterates lie in a fixed
slightly larger section neighborhood. The compact section annulus
between these two radii has positive distance from \(\Gamma\).
At the corresponding return times the original orbit is therefore
a fixed positive distance from \(\Gamma\), proving instability.
The planar statements follow from the positive multiplier formula.
\end{proof}

\begin{example}[A computable multiplier]\label{d13:14:radial-example}
For the oscillator of Example~\ref{d13:13:annulus}, \(\Gamma=\{r=1\}\)
has period \(2\pi\). The radial variation satisfies
\(\rho'=(1-3r^2)\rho=-2\rho\) on that orbit, so its transverse multiplier
is \(e^{-4\pi}<1\). The tangential multiplier is \(1\).
The orbit is attracting although the full monodromy cannot have all
multipliers inside the unit circle. This is why equilibrium and
periodic-orbit stability criteria must not be confused.
\end{example}

\paragraph{Source and coverage.}
Teschl, Sections 6.4 and 12.1--12.2, and the return-map discussion in
Hirsch--Smale--Devaney provide the main source material. The local nature
of sections, common existence domains, neutral time direction and
conversion from iteration count to real time are made explicit.

% END chapters/M13.14


% BEGIN chapters/M13.15
\chapter{Discrete dynamical systems}\label{ch:d13-15}

\section{Iteration, invariance and changes of coordinates}

\begin{definition}[Discrete system and periodic point]\label{d13:15:discrete-system}
A discrete dynamical system is a map \(F:X\to X\); its forward orbit at \(x\)
is \((F^n x)_{n\geq0}\), with \(F^0=\id\) and \(F^{n+1}=F\circ F^n\).
A set \(A\) is positively invariant if \(F(A)\subset A\), and invariant if
\(F(A)=A\). A point has least period \(p\geq1\) if \(F^p x=x\) and no
smaller positive integer has this property. Negative iterates are used only
when \(F\) is invertible.
\end{definition}

\begin{definition}[Conjugacy and semiconjugacy]\label{d13:15:conjugacy-definition}
Maps \(F:X\to X\), \(G:Y\to Y\) are topologically conjugate if a
homeomorphism \(h:X\to Y\) satisfies \(hF=Gh\). A continuous surjection
satisfying that equation is a semiconjugacy. Injectivity is part of the
first definition and is not part of the second.
\end{definition}

\begin{proposition}[Transport by a conjugacy]\label{d13:15:conjugacy-transport}
If \(hF=Gh\), then \(hF^n=G^nh\) for every \(n\geq0\). For a conjugacy,
periodic points and their least periods, invariant sets and dense forward
orbits correspond. At a fixed point, Lyapunov stability and asymptotic
stability are preserved by a conjugacy.
\end{proposition}
\begin{proof}
The iterate identity follows by induction. Bijectivity permits the same
identity for \(h^{-1}\), so equality \(F^nx=x\) is equivalent to
\(G^nh(x)=h(x)\), including the absence of smaller periods. Images and
inverse images give the invariant-set assertion. A homeomorphism maps
closures to closures, proving the dense-orbit assertion. For stability,
use neighborhoods: given a neighborhood \(V\) of \(h(x)\), stability at
\(x\) supplies a neighborhood \(U\) with every \(F^nU\subset h^{-1}V\).
Then \(h(U)\) has every \(G^nh(U)\subset V\). The reverse implication uses
\(h^{-1}\). Convergence of an orbit to the fixed point is preserved by
continuity in both directions, proving the asymptotic assertion.
\end{proof}

\begin{proposition}[An invariant interval]\label{d13:15:interval-fixed-point}
Every continuous map \(F:[a,b]\to[a,b]\), \(a<b\), has a fixed point.
If it is Lipschitz with constant \(q<1\), the fixed point \(p\) is unique,
every forward orbit converges to it, and
\(\abs{F^nx-p}\leq q^n\abs{x-p}\).
\end{proposition}
\begin{proof}
The continuous function \(F(x)-x\) is nonnegative at \(a\) and nonpositive
at \(b\); apply the intermediate value theorem. If \(p,p'\) are fixed,
\(\abs{p-p'}\leq q\abs{p-p'}\), forcing equality. Repeated use of the
Lipschitz inequality gives the displayed estimate and convergence.
\end{proof}

\section{Stability and multipliers}

\begin{definition}[Stability for a map]\label{d13:15:stability-definition}
A fixed point \(p\) of a continuous map on a metric space is Lyapunov stable
if, for every \(\varepsilon>0\), some \(\delta>0\) makes
\(\dist(F^nx,p)<\varepsilon\) for all \(n\geq0\) whenever
\(\dist(x,p)<\delta\). It is asymptotically stable if it is stable and attracts
a neighborhood. It is locally exponentially stable if there are \(C>0\),
\(0<q<1\) and a neighborhood on which
\(\norm{F^nx-p}\leq Cq^n\norm{x-p}\).
\end{definition}

\begin{lemma}[Adapted finite-dimensional norms]\label{d13:15:adapted-norm}
Let \(A\) be a real or complex finite-dimensional linear operator. If
\(\rho(A)<a\), there is an equivalent norm \(\norm{\cdot}_*\) with
\(\norm{A}_*\leq a\). If \(A\) is invertible and all its eigenvalues have
modulus greater than \(a>0\), an equivalent norm satisfies
\(\norm{Av}_*\geq a\norm{v}_*\).
\end{lemma}
\begin{proof}
Choose \(b\) strictly between \(\rho(A)\) and \(a\). Jordan normal form
gives \(\norm{A^n}\leq Cb^n\) after enlarging \(C\): each block contributes
an exponential times a polynomial, absorbed by choosing an intermediate
rate if necessary. Thus
\[
 \norm{v}_*=\sum_{n=0}^\infty a^{-n}\norm{A^nv}
\]
is finite, is a norm, and lies between \(\norm{v}\) and a constant times
\(\norm{v}\). Shifting the sum shows
\(\norm{Av}_*=a(\norm{v}_*-\norm{v})\leq a\norm{v}_*\).
For the second assertion apply the first one to \(A^{-1}\), choosing
an upper bound strictly smaller than \(a^{-1}\), and invert its inequality.
\end{proof}

\begin{theorem}[Linearized attraction]\label{d13:15:linearized-stability}
Let \(F\) be \(C^1\) near a fixed point \(p\in\R^d\). If
\(\rho(DF(p))<1\), then \(p\) is locally exponentially stable.
\end{theorem}
\begin{proof}
Translate \(p\) to zero. By Lemma~\ref{d13:15:adapted-norm} choose an
equivalent norm and \(q<1\) with \(\norm{DF(0)}_*<q\). Continuity of the
derivative gives a convex ball \(B_*(0,r)\) on which
\(\norm{DF(x)}_*\leq q\). Integration along line segments gives
\(\norm{F(x)}_*\leq q\norm{x}_*\). The ball is positively invariant, and
induction gives \(\norm{F^nx}_*\leq q^n\norm{x}_*\). Norm equivalence
supplies the constant \(C\) for the original norm and proves stability.
\end{proof}

\begin{theorem}[Linearized instability]\label{d13:15:linearized-instability}
Let \(F\) be \(C^1\) near a fixed point \(p\in\R^d\). If \(DF(p)\) has
an eigenvalue of modulus greater than one, \(p\) is Lyapunov unstable.
Eigenvalues of modulus one in other directions do not invalidate this
sufficient criterion.
\end{theorem}
\begin{proof}
Translate \(p\) to zero. The real generalized spectral decomposition
splits the space as \(E_u\oplus E_c\), where \(E_u\neq0\) contains precisely
the eigenvalues of modulus greater than one and \(E_c\) the remaining
ones. This includes conjugate complex pairs as real invariant subspaces.
Choose \(a>1\) below the smallest unstable modulus and \(b<a\) above all
remaining moduli. Adapted norms give
\(\norm{A_uu}\geq a\norm{u}\) and \(\norm{A_cv}\leq b\norm{v}\).
Write
\[
 F(u,v)=(A_uu,A_cv)+R(u,v),\qquad
 \norm{R_u},\norm{R_c}\leq\eta(\norm{u}+\norm{v})
\]
on a sufficiently small product ball. Such an estimate follows from
\(R(0)=DR(0)=0\). Choose \(\eta>0\) with
\(a-2\eta>1\) and \(b+2\eta<a-2\eta\).
In the cone \(\norm{v}\leq\norm{u}\), as long as the orbit stays in this
ball, its next coordinates satisfy
\[
 \norm{u'}\geq(a-2\eta)\norm{u},\qquad
 \norm{v'}\leq(b+2\eta)\norm{u}\leq\norm{u'}.
\]
Start at an arbitrarily small nonzero \((u,0)\). If the orbit stayed
forever in the fixed ball, its unstable coordinate would grow
geometrically without bound. Consequently it leaves that ball, which
contradicts Lyapunov stability.
\end{proof}

\begin{proposition}[Periodic-orbit test]\label{d13:15:periodic-multipliers}
For a \(C^1\) map, let \(p\) have least period \(m\). If
\(\rho(D(F^m)(p))<1\), its finite periodic orbit is locally exponentially
attracting and stable. If this derivative has an eigenvalue of modulus
greater than one, the periodic orbit is unstable.
\end{proposition}
\begin{proof}
Apply the preceding two theorems to the fixed point of \(F^m\).
For attraction, write \(n=km+j\), \(0\leq j<m\); the finitely many maps
\(F^j\) are Lipschitz on sufficiently small neighborhoods of the orbit,
so their constants multiply the exponential estimate for \((F^m)^k\).
For instability take disjoint small neighborhoods of the distinct orbit
points. Continuity ensures that, while an orbit remains in their union
and starts close to \(p\), its itinerary follows these neighborhoods
cyclically. Stability of the orbit would therefore give stability of
\(p\) under \(F^m\), contrary to the instability theorem.
\end{proof}

\begin{example}[The unit-circle boundary needs nonlinear terms]
At zero the maps \(x\mapsto x-x^3\) and \(x\mapsto x+x^3\) both have
derivative one. For sufficiently small \(0<\abs{x}<1\), the first
strictly decreases \(\abs{x}\); its limiting value must be zero, while
the second strictly increases \(\abs{x}\) until it exits any fixed
small interval. Thus a multiplier on the unit circle alone settles
neither attraction nor instability.
\end{example}

\section{Circle maps and rotation}

\begin{definition}[Circle and lift]\label{d13:15:circle-lift}
Write \(\mathbb T=\R/\Z\), with distance
\(\dist([x],[y])=\inf_{k\in\Z}\abs{x-y-k}\).
An increasing homeomorphism \(L:\R\to\R\) with \(L(x+1)=L(x)+1\)
induces an orientation-preserving circle homeomorphism
\(F([x])=[L(x)]\). Such an \(L\) is called a lift.
\end{definition}

\begin{theorem}[Rotation number of an increasing lift]\label{d13:15:rotation-number}
For an increasing lift \(L\), the limit
\[
 \rho(L)=\lim_{n\to\infty}\frac{L^n(x)-x}{n}
\]
exists and is independent of \(x\). Replacing \(L\) by \(L+k\),
\(k\in\Z\), replaces \(\rho(L)\) by \(\rho(L)+k\).
The induced circle homeomorphism has a periodic point if and only if
\(\rho(L)\) is rational. More precisely, \(\rho(L)=p/q\), \(q\geq1\),
if and only if \(L^q(x)=x+p\) for some \(x\).
\end{theorem}
\begin{proof}
Put \(a_n=\sup_x(L^n(x)-x)\), \(b_n=\inf_x(L^n(x)-x)\); periodicity
reduces both extrema to a compact unit interval. For \(x\leq y\leq x+1\),
\(L^n(x)\leq L^n(y)\leq L^n(x)+1\), so \(a_n-b_n\leq1\).
Composition gives \(a_{n+m}\leq a_n+a_m\) and
\(b_{n+m}\geq b_n+b_m\). For fixed \(q\), write \(n=kq+r\),
\(0\leq r<q\), and use subadditivity to obtain
\(\limsup a_n/n\leq a_q/q\). Since every \(a_n/n\) is at least
\(\inf_q a_q/q\), the limit of \(a_n/n\) exists and equals that infimum.
Apply the same argument to \(-b_n\). Their difference is at most \(1/n\),
so the limits agree; squeezing proves the assertion for every \(x\).
Integer translations commute with \(L\), proving the lift-change rule.

If \(L^q(x)=x+p\), iterating this identity gives the rotation limit
\(p/q\). Conversely assume \(\rho(L)=p/q\) and put \(G=L^q-p\).
Then \(\rho(G)=q\rho(L)-p=0\), by its iterate formula. If \(G(x)-x\)
never vanished, continuity and periodicity would give either
\(G(x)-x\geq c>0\) everywhere or \(G(x)-x\leq-c<0\) everywhere.
Iteration would give \(\rho(G)\geq c\) or \(\rho(G)\leq-c\), a
contradiction. Thus \(G\) has a fixed point.
\end{proof}

\begin{theorem}[Rigid rotations]\label{d13:15:rigid-rotation}
For \(R_\alpha([x])=[x+\alpha]\), if \(\alpha=p/q\) in lowest terms,
every point has least period \(q\). If \(\alpha\) is irrational, every
forward orbit is dense, no point is periodic, and all orbit distances
are preserved.
\end{theorem}
\begin{proof}
The periodicity assertion is the equivalence
\(n\alpha\in\Z\), with minimal such \(n\). For irrational \(\alpha\),
pigeonhole the \(N+1\) points \([j\alpha]\), \(0\leq j\leq N\), into
\(N\) equal arcs. Their difference supplies a nonzero integer multiple
of \(\alpha\) at circle distance less than \(1/N\) from zero.
Hence the subgroup generated by \([\alpha]\) has positive representatives
arbitrarily close to zero, using an inverse when needed. The multiples
of such a representative form a mesh with maximum gap at most that
representative. The subgroup is therefore dense.

The nonnegative multiples have the same closure as the subgroup.
Indeed choose positive integers \(n_j\to\infty\) with
\([n_j\alpha]\to0\) (the preceding pigeonhole construction supplies
unbounded positive denominators). Then
\([(n_j-1)\alpha]\to[-\alpha]\). The closure of the nonnegative
semigroup is closed under addition and contains both generators
\([\alpha]\) and \([-\alpha]\), hence the whole subgroup. Translation
gives density of every orbit. The distance formula is translation
invariant, proving the final assertion.
\end{proof}

\section{A nonlinear interval family}

\begin{proposition}[Logistic fixed points]\label{d13:15:logistic-fixed-points}
Let \(f_\mu(x)=\mu x(1-x)\). The interval \([0,1]\) is positively
invariant for \(0\leq\mu\leq4\). Its fixed points are zero and, when
\(\mu\neq0\), \(1-\mu^{-1}\), the latter lying in \([0,1]\) for
\(\mu\geq1\). Zero is locally attracting for \(0\leq\mu<1\) and
repelling for \(\mu>1\); the nonzero fixed point is locally attracting
for \(1<\mu<3\) and repelling for \(\mu>3\).
\end{proposition}
\begin{proof}
On \([0,1]\), \(0\leq x(1-x)\leq1/4\). Solve
\(\mu x(1-x)=x\) without dividing by \(x\): either \(x=0\) or
\(\mu(1-x)=1\). Finally \(f_\mu'(0)=\mu\) and
\(f_\mu'(1-\mu^{-1})=2-\mu\), so the stability tests apply.
At the threshold values their strict hypotheses fail and no conclusion
is being asserted by this test.
\end{proof}

\begin{proposition}[An exact change of coordinates]\label{d13:15:logistic-tent}
The increasing homeomorphism \(h:[0,1]\to[0,1]\),
\(h(t)=\sin^2(\pi t/2)\), conjugates the tent map
\[
 \tau(t)=
 \begin{cases}2t,&0\leq t\leq1/2,\\2-2t,&1/2\leq t\leq1,\end{cases}
\]
to \(f_4\): \(h\circ\tau=f_4\circ h\).
\end{proposition}
\begin{proof}
The function \(h\) is continuous, strictly increasing, and takes the
endpoints to zero and one. The double-angle identity gives
\(4h(t)(1-h(t))=\sin^2(\pi t)\). For either branch of \(\tau\),
\(h(\tau(t))=\sin^2(\pi t)\), including the common endpoint.
\end{proof}

\paragraph{Sources and coverage.}
The map viewpoint, multiplier tests, interval iteration and conjugacy
are developed in Teschl, Chapters 10--11, and Hirsch--Smale--Devaney,
Chapter 15. The rotation-number proof supplies the circle-map bridge
explicitly. The logistic identities here are algebraic assertions,
not numerical assertions about an orbit diagram. Symbolic coding and
rigorous chaotic examples are proved in Chapter~\ref{ch:d13-17}.

% END chapters/M13.15


% BEGIN chapters/M13.16
\chapter{Local bifurcations and bounded center reductions}\label{ch:d13-16}

\section{What changes at a bifurcation}

\begin{definition}[A local family]\label{d13:16:family-definition}
Here a scalar flow family is \(x'=f(x,\mu)\), with real parameter \(\mu\);
a map family is \(x_{n+1}=F_\mu(x_n)\). We study equilibria or periodic
orbits in a fixed neighborhood of the distinguished state and parameter.
Their number, stability and conjugacy type may change at a bifurcation.
Statements below specify actual equations and nondegeneracy conditions;
a sketch of branches is not the definition of the conclusion.
\end{definition}

\begin{proposition}[Persistence of a simple equilibrium]\label{d13:16:simple-persistence}
If \(f\) is \(C^1\), \(f(0,0)=0\), and \(f_x(0,0)\neq0\), then all nearby
equilibria form a unique \(C^1\) graph \(x=x(\mu)\). Its scalar
linearized stability has constant type for sufficiently small \(\mu\).
For a map, the same persistence assertion holds if
\(\partial_xF_0(0)\neq1\), applied to \(F_\mu(x)-x\).
\end{proposition}
\begin{proof}
Apply the implicit function theorem to the displayed equation.
Continuity keeps the sign of \(f_x\), or the strict multiplier inequalities
when applicable, unchanged. For a scalar flow a negative derivative gives
attraction and a positive derivative gives repulsion by the sign of
\((f(x)-f(p))/(x-p)\) in a small interval.
\end{proof}

\begin{theorem}[Generic scalar saddle-node]\label{d13:16:saddle-node}
Suppose \(f\) is \(C^3\) near \((0,0)\) and
\[
 f(0,0)=f_x(0,0)=0,\qquad f_\mu(0,0)\neq0,\qquad f_{xx}(0,0)\neq0.
\]
There is a local state coordinate \(y=y(x)\), \(y'(0)=1\), in which the
equilibrium set is exactly \(\mu=a y^2\), where
\(a=-f_{xx}(0,0)/(2f_\mu(0,0))\neq0\).
Thus sufficiently close to the origin there are two equilibria on the
side \(\mu/a>0\), none on the other side, and a double equilibrium at
\(\mu=0\). The two simple equilibria have opposite scalar stability.
This conclusion describes the equilibrium set, not a claim of an exact
time-preserving normal form for the whole vector field.
\end{theorem}
\begin{proof}
The implicit function theorem in \(\mu\) solves \(f(x,\mu)=0\) uniquely
as \(\mu=\psi(x)\). Differentiating that identity gives
\(\psi(0)=\psi'(0)=0\) and
\(\psi''(0)=-f_{xx}(0,0)/f_\mu(0,0)=2a\).
Taylor's integral formula gives
\[
 \psi(x)=x^2h(x),\qquad
 h(x)=\int_0^1(1-t)\psi''(tx)\,\dd t,\qquad h(0)=a.
\]
After shrinking the interval, \(h/a>0\), and
\(y=x\sqrt{h(x)/a}\) is a \(C^1\) local coordinate with derivative one
at zero. This proves the exact description of the zero set. Finally,
along that set,
\(f_x(x,\psi(x))=-f_\mu(x,\psi(x))\psi'(x)\).
For small nonzero \(x\), \(\psi'(x)=2ax+o(x)\), so these derivatives
have opposite nonzero signs on the two branches.
\end{proof}

\begin{theorem}[Transcritical exchange]\label{d13:16:transcritical}
Let \(f\) be \(C^3\), with \(f(0,\mu)=0\) for all small \(\mu\), and
\[
 f_x(0,0)=0,\qquad a=f_{x\mu}(0,0)\neq0,\qquad
 b=\tfrac12f_{xx}(0,0)\neq0 .
\]
All nearby equilibria are \(x=0\) and a unique branch
\(x=\chi(\mu)=-a\mu/b+o(\mu)\). The branches meet only at \(\mu=0\)
and exchange stability: their derivatives have respective signs
\(\operatorname{sgn}(a\mu)\) and \(-\operatorname{sgn}(a\mu)\).
\end{theorem}
\begin{proof}
Factor without losing the root zero:
\[
 f(x,\mu)=xG(x,\mu),\qquad
 G(x,\mu)=\int_0^1f_x(tx,\mu)\,\dd t.
\]
Then \(G(0,0)=0\), \(G_\mu(0,0)=a\), \(G_x(0,0)=b\).
The implicit function theorem applied to \(G=0\) gives precisely
\(\chi\) and its derivative \(-a/b\). On the zero branch
\(f_x(0,\mu)=a\mu+o(\mu)\). On the other branch,
\(f_x(\chi(\mu),\mu)=\chi(\mu)G_x(\chi(\mu),\mu)
=-a\mu+o(\mu)\). Both are nonzero for sufficiently small nonzero
\(\mu\), proving the stated exchange.
\end{proof}

\begin{theorem}[A symmetric pitchfork]\label{d13:16:pitchfork}
Suppose \(f\) is \(C^5\), odd in \(x\), and
\[
 f_x(0,0)=0,\qquad a=f_{x\mu}(0,0)\neq0,\qquad
 b=\tfrac16f_{xxx}(0,0)\neq0 .
\]
Besides \(x=0\), two small equilibria exist exactly on the side
\(-a\mu/b>0\). They are \(x=\pm\sqrt{s(\mu)}\), where
\(s(0)=0\), \(s'(0)=-a/b\).
The derivative at zero has sign \(a\mu\), and that on either nonzero
branch has sign \(-a\mu\), for sufficiently small nonzero parameters.
\end{theorem}
\begin{proof}
As above write \(f=xG\); oddness makes \(G\) even in \(x\).
For \(s\geq0\) set \(H(s,\mu)=G(\sqrt{s},\mu)\).
This is \(C^1\) up to \(s=0\): for \(x>0\),
\[
 H_s(x^2,\mu)=\frac{G_x(x,\mu)}{2x}
 =\frac12\int_0^1G_{xx}(tx,\mu)\,\dd t ,
\]
which extends continuously to \(G_{xx}(0,\mu)/2\);
\(H_\mu\) extends in the same way. The \(C^5\) assumption also makes
this boundary value \(C^1\) in \(\mu\), so extending \(H\) linearly
in \(s\) for negative \(s\) gives a \(C^1\) function across the boundary.
At \((0,0)\), \(H=0\), \(H_\mu=a\), \(H_s=b\). The implicit function
theorem gives the unique small zero \(s=s(\mu)\). For small nonzero
\(\mu\), its sign is that of \(-a\mu/b\); nonnegative zeros correspond
exactly to the two square roots. At a nonzero root,
\[
 f_x=xG_x=2sH_s=-2a\mu+o(\mu).
\]
At zero, \(f_x=a\mu+o(\mu)\), proving the stability statement.
\end{proof}

\begin{example}[The symmetry is a genuine hypothesis]
The equation \(x'=\mu x-x^3\) has the stable branches
\(\pm\sqrt{\mu}\) for \(\mu>0\). Adding a constant perturbation
\(\varepsilon\neq0\) destroys the branch \(x=0\) and the odd symmetry.
Its equilibria solve a different cubic equation; the pitchfork theorem
cannot be applied while ignoring that perturbation.
\end{example}

\begin{proposition}[A period-doubling normal form]\label{d13:16:flip-normal-form}
For \(F_\mu(x)=-(1+\mu)x+x^3\), zero is locally attracting for small
\(\mu<0\) and unstable for small \(\mu>0\). For \(0<\mu<1\),
the two points \(\pm\sqrt{\mu}\) form a stable orbit of least period two,
with multiplier \((-1+2\mu)^2\). Its amplitude tends to zero as
\(\mu\downarrow0\).
\end{proposition}
\begin{proof}
The derivative at zero is \(-(1+\mu)\).
For \(a=\sqrt{\mu}>0\),
\(F_\mu(a)=-(1+\mu)a+a^3=-a\), and oddness gives \(F_\mu(-a)=a\);
the points are distinct. The derivative at either is \(-1+2\mu\).
The chain rule and Proposition~\ref{d13:15:periodic-multipliers}
give attraction since \(0<\mu<1\) implies \(\abs{-1+2\mu}<1\).
This proves a selected nondegenerate flip example, not a classification
of every family with multiplier \(-1\).
\end{proof}

\section{Generic planar Hopf bifurcation}

\begin{lemma}[A finite-order complex normal form]\label{d13:16:cubic-normal-form}
Let a real \(C^5\) planar family have an equilibrium at zero for each
small \(\mu\). Suppose the eigenvalues of its linearization are
\(a(\mu)\pm i b(\mu)\), with \(a(0)=0\) and \(b(0)=\omega>0\).
There are parameter-smooth real local coordinates, written as a complex
coordinate \(w\), in which
\[
 w'=(a(\mu)+ib(\mu))w+c(\mu)w^2\overline w+
       R(w,\overline w,\mu),\qquad R=O(\abs{w}^4),\quad
 D_wR=O(\abs{w}^3).
\]
The estimates are uniform for small \(\mu\), with the corresponding
parameter derivatives. The coefficient \(c(0)\) is obtained by the
explicit homological procedure in the proof. In these coordinates put
\(\ell=\Re c(0)\); its nonvanishing is the cubic nondegeneracy condition.
\end{lemma}
\begin{proof}
Let \(A(\mu)\) be the linearization. Choose a fixed real vector \(v\neq0\).
The vectors \(v\) and \((A-aI)v/b\) are independent near zero: dependence
would give a real eigenvector for a nonreal eigenvalue pair.
In this parameter-smooth basis \(A\) is multiplication by
\(\lambda=a+ib\). Taylor expansion in \(z,\overline z\) gives homogeneous
quadratic and cubic polynomials, with a \(C^1\)-controlled fourth-order
remainder.

For a monomial \(h_{jk}w^j\overline w^k\), the operator
\(H\mapsto DH(\lambda w)-\lambda H(w)\) multiplies its coefficient by
\[
 D_{jk}=j\lambda+k\overline\lambda-\lambda
       =(j+k-1)a+i(j-k-1)b.
\]
For \(j+k=2\) every such divisor is nonzero near \(\mu=0\).
In \(z=w+H_2(w,\overline w)\), choose each coefficient of \(H_2\)
to be the corresponding quadratic coefficient divided by \(D_{jk}\).
Indeed substitution gives the new quadratic term
\(f_2+\lambda H_2-DH_2(\lambda w)=f_2-\mathcal L H_2=0\).
The inverse function theorem applies because this coordinate map has
derivative the identity at zero. Expanding its inverse Jacobian through
degree three leaves a uniquely determined cubic polynomial \(g_3\);
this is an actual finite polynomial operation, with all fourth-order
terms retained in the remainder.

For \(j+k=3\), the divisors are nonzero near zero except possibly for
\((j,k)=(2,1)\). A second change \(w=u+H_3(u,\overline u)\), with
coefficients \((g_3)_{jk}/D_{jk}\) for the three nonresonant monomials
and zero for \((2,1)\), removes exactly those three terms.
Since the quadratic term is already zero, it does not introduce
additional cubic corrections. The remaining coefficient is
\(c(\mu)=(g_3)_{21}\). Thus the procedure specifies the coefficient
without requiring a memorized normalization-dependent formula.
Taylor's theorem and the two smooth near-identity inverses give the
stated remainder and derivative estimates. Replacing \(u\) by \(w\)
finishes the construction.
\end{proof}

\begin{lemma}[Return-map expansion]\label{d13:16:hopf-return-expansion}
For the normal form above, the first return to the positive real ray
has radial map, for sufficiently small \(r\geq0\),
\[
 P(r,\mu)=e^{2\pi a/b}r+K(\mu)r^3+O(r^4),\qquad
 K(0)=\frac{2\pi\ell}{\omega}.
\]
The remainder is \(C^1\)-controlled: its radial derivative is \(O(r^3)\)
and its parameter derivative is \(O(r^4)\). The return is defined
locally; it is not a global section for all trajectories.
\end{lemma}
\begin{proof}
Write \(w=re^{i\theta}\). The equations, for \(r>0\), are
\[
 r'=ar+(\Re c)r^3+O(r^4),\qquad
 \theta'=b+(\Im c)r^2+O(r^3).
\]
For small \(r,\mu\) the second expression exceeds \(\omega/2\).
Division gives an equation on the fixed angular interval \([0,2\pi]\),
\[
 \frac{\dd r}{\dd\theta}
 =Ar+Br^3+O(r^4),\qquad
 A=a/b,\quad B=\frac{\Re c}{b}-\frac{a\Im c}{b^2}.
\]
The fourth-order terms and their first derivatives are uniform on
this compact angular interval. Existence and Gronwall estimates keep
a sufficiently small initial radius in the coordinate neighborhood
throughout one turn. Variation of constants first gives
\(r(\theta)=e^{A\theta}r_0+O(r_0^3)\). Substituting this back into the
cubic integral gives
\[
 r(2\pi)=e^{2\pi A}r_0+
 B e^{2\pi A}\left(\int_0^{2\pi}e^{2As}\,\dd s\right)r_0^3
 +O(r_0^4).
\]
Differentiating the angular integral equation in \(r_0\) and in \(\mu\),
and using the same uniform bounds, gives the stated derivative
remainders. Hence \(K=B e^{2\pi A}\int_0^{2\pi}e^{2As}\,\dd s\);
at zero it is \(2\pi\ell/\omega\).
One can extend the signed-radius angular equation across \(r=0\)
by Taylor expansion, so division by \(r_0\) below has a \(C^2\) extension.
\end{proof}

\begin{theorem}[Nondegenerate planar Hopf]\label{d13:16:planar-hopf}
Under Lemma~\ref{d13:16:cubic-normal-form}, suppose in addition
\(a'(0)\neq0\) and \(\ell\neq0\).
For sufficiently small parameters, there is a unique small periodic
orbit precisely on the side \(a(\mu)\ell<0\). Its positive-ray radius
satisfies
\[
 \mu=-\frac{\ell}{a'(0)}r^2+O(r^3),\qquad
 r^2=-\frac{a'(0)}{\ell}\mu+o(\abs{\mu}).
\]
Its period tends to \(2\pi/\omega\).
It is asymptotically orbitally stable if \(\ell<0\) and unstable if
\(\ell>0\); the equilibrium has the stability type determined by
the sign of \(a(\mu)\). Uniqueness concerns periodic orbits contained
in a sufficiently small fixed neighborhood.
\end{theorem}
\begin{proof}
Set \(Q(r,\mu)=(P(r,\mu)-r)/r\) for \(r>0\) and use its smooth extension
at zero. The preceding expansion gives
\[
 Q=e^{2\pi a/b}-1+K(\mu)r^2+O(r^3),\qquad
 Q_\mu(0,0)=2\pi a'(0)/\omega\neq0.
\]
The implicit function theorem gives all small zeros as
\(\mu=\psi(r)\), with \(\psi(0)=\psi'(0)=0\) and
\(\psi(r)=-\ell r^2/a'(0)+O(r^3)\).
The derivative estimate is
\(\psi'(r)=-2\ell r/a'(0)+O(r^2)\), so this branch is strictly monotone
for small \(r>0\) and covers just the asserted side of parameter space.
No extra small periodic orbit is missed: angular velocity is positive,
so every such orbit crosses the positive ray and gives a fixed point of
its first-return map. A fixed point of an increasing local scalar return
map also rules out a distinct higher-cycle interpretation.

At a nonzero fixed radius, \(P/r=1\), so
\[
 P_r-1=rQ_r=2K(0)r^2+O(r^3).
\]
For \(\ell<0\) this multiplier lies between zero and one; for
\(\ell>0\) it exceeds one. The local section and the smooth flow carry
the return-map stability to orbital stability, as in
Chapter~\ref{ch:d13-14}. Integrating
\(\dd t/\dd\theta=1/(b+O(r^2))\) over a turn gives the limiting period.
Finally the equilibrium has real eigenvalue part \(a(\mu)\), so the
strict linearized flow stability or instability criterion applies.
\end{proof}

\begin{example}[A check on the normalization]
For \(z'=(\mu+i\omega)z-z\abs{z}^2\), the quadratic elimination is
empty and \(c=-1\). The exact polar equations are
\(r'=\mu r-r^3\), \(\theta'=\omega\). The stable cycle has
\(r=\sqrt{\mu}\) for \(\mu>0\), verifying the signs, amplitude and
period in the theorem. If the cubic term is absent, \(\ell=0\);
the nondegenerate theorem deliberately makes no conclusion.
\end{example}

\section{Invariant graphs: a bounded center-manifold bridge}

\begin{proposition}[The invariance equation]\label{d13:16:invariant-graph}
For a \(C^1\) system \(x'=f(x,y)\), \(y'=g(x,y)\), a \(C^1\) graph
\(y=h(x)\) is locally invariant if
\[
 Dh(x)f(x,h(x))=g(x,h(x)).
\]
More precisely this identity implies that trajectories of
\(x'=f(x,h(x))\), lifted to \((x,h(x))\), are trajectories of the
full system while they remain in the graph domain; uniqueness gives
the local invariance. Conversely local invariance implies the identity.
\end{proposition}
\begin{proof}
Differentiate \(h(x(t))\) and use the identity. This proves both component
equations for the lift. Uniqueness of the full initial-value problem
identifies it with the full trajectory. Conversely differentiate a
full trajectory starting at a graph point at time zero.
\end{proof}

\begin{theorem}[A proved triangular center graph]\label{d13:16:triangular-center}
Consider
\[
 x'=f(x),\qquad y'=-By+g(x),\qquad (x,y)\in\R^m\times\R^k.
\]
Assume \(f\) is \(C^1\), globally Lipschitz with constant \(L\), and
\(f(0)=Df(0)=0\). Assume \(g\) is \(C^1\), bounded with bounded
derivative, and \(g(0)=Dg(0)=0\). If
\(\norm{e^{-Bt}}\leq M e^{-\beta t}\) for \(t\geq0\), where
\(\beta>L\), then
\[
 h(x)=\int_0^\infty e^{-Bs}g(\phi_{-s}(x))\,\dd s
\]
is \(C^1\), \(h(0)=Dh(0)=0\), and its graph is invariant.
For any full solution,
\[
 y(t)-h(x(t))=e^{-Bt}(y(0)-h(x(0)))\quad(t\geq0).
\]
Thus the graph is tangent to the center subspace and attracts each
vertical fiber exponentially. This is a triangular special theorem;
it does not assert a general coupled center-manifold theorem.
\end{theorem}
\begin{proof}
Global Lipschitz continuity gives a complete slow flow and
\(\norm{D\phi_{-s}(x)}\leq e^{Ls}\), by its variational equation.
The integral for \(h\) converges absolutely, uniformly in \(x\).
Its derivative integrand has norm at most
\(M\norm{Dg}_\infty e^{-(\beta-L)s}\), which is integrable.
Differentiation under the integral, and dominated convergence for the
continuous derivative, prove \(C^1\) regularity. At zero the integrands
for \(h\) and \(Dh\) vanish.

For every \(t\geq0\), change variables in the integral to obtain
\[
 h(\phi_t x)
 =e^{-Bt}h(x)+\int_0^t e^{-B(t-u)}g(\phi_u x)\,\dd u.
\]
This is exactly the variation-of-constants expression for the fast
equation with initial value \(h(x)\); invariance follows.
Subtracting it from the same expression for an arbitrary fast initial
value proves the tracking identity and the exponential estimate.
\end{proof}

\begin{remark}[Scope of the reduction]
The invariance equation alone neither constructs a graph nor proves its
regularity. The triangular theorem supplies one genuine construction
under stated gap assumptions. Hyperbolic stable and unstable manifolds
are treated separately in Chapter~\ref{ch:d13-20}. General high-dimensional
coupled center-manifold existence, higher-codimension bifurcation theory
and global bifurcations are further developments, not silently assumed
lemmas in the preceding planar proofs.
\end{remark}

\paragraph{Sources and coverage.}
The scalar examples and elementary flip mechanism are compared with
Teschl, Chapters 7 and 11, and Hirsch--Smale--Devaney, Chapters 8, 12
and 15. Guckenheimer--Holmes, Chapter 3, is the reference branch for
local bifurcations, center reductions and normal forms. The finite
polynomial elimination and return-map estimates have been written
out here so that the generic planar Hopf statement has its own
proof bridge rather than being inferred from its polar example.

% END chapters/M13.16


% BEGIN chapters/M13.17
\chapter{Symbolic dynamics and rigorous chaotic examples}\label{ch:d13-17}

\section{A precise convention for chaos}

\begin{definition}[Transitivity, density and sensitivity]\label{d13:17:chaos-definition}
For a continuous map \(F:X\to X\) on a metric space, topological
transitivity means that, for nonempty open \(U,V\subset X\), some
\(n\geq0\) satisfies \(F^n(U)\cap V\neq\varnothing\).
Sensitive dependence means that some \(\delta>0\) has the following
property: for every \(x\) and every neighborhood \(U\) of \(x\), there
are \(y\in U\) and \(n\geq0\) with
\(\dist(F^nx,F^ny)>\delta\).
In this chapter \emph{chaotic} means transitive, with dense periodic
points and sensitive dependence. Each of the three properties will
be checked for the examples; this convention is not a claim that
all definitions of chaos are equivalent on all spaces.
\end{definition}

\begin{definition}[One-sided and two-sided shifts]\label{d13:17:shift-definition}
Let \(\Sigma^+=\{0,1\}^{\N_0}\) and \(\Sigma=\{0,1\}^{\Z}\).
The shift is \((\sigma s)_j=s_{j+1}\).
On \(\Sigma^+\), set \(d_+(s,t)=2^{-N}\), where \(N\) is the first
index at which \(s,t\) differ, and set \(d_+(s,s)=0\).
On \(\Sigma\), use
\[
 d_{\Z}(s,t)=\sum_{j\in\Z}2^{-\abs{j}-2}\abs{s_j-t_j}.
\]
These metrics induce the product topology. Specifying finitely many
coordinates gives a cylinder set; initial cylinders suffice on
\(\Sigma^+\). The one-sided shift is surjective but not injective;
the two-sided shift is a homeomorphism.
\end{definition}

\begin{lemma}[Compactness of the symbolic spaces]\label{d13:17:shift-compactness}
Both symbolic spaces are compact, have no isolated points, and are
totally disconnected.
\end{lemma}
\begin{proof}
Given a sequence of symbolic points, successively choose subsequences
constant in the first, second, and subsequent coordinates in an
enumeration of the countable index set. Its diagonal subsequence
converges in every coordinate. Uniformly small tails of either metric
then give metric convergence. Sequential compactness is compactness
for metric spaces. A finite cylinder leaves another coordinate free,
so it contains more than one point. Distinct points differ in one
coordinate and are separated by the two complementary clopen
cylinders specifying that coordinate; a connected set can therefore
contain at most one point.
\end{proof}

\begin{theorem}[Full-shift chaos]\label{d13:17:full-shift-chaos}
Both shifts are chaotic in Definition~\ref{d13:17:chaos-definition}.
For every \(n\geq1\), \(\sigma^n\) has exactly \(2^n\) fixed points.
Each shift also has a dense forward orbit.
\end{theorem}
\begin{proof}
For transitivity, a cylinder in \(U\) prescribes a finite block and a
cylinder in \(V\) prescribes another. Place the second block sufficiently
far to the right of the first and fill all other coordinates arbitrarily.
This point lies in \(U\) and shifts into \(V\).
For density of periodic points, extend any prescribed finite block
periodically; on the two-sided space use a block containing all the
specified coordinates. The \(n\)-fixed points are exactly the
periodic repetitions of binary blocks of length \(n\), giving \(2^n\)
distinct points. This count concerns fixed points of \(\sigma^n\),
not points of least period \(n\).

For sensitivity on the one-sided space, change a single coordinate
beyond the finite prefix controlled by the neighborhood. After shifting
to that coordinate the distance is one. On the two-sided space make
the change at a sufficiently large positive coordinate outside the
specified finite window; at the corresponding time the distance is
at least \(1/4\). Thus \(\delta=1/2\) and \(\delta=1/8\), respectively,
work. Finally concatenate an enumeration of all finite binary words
along the nonnegative coordinates. Its shifts meet every finite
cylinder, including any finite two-sided window after centering the
corresponding occurrence. This gives a dense forward orbit.
\end{proof}

\section{Contracting inverse branches and polynomial dynamics}

\begin{theorem}[A coding theorem with disjoint branches]\label{d13:17:inverse-branch-coding}
Let \(K\) be a nonempty compact metric space of finite diameter, and
let \(b_0,b_1:K\to K\) be injective contractions with common constant
\(q<1\) and disjoint images. Suppose a continuous map \(F\), defined
on a neighborhood containing those images, satisfies \(Fb_i=\id_K\).
Set
\[
 K_n=\bigcup_{s_0,\ldots,s_{n-1}\in\{0,1\}}
          b_{s_0}\cdots b_{s_{n-1}}K,\qquad
 \Lambda=\bigcap_{n\geq1}K_n .
\]
There is a homeomorphism \(\pi:\Sigma^+\to\Lambda\) satisfying
\(F\pi=\pi\sigma\). In particular \(\Lambda\) is compact, invariant,
has no isolated points, is totally disconnected, and \(F|_\Lambda\)
is chaotic.
\end{theorem}
\begin{proof}
For \(s\in\Sigma^+\), its successive cylinder images are nested,
nonempty compact sets of diameter at most \(q^n\operatorname{diam}K\).
Compactness and the diameter bound give a unique point \(\pi(s)\).
Agreement of \(n\) initial symbols gives the same diameter bound on
the two image points, so \(\pi\) is continuous.
Disjointness of the first images and injectivity of each branch show
that two different first-divergence words cannot code the same point.
Every point of \(\Lambda\) has a compatible branch at every depth:
its first branch is unique, and applying \(F\) successively gives the
next unique branch. Hence it equals \(\pi(s)\) for its itinerary.
The map is therefore a continuous bijection from a compact space to a
Hausdorff space and is a homeomorphism.
Applying \(F\) removes the first branch, proving the intertwining
identity. Surjectivity of \(\sigma\) proves \(F(\Lambda)=\Lambda\).
The topological properties and transitivity/periodic density transport
by the homeomorphism. For sensitivity, uniform continuity of
\(\pi^{-1}\) on the compact image supplies \(\varepsilon>0\) such that
image distance less than \(\varepsilon\) forces symbolic distance
less than the shift sensitivity constant. The contrapositive gives
image distance at least \(\varepsilon\); choosing the final
sensitivity constant \(\varepsilon/2\) proves the strict inequality
required in Definition~\ref{d13:17:chaos-definition}.
\end{proof}

\begin{theorem}[A chaotic invariant set for a quadratic polynomial]\label{d13:17:logistic-cantor}
For every \(\mu\geq5\), \(f_\mu(x)=\mu x(1-x)\) has a compact invariant
Cantor set \(\Lambda_\mu\subset[0,1]\) on which it is conjugate to
the one-sided binary shift and hence chaotic.
This does not say that \(f_\mu\) maps all of \([0,1]\) into itself.
\end{theorem}
\begin{proof}
For \(u\in[0,1]\) the two inverse branches are
\[
 b_0(u)=\tfrac12(1-\sqrt{1-4u/\mu}),\qquad
 b_1(u)=\tfrac12(1+\sqrt{1-4u/\mu}).
\]
Their images lie on opposite sides of \(1/2\), separated by a
positive gap, and \(f_\mu b_i(u)=u\). Their derivative moduli are
at most
\[
 q=\frac1{\mu\sqrt{1-4/\mu}}
   =\frac1{\sqrt{\mu(\mu-4)}}\leq\frac1{\sqrt5}<1.
\]
The coding theorem applies with \(K=[0,1]\).
The full interval fails to be invariant since \(f_\mu(1/2)=\mu/4>1\).
\end{proof}

\begin{proposition}[Tent-map exactness and chaos]\label{d13:17:tent-chaos}
The tent map \(\tau\) of Proposition~\ref{d13:15:logistic-tent} is
chaotic. Every nonempty relatively open interval \(I\subset[0,1]\)
has \(\tau^n(I)=[0,1]\) for some \(n\). Consequently \(f_4\) is
chaotic on its full invariant interval.
\end{proposition}
\begin{proof}
Induction shows that \(\tau^n\) has \(2^n\) monotone affine branches,
each on a closed interval of length \(2^{-n}\), and each branch
maps onto \([0,1]\). Every nonempty open interval contains such a
closed branch interval for sufficiently large \(n\), proving exactness
and therefore transitivity.
On a branch interval \(J=[a,b]\subset I\), the endpoint values of
\(\tau^n\) are zero and one, in one order or the other.
The continuous function \(\tau^n(x)-x\) has opposite weak signs at
the endpoints, so it has a zero in \(J\). Periodic points are dense.
For sensitivity, fix \(x\) and its neighborhood \(I\). Choose a
full branch interval \(J\subset I\). One of the two endpoint images
zero or one is at distance at least \(1/2\) from \(\tau^n(x)\);
the corresponding point \(y\in J\) witnesses sensitivity with any
fixed constant less than \(1/2\). Finally transport these properties
through the homeomorphism \(h\) of
Proposition~\ref{d13:15:logistic-tent}; compactness supplies the
sensitivity transport as in the coding theorem.
\end{proof}

\begin{remark}[Why not identify the interval with the binary shift?]
At the tent-map turning point, the two inverse branches meet.
Itineraries at preimages of that point need boundary conventions,
and the coding need not be injective. The full interval is connected,
whereas the full symbolic space is totally disconnected, so they
cannot be homeomorphic. The disjoint-branch hypothesis above is
essential to its stronger conjugacy conclusion.
\end{remark}

\section{An explicit two-sided horseshoe model}

\begin{theorem}[Affine strip horseshoe]\label{d13:17:affine-horseshoe}
Fix \(0<r<1/2\), and let
\[
 C_r=\left\{(1-r)\sum_{j=0}^\infty s_jr^j:
                       s_j\in\{0,1\}\right\}\subset[0,1].
\]
On the disjoint vertical strips
\(V_0=[0,r]\times[0,1]\),
\(V_1=[1-r,1]\times[0,1]\), define
\[
 H(x,y)=\left(\frac{x-(1-r)i}{r},\,ry+(1-r)i\right)
             \quad\hbox{for }(x,y)\in V_i .
\]
This map is a homeomorphism from \(V_0\cup V_1\) onto two disjoint
horizontal strips. Its maximal set with a bi-infinite orbit in the
unit square is \(\Lambda=C_r\times C_r\), and \(H|_\Lambda\) is
conjugate to the two-sided binary shift. On each strip its derivative
has expanding and contracting factors \(r^{-1}\) and \(r\).
The statement concerns this explicitly defined local strip model;
no unspecified smooth extension over the middle gap is assumed.
\end{theorem}
\begin{proof}
The two strip images are
\([0,1]\times[0,r]\) and
\([0,1]\times[1-r,1]\), with a positive gap.
The displayed affine inverses are continuous, proving
the homeomorphism assertion on the disconnected domain.
For a bi-infinite sequence \(s\), put
\[
 x(s)=(1-r)\sum_{j=0}^\infty s_jr^j,\qquad
 y(s)=(1-r)\sum_{j=1}^\infty s_{-j}r^{j-1}.
\]
These series converge uniformly. Because \(r<1/2\), a leading-digit
difference outweighs the sum of all following differences:
\((1-r)-r=1-2r>0\), after multiplication by the corresponding
power of \(r\). Each coordinate thus has a unique binary expansion
of this form. The map \(s\mapsto(x(s),y(s))\) is a continuous bijection
from \(\Sigma\) to \(C_r\times C_r\), hence a homeomorphism.
Removing the leading digit of \(x\) and adding it at the front of
\(y\) is exactly the formula for \(H\); therefore this coding
intertwines \(H\) with \(\sigma\).

Conversely, a point whose forward iterates stay in the square must
have a strip choice \(s_j\) at each nonnegative time. Solving the
expanding-coordinate recurrence for \(x\) leaves a remainder bounded
by \(r^n\), so taking \(n\to\infty\) forces the displayed future
series. Its backward strip choices similarly force the past series
for \(y\). Thus no other full orbit in the square is omitted.
Finally the derivative is \(\operatorname{diag}(r^{-1},r)\) on each
branch, proving the expansion and contraction statements.
\end{proof}

\begin{example}[Ergodic rotation is not sensitive]
An irrational circle rotation has dense orbits by
Theorem~\ref{d13:15:rigid-rotation}, but it is an isometry and has no
periodic points. Given a proposed sensitivity constant, take a
smaller-radius neighborhood; all paired distances remain smaller
for every time. Thus density of orbits alone does not prove chaos.
Its measure-theoretic ergodicity will be a separate assertion in
Chapter~\ref{ch:d13-19}.
\end{example}

\paragraph{Sources and coverage.}
The itinerary construction follows the mathematical themes of
Teschl, Chapters 11 and 13, especially its explicit strip construction,
and Hirsch--Smale--Devaney, Chapters 15--16. Guckenheimer--Holmes,
Chapter 5, supplies the broader hyperbolic and symbolic context.
Only the explicit coding and horseshoe theorems proved here are
adopted; a transverse-homoclinic theorem or a numerically observed
Lorenz attractor is not being smuggled in as a proved consequence.

% END chapters/M13.17


% BEGIN chapters/M13.18
\chapter{Growth exponents, attractors and Cantor dimensions}\label{ch:d13-18}

\section{Finite-time growth and actual asymptotic limits}

\begin{definition}[Derivative growth]\label{d13:18:exponent-definition}
For a \(C^1\) map \(F:U\to U\), \(U\subset\R^d\) open, \(x\in U\)
and \(v\neq0\), the finite-time directional growth rate is
\[
 \lambda_n(x,v)=\frac1n\log\frac{\norm{D(F^n)(x)v}}{\norm{v}} .
\]
The directional Lyapunov exponent is its limit when that limit exists,
possibly \(-\infty\); \(\log0=-\infty\).
An operator-norm upper rate uses
\(\limsup_n n^{-1}\log\norm{D(F^n)(x)}\).
For a differentiable flow replace \(n\) by \(t\to+\infty\).
For a matrix cocycle \(A(x)\) over a map \(T\), replace \(D(F^n)(x)\)
by \(A(T^{n-1}x)\cdots A(x)\).
A limsup is not, by definition alone, an existing limit.
\end{definition}

\begin{proposition}[Norm independence]\label{d13:18:norm-independence}
In a fixed finite-dimensional space the existence and value of each
of these limits, and of the corresponding limsup, are independent
of the chosen equivalent norm.
\end{proposition}
\begin{proof}
If \(c\norm{w}_1\leq\norm{w}_2\leq C\norm{w}_1\), the logarithms of
the two growth ratios differ in absolute value by at most
\(\log(C/c)\), uniformly in the nonzero image vector.
Division by \(n\), or by \(t\), makes that difference tend to zero.
A zero image has rate \(-\infty\) in both norms and remains zero
under every subsequent linear factor. Operator norms differ by
fixed multiplicative constants as well.
\end{proof}

\begin{theorem}[Constant linear-flow exponents]\label{d13:18:linear-exponents}
Let \(A\) be a real or complex finite-dimensional matrix and \(v\neq0\).
In the complex generalized eigenspace decomposition let
\[
 \alpha(v)=\max\{\Re\lambda:\text{the component of \(v\) in \(E_\lambda\)
                                        is nonzero}\}.
\]
Then
\(\lim_{t\to\infty}t^{-1}\log\norm{e^{tA}v}=\alpha(v)\).
In particular the operator-norm limit is the largest real part of
an eigenvalue. Nontrivial Jordan blocks do not change these limits.
\end{theorem}
\begin{proof}
On \(E_\lambda\), \(A=\lambda I+N\) with \(N\) nilpotent, so
\(e^{tA}=e^{t\lambda}\sum_{j=0}^{m-1}t^jN^j/j!\).
Summing the finitely many block estimates gives
\(\norm{e^{tA}v}\leq C e^{\alpha(v)t}(1+t)^m\).
Choose a nonzero component \(v_\lambda\) with
\(\Re\lambda=\alpha(v)\). Since \(e^{-tN}\) has polynomial norm,
\[
 \norm{e^{tN}v_\lambda}
 \geq\frac{\norm{v_\lambda}}{\norm{e^{-tN}}}
 \geq c(1+t)^{-m}.
\]
The bounded spectral projection onto this eigenspace then yields
\(\norm{e^{tA}v}\geq c'e^{\alpha(v)t}(1+t)^{-m}\).
Taking logarithms and dividing by \(t\) squeezes the limit.
The same upper bound for all vectors and a vector in a maximizing
eigenspace prove the operator-norm assertion. Real spaces may be
complexified, using equivalent finite-dimensional norms.
\end{proof}

\begin{proposition}[One-dimensional and volume exponents]\label{d13:18:ergodic-exponents}
Let \(F\) be \(C^1\) on an open subset of \(\R^d\), and let
\(\mu\) be an invariant probability measure on a compact invariant
set \(K\) in that open set. If
\(g(x)=\log\abs{\det DF(x)}\) is in \(L^1(\mu)\), then
\[
 \frac1n\log\abs{\det D(F^n)(x)}
       \longrightarrow g^*(x)
 \quad\text{for almost every }x,
\]
where \(g^*\) is invariant and \(\int g^*\,\dd\mu=\int g\,\dd\mu\).
If \(\mu\) is ergodic, the limit is the constant \(\int g\,\dd\mu\).
For \(d=1\) this is the directional derivative exponent.
For \(d>1\) it is a volume-growth assertion, not an assertion that
every individual directional exponent exists.
\end{proposition}
\begin{proof}
Integrability means the determinant is nonzero almost everywhere.
Remove the union of the inverse images of its null zero set; measure
preservation keeps this union null. The chain rule and multiplicativity
of the determinant give, at every remaining point,
\[
 \log\abs{\det D(F^n)(x)}
       =\sum_{j=0}^{n-1}g(F^jx).
\]
The pointwise and ergodic conclusions now follow from
Theorem~\ref{d13:19:birkhoff} below. In dimension one a nonzero
direction contributes only a canceling fixed scalar factor.
\end{proof}

\begin{example}[A bounded cocycle without a pointwise exponent]
Over the one-sided shift take the scalar cocycle \(A(s)=e^{s_0}\).
It is continuous, bounded and invertible. Choose a sequence consisting
of alternating all-zero and all-one blocks, with lengths
\(L_j=2^{2^j}\). At the end of block \(j\), the proportion occupied
by earlier blocks is at most
\((L_1+\cdots+L_{j-1})/L_j\to0\).
Thus the logarithmic growth average has subsequences tending to zero
and to one. It has no limit at that base point.
This does not contradict the preceding almost-everywhere theorem:
it neither supplies an invariant ergodic measure typical for that
point nor verifies a full-measure claim about it.
\end{example}

\section{Attractors from trapping sets}

\begin{definition}[A stable attractor]\label{d13:18:attractor-definition}
For a continuous map on a metric space, a stable attractor is a
nonempty compact set \(A\) with \(F(A)=A\), which is Lyapunov stable
and attracts a neighborhood: \(\dist(F^nx,A)\to0\) for every point
in that neighborhood. Stability uses neighborhoods of the whole set.
The basin is the set of points whose forward iterates approach \(A\).
This definition does not require transitivity or positive exponents.
\end{definition}

\begin{theorem}[A compact trapping-set construction]\label{d13:18:trapping-attractor}
Let \(F:X\to X\) be continuous on a metric space. If \(K\neq\varnothing\)
is compact and \(F(K)\subset\operatorname{int}K\), then
\[
 A=\bigcap_{n\geq0}F^n(K)
\]
is a stable attractor. Attraction is uniform for starting points in
\(K\): every neighborhood \(V\) of \(A\) contains \(F^n(K)\) for
all sufficiently large \(n\). Its basin contains the open set
\(\bigcup_{n\geq0}F^{-n}(\operatorname{int}K)\).
\end{theorem}
\begin{proof}
The sets \(F^n(K)\) are nonempty compact and decreasing, so their
intersection is nonempty and compact. Continuity gives \(F(A)\subset A\).
If \(y\in A\), choose \(x_n\in F^n(K)\) with \(F(x_n)=y\).
A convergent subsequence in \(K\) has limit \(x\); for every fixed
\(m\), eventually \(x_n\in F^m(K)\), hence \(x\in A\) and \(F(x)=y\).
Thus \(F(A)=A\).
If the uniform neighborhood assertion failed, there would be
\(x_j\in F^{n_j}(K)\setminus V\), \(n_j\to\infty\). A convergent
subsequence would have a limit in every \(F^m(K)\), but outside the
open set \(V\), contradicting \(A\subset V\).

Given \(\varepsilon>0\), choose \(m\) with
\(F^m(K)\subset N_\varepsilon(A)\). Since
\(A\subset F(K)\subset\operatorname{int}K\), the open neighborhood
\[
 W=\operatorname{int}K\cap
       \bigcap_{j=0}^{m-1}F^{-j}(N_\varepsilon(A))
\]
contains \(A\). Its first \(m\) iterates lie in the required
neighborhood by construction; later iterates lie in
\(F^m(K)\). This proves stability. The uniform assertion gives
attraction from \(K\), and any point entering \(\operatorname{int}K\)
subsequently has the same property. The displayed union is open.
\end{proof}

\begin{example}[Attraction need not mean chaos]
For \(F(x)=x/2\) on \(\R\), \(K=[-1,1]\) is trapping and
\(A=\{0\}\). Its derivative exponent is \(-\log2\).
This is a stable attractor, but it has neither the sensitive
dependence nor the large orbit structure of the horseshoe.
\end{example}

\section{Hausdorff dimension in a closed example}

\begin{definition}[Hausdorff measure and dimension]\label{d13:18:hausdorff-definition}
For \(s>0\) and \(\delta>0\), define
\[
 \mathcal H^s_\delta(E)=
 \inf\left\{\sum_i(\operatorname{diam}U_i)^s:
 E\subset\bigcup_iU_i,\ \operatorname{diam}U_i\leq\delta\right\}.
\]
Set \(\mathcal H^s(E)=\lim_{\delta\downarrow0}\mathcal H^s_\delta(E)\)
and
\(\dim_H E=\inf\{s>0:\mathcal H^s(E)=0\}\).
For a compact subset of the line, let \(N_\varepsilon(E)\) be the
least number of intervals of length at most \(\varepsilon\) that
cover it. Its box dimension is
\(\lim_{\varepsilon\downarrow0}\log N_\varepsilon(E)/
 \log(1/\varepsilon)\), when that limit exists.
\end{definition}

\begin{lemma}[The mass-distribution estimate]\label{d13:18:mass-distribution}
If a Borel probability measure \(\nu\), supported on \(E\subset\R\),
satisfies \(\nu(J)\leq C\abs{J}^s\) for every sufficiently short
interval \(J\), then \(\mathcal H^s(E)\geq C^{-1}\).
\end{lemma}
\begin{proof}
Every bounded subset of the line is contained in a closed interval
of length equal to its diameter. Replace a sufficiently fine
countable cover by these intervals. Countable subadditivity gives
\(1\leq\sum_i\nu(J_i)\leq C\sum_i(\operatorname{diam}U_i)^s\).
Take the infimum and then let the mesh tend to zero.
\end{proof}

\begin{theorem}[Dimension of a two-branch Cantor set]\label{d13:18:cantor-dimension}
For \(0<r<1/2\), the set \(C_r\) of
Theorem~\ref{d13:17:affine-horseshoe} has
\[
 \dim_H C_r=\dim_{\mathrm{box}}C_r
       =d=\frac{\log2}{\log(1/r)},\qquad
          0<\mathcal H^d(C_r)<\infty .
\]
In particular the middle-thirds set has dimension \(\log2/\log3\).
\end{theorem}
\begin{proof}
At level \(n\), \(C_r\) lies in \(2^n\) disjoint closed intervals
of length \(r^n\). Therefore
\(\mathcal H^s_{r^n}(C_r)\leq 2^nr^{ns}\).
This tends to zero for \(s>d\), and equals one for \(s=d\);
it proves the upper dimension bound and finiteness at \(d\).

Construct a probability measure by taking binary digits of
\(t\in[0,1)\) with Lebesgue measure and sending them to
\((1-r)\sum_{j\geq1}\varepsilon_j(t)r^{j-1}\).
The digit functions are measurable; uniform convergence makes the
map measurable. Every prescribed first \(n\) digits occupy a
half-open binary interval of length \(2^{-n}\); the possible
dyadic endpoint choices form a null set. Strong separation for
\(r<1/2\) therefore gives mass exactly \(2^{-n}\) to each level
\(n\) Cantor cylinder interval.

Let \(J\) have length \(u\), and choose \(n\) with
\(r^n\leq u<r^{n-1}\). Distinct level \(n\) intervals have gaps
at least \((1-2r)r^{n-1}\). Consequently \(J\) meets at most
\(C_r'=2+1/(1-2r)\), rounded upward, of them. Hence
\[
 \nu(J)\leq C_r'2^{-n}=C_r'(r^n)^d\leq C_r'u^d .
\]
The mass-distribution lemma gives positivity of \(\mathcal H^d\).
For \(s<d\) any \(\delta\)-cover satisfies
\(\sum_i\operatorname{diam}(U_i)^s
\geq \delta^{s-d}\sum_i\operatorname{diam}(U_i)^d\);
the positive lower bound at \(d\) forces divergence as
\(\delta\downarrow0\), proving the lower dimension bound.

The same mass estimate gives
\(N_\varepsilon(C_r)\geq (C_r')^{-1}\varepsilon^{-d}\).
For \(r^n\leq\varepsilon<r^{n-1}\), the level \(n\) intervals
give \(N_\varepsilon(C_r)\leq2^n\leq2\varepsilon^{-d}\).
Taking logarithms squeezes the box-dimension limit to \(d\).
\end{proof}

\begin{remark}[Limits of the present conclusions]
A numerical rate, a limsup, and an almost-everywhere Lyapunov
exponent are different assertions. Likewise a box-dimension
estimate is not automatically a Hausdorff-dimension theorem.
The general multiplicative ergodic theorem and general dimension
formulae are outside this bounded bridge chapter. The linear,
integrable scalar/volume and separated Cantor results above are
actual proved asymptotic statements within it.
\end{remark}

\paragraph{Sources and coverage.}
The attractor and explicit Cantor themes are compared with Teschl,
Chapters 7, 11 and 13, and Hirsch--Smale--Devaney, Chapters 14--16.
Walters, Sections 1.6 and 10.2, separates ergodic growth theorems
from pointwise computations. Our use of the pointwise theorem is
proved in the next chapter and does not invoke a general
multiplicative theorem without its hypotheses or proof.

% END chapters/M13.18


% BEGIN chapters/M13.19
\chapter{Invariant measures and introductory ergodic theory}\label{ch:d13-19}

\section{Preservation, recurrence and the averaging operator}

\begin{definition}[Measure-preserving and ergodic maps]\label{d13:19:preserving-definition}
Let \((X,\mathcal A,\mu)\) be a probability space. A measurable map
\(T:X\to X\) is measure preserving if
\(\mu(T^{-1}E)=\mu(E)\) for every measurable \(E\).
A measurable set is invariant modulo null sets if
\(\mu(T^{-1}E\mathbin{\triangle}E)=0\).
The map is ergodic if every such set has measure zero or one.
It is mixing if
\(\mu(A\cap T^{-n}B)\to\mu(A)\mu(B)\) for all measurable \(A,B\).
No inverse map is required in these definitions.
\end{definition}

\begin{lemma}[Integration and composition]\label{d13:19:composition-isometry}
For a measure-preserving map, \(\int h\circ T\,\dd\mu=\int h\,\dd\mu\)
for nonnegative measurable \(h\), and for every integrable \(h\).
Consequently \(Uh=h\circ T\) is a well-defined linear isometry on
\(L^p(\mu)\), \(1\leq p\leq\infty\).
The averages \(A_nh=n^{-1}\sum_{j=0}^{n-1}U^jh\) have norm at most one
as operators on each of these spaces.
\end{lemma}
\begin{proof}
The integral identity is the definition of preservation for indicators,
then follows for nonnegative simple functions and increasing limits.
Apply it to positive and negative parts, or real and imaginary parts,
to obtain the integrable case. Null functions remain null because
inverse images of null sets are null. Apply the same identity to
\(\abs{h}^p\). For \(p=\infty\), preservation gives the same measures
to the level sets \(\{\abs{h}>c\}\) and their inverse images.
The triangle inequality proves the averaging bound.
\end{proof}

\begin{theorem}[Poincare recurrence]\label{d13:19:poincare-recurrence}
For a measurable \(E\), almost every \(x\in E\) returns to \(E\)
infinitely many times: \(T^nx\in E\) for infinitely many \(n\geq1\).
The assertion also holds on any finite measure space.
\end{theorem}
\begin{proof}
Let \(B=E\setminus\bigcup_{n\geq1}T^{-n}E\) be the set with no future
return. The sets \(T^{-j}B\), \(j\geq0\), are disjoint: membership in
two, with indices \(i<j\), would send a point of \(B\subset E\)
back into \(E\) after \(j-i\) steps. They all have measure \(\mu(B)\);
finiteness forces \(\mu(B)=0\). A point of \(E\) with only finitely
many returns has a last visit, and some iterate lies in \(B\).
All these points lie in the null set \(\bigcup_{j\geq0}T^{-j}B\).
\end{proof}

\begin{corollary}[Metric recurrence]\label{d13:19:metric-recurrence}
If \(X\) is a second-countable metric space with Borel probability
\(\mu\), almost every \(x\) has a sequence \(n_k\to\infty\) with
\(T^{n_k}x\to x\).
\end{corollary}
\begin{proof}
Apply recurrence to every member of a countable base and remove
the countable union of their exceptional null sets.
For a remaining point, every basis neighborhood containing it is
visited infinitely often. Choose an increasing sequence of visit
times to basis neighborhoods of diameters tending to zero.
\end{proof}

\begin{example}[Finiteness cannot be dropped]
Translation \(T(n)=n+1\) preserves counting measure on \(\Z\), but
the set \(\{0\}\) never returns to itself. This does not contradict
the theorem because the whole space has infinite measure.
\end{example}

\section{A complete Hilbert mean-ergodic argument}

\begin{lemma}[Closed-subspace projection]\label{d13:19:hilbert-projection}
In a real or complex Hilbert space, every closed linear subspace
\(M\) has an orthogonal decomposition \(H=M\oplus M^\perp\).
\end{lemma}
\begin{proof}
For \(x\in H\), let \(d=\inf_{m\in M}\norm{x-m}\), and choose \(m_n\)
with \(\norm{x-m_n}^2\to d^2\). The parallelogram identity gives
\[
 \norm{m_n-m_k}^2
 =2\norm{x-m_n}^2+2\norm{x-m_k}^2
       -4\norm{x-(m_n+m_k)/2}^2\longrightarrow0.
\]
Completeness and closedness give \(m_n\to m\in M\).
Minimality of \(\norm{x-m}\) in the directions \(tv\), \(v\in M\),
implies \(\Re\ip{x-m}{v}=0\); in the complex case use \(iv\) as well.
Thus \(x-m\perp M\). Orthogonality proves uniqueness.
\end{proof}

\begin{theorem}[Mean ergodic theorem for an isometry]\label{d13:19:mean-ergodic}
Let \(U\) be a linear isometry of a Hilbert space into itself, not
necessarily onto. If \(P\) is orthogonal projection onto
\(\operatorname{Fix}U=\{h:Uh=h\}\), then
\[
 \frac1n\sum_{j=0}^{n-1}U^jh\longrightarrow Ph
       \quad\hbox{in Hilbert norm}.
\]
In particular this applies to composition by a probability-preserving
map on \(L^2(\mu)\).
\end{theorem}
\begin{proof}
Put \(M=\overline{\operatorname{ran}(I-U)}\).
Its orthogonal complement is \(\operatorname{Fix}U\).
Indeed, \(h\perp M\) means
\(\ip{h}{Uv}=\ip{h}{v}\) for all \(v\). Taking \(v=h\) and using
\(\norm{Uh}=\norm{h}\) gives
\(\norm{Uh-h}^2=0\), so \(Uh=h\).
Conversely, if \(Uh=h\), preservation of inner products by an
isometry gives
\(\ip{h}{Uv}=\ip{Uh}{Uv}=\ip{h}{v}\).
The projection lemma therefore gives
\(H=\operatorname{Fix}U\oplus M\).

On the fixed subspace the averages are the identity. On a range
vector \(h=(I-U)g\), telescoping gives
\(A_nh=(g-U^ng)/n\), of norm at most \(2\norm{g}/n\).
Since \(\norm{A_n}\leq1\), approximation extends convergence to
the closure \(M\). Add the two components. For \(L^2\), the required
Hilbert completeness is the usual \(L^2\) completeness theorem
from the measure-and-integration foundations; the operator is
an isometry by Lemma~\ref{d13:19:composition-isometry}.
\end{proof}

\section{The maximal estimate and pointwise convergence}

\begin{lemma}[Finite and infinite maximal ergodic estimate]\label{d13:19:maximal-lemma}
Let \(f\in L^1(\mu)\) be real and put
\(S_kf=\sum_{j=0}^{k-1}f\circ T^j\).
For \(E_N=\{\max_{1\leq k\leq N}S_kf>0\}\),
\(\int_{E_N}f\,\dd\mu\geq0\).
The same holds for \(E=\{\sup_{k\geq1}S_kf>0\}\).
Consequently, for \(h\in L^1\) and \(\alpha>0\),
\[
 \mu\{\sup_{n\geq1}A_n\abs{h}>\alpha\}
       \leq\frac{\norm{h}_1}{\alpha}.
\]
\end{lemma}
\begin{proof}
Let \(M_N=\max(0,S_1f,\ldots,S_Nf)\). On \(E_N\) the identity
\(S_kf=f+S_{k-1}f\circ T\) gives
\(M_N\leq f+M_N\circ T\).
Outside \(E_N\), \(M_N=0\), so in both cases
\[
 f1_{E_N}\geq M_N-M_N\circ T.
\]
The function \(M_N\) is integrable, bounded by
\(\sum_{j=0}^{N-1}\abs{f}\circ T^j\). Integrating and using preservation
proves the finite assertion. The sets \(E_N\) increase to \(E\);
dominated convergence applied to \(f1_{E_N}\) proves the infinite one.
Apply it to \(f=\abs{h}-\alpha\): its set \(E\) is exactly the
maximal-average set in the display, and
\(\alpha\mu(E)\leq\int_E\abs{h}\leq\norm{h}_1\).
\end{proof}

\begin{lemma}[Dense bounded coboundaries]\label{d13:19:dense-coboundaries}
In real \(L^1\) on a probability space, the class
\[
 \mathcal D=\{v+(I-U)g:\ v,g\in L^\infty,\ Uv=v\}
\]
is dense. For \(d=v+(I-U)g\) in this class, \(A_nd\to v\) almost
everywhere and in \(L^1\).
\end{lemma}
\begin{proof}
The orthogonal decomposition in the mean ergodic proof says that
fixed \(L^2\) functions and coboundaries with \(L^2\) transfer
functions have dense span in \(L^2\).
Truncating a fixed real function gives bounded fixed functions,
converging in \(L^2\). Truncating an arbitrary \(L^2\) transfer
function and using \(\norm{I-U}\leq2\) approximates its coboundary
by bounded ones. Thus \(\mathcal D\) is dense in \(L^2\), hence in
\(L^1\): bounded truncations show \(L^2\) is \(L^1\)-dense, and on
a probability space \(\norm{h}_1\leq\norm{h}_2\).
Telescoping gives
\(A_nd=v+(g-U^ng)/n\); the error is at most
\(2\norm{g}_\infty/n\) outside a common null set.
\end{proof}

\begin{theorem}[Pointwise and \(L^1\) ergodic theorem]\label{d13:19:birkhoff}
For a probability-preserving \(T\) and \(f\in L^1(\mu)\), the averages
\[
 A_nf(x)=\frac1n\sum_{j=0}^{n-1}f(T^jx)
\]
converge almost everywhere and in \(L^1\) to an integrable function
\(f^*\). It satisfies \(f^*\circ T=f^*\) almost everywhere and
\[
 \int_E f^*\,\dd\mu=\int_E f\,\dd\mu
\]
for every invariant measurable set \(E\), in particular for \(E=X\).
If \(T\) is ergodic, \(f^*=\int f\,\dd\mu\) almost everywhere.
Complex-valued functions are treated componentwise.
\end{theorem}
\begin{proof}
First let \(f\) be real. Choose \(d_k\in\mathcal D\) with
\(\norm{f-d_k}_1\leq2^{-3k}\). Since
\(\abs{A_n(f-d_k)}\leq A_n\abs{f-d_k}\), the maximal lemma gives
\[
 \mu\{\sup_n\abs{A_n(f-d_k)}>2^{-k}\}\leq2^{-2k}.
\]
The sum is finite, so almost every point belongs to only finitely
many of these exceptional sets. This elementary Borel--Cantelli
conclusion follows directly from bounding the measure of a union
of all sets after index \(k\) by the tail of their measures.
At a point also satisfying convergence for every \(d_k\),
the upper and lower limits of \(A_nf\) differ by at most \(2^{1-k}\)
for every sufficiently large \(k\). They are finite, because for
one such \(k\) the entire sequence is uniformly close to the
convergent sequence \(A_nd_k\). Thus \(A_nf\to f^*\) almost everywhere.

Write \(d_k^*=\lim_nA_nd_k\). Fatou's lemma and the \(L^1\)
contraction bound give
\[
 \norm{f^*-d_k^*}_1
 \leq\liminf_n\norm{A_n(f-d_k)}_1\leq\norm{f-d_k}_1.
\]
In particular \(f^*\) is integrable, and
\[
 \limsup_n\norm{A_nf-f^*}_1\leq2\norm{f-d_k}_1.
\]
Letting \(k\to\infty\) proves \(L^1\) convergence.

To check invariance, for every \(\varepsilon>0\),
\[
 \sum_{n\geq1}\mu\{\abs{f(T^nx)}>\varepsilon n\}
 =\sum_{n\geq1}\mu\{\abs{f}>\varepsilon n\}
 \leq\varepsilon^{-1}\norm{f}_1.
\]
The same tail argument gives \(f(T^nx)/n\to0\) almost everywhere.
Now
\(A_nf(Tx)-A_nf(x)=(f(T^nx)-f(x))/n\to0\).
The inverse image of the convergence full-measure set is still
full measure; taking limits proves \(f^*\circ T=f^*\).

For invariant \(E\), its indicator satisfies
\(1_E\circ T^j=1_E\) almost everywhere. Therefore
\[
 \int_E f\circ T^j\,\dd\mu
 =\int (1_Ef)\circ T^j\,\dd\mu=\int_Ef\,\dd\mu.
\]
Average and pass to the \(L^1\) limit.
Finally every invariant real measurable function is almost surely
constant in an ergodic system: each rational superlevel set has
measure zero or one, and these nested level sets determine a single
finite value for an almost surely finite function.
Applying this to \(f^*\) and using its integral determines the
constant. Real and imaginary parts complete the complex case.
\end{proof}

\begin{corollary}[An equivalent ergodicity criterion]\label{d13:19:ergodic-criterion}
A probability-preserving map is ergodic if and only if for every
integrable \(f\), its time averages converge almost everywhere to
the constant \(\int f\,\dd\mu\). Mixing implies ergodicity.
\end{corollary}
\begin{proof}
The forward implication is the theorem. Conversely for invariant
\(E\), the averages of \(1_E\) are \(1_E\) almost everywhere.
The asserted constant limit forces \(\mu(E)\in\{0,1\}\).
If the map is mixing, apply its definition with \(A=B=E\) invariant:
the left side is always \(\mu(E)\), and the limit is \(\mu(E)^2\).
\end{proof}

\section{Invariant probabilities on a compact metric space}

\begin{lemma}[Weak compactness of probabilities]\label{d13:19:probability-compactness}
On a nonempty compact metric space \(K\), every sequence of Borel
probability measures has a subsequence for which the integrals of
every real continuous function converge to those for a Borel
probability measure.
\end{lemma}
\begin{proof}
The space \(C(K,\R)\) is separable in uniform norm. One concrete
proof first approximates a continuous \(f\) uniformly by the
Lipschitz functions
\(\inf_y(f(y)+L\dist(x,y))\) as \(L\to\infty\).
Uniform continuity proves this convergence: near \(x\) the loss
in \(f(y)\) is small, while outside a fixed small ball the term
\(L\dist(x,y)\) dominates \(2\norm{f}_\infty\).
Using finite nets from a countable dense subset of \(K\), rational
function values and integer Lipschitz constants then gives a
countable uniformly dense family of functions of the form
\(\min_i(q_i+L\dist(x,x_i))\).

Diagonal extraction on that countable family gives a subsequence
whose integrals converge there. The bound
\(\abs{\int f\,\dd\nu}\leq\norm{f}_\infty\) extends convergence to
every continuous \(f\). The resulting functional \(I\) is linear,
positive, satisfies \(I(1)=1\), and is bounded.
Invoke the positive Riesz--Markov representation theorem from the
measure-and-integration foundations: a positive real linear
functional on \(C_c(K,\R)\) is represented uniquely by a regular
Borel measure on locally compact Hausdorff \(K\).
Here compactness identifies \(C_c(K,\R)=C(K,\R)\), and \(I(1)=1\)
makes the representing measure a probability. This is the
measure representation theorem, not finite-dimensional inner-product
representation.
\end{proof}

\begin{theorem}[Existence of an invariant probability]\label{d13:19:invariant-probability}
Every continuous map \(T:K\to K\) on a nonempty compact metric
space has an invariant Borel probability measure.
\end{theorem}
\begin{proof}
Choose \(x\in K\) and the empirical probabilities
\(\nu_n=n^{-1}\sum_{j=0}^{n-1}\delta_{T^jx}\).
Take a subsequential weak limit \(\nu\) using the lemma.
For \(g\in C(K,\R)\), also \(g\circ T\) is continuous, and
\[
 \int(g\circ T-g)\,\dd\nu_n
     =\frac{g(T^nx)-g(x)}n\longrightarrow0 .
\]
Thus \(\int g\circ T\,\dd\nu=\int g\,\dd\nu\).
The pushforward \(T_*\nu\) is a regular Borel probability on this
compact metric space; uniqueness in Riesz--Markov gives
\(T_*\nu=\nu\), which is precisely preservation.
\end{proof}

\begin{theorem}[Unique ergodicity gives uniform continuous averages]\label{d13:19:unique-ergodic-uniform}
If the continuous map \(T:K\to K\) above has a unique invariant
Borel probability \(\nu\), then for every \(g\in C(K,\R)\),
\[
 \sup_{x\in K}\left|\frac1n\sum_{j=0}^{n-1}g(T^jx)
                         -\int g\,\dd\nu\right|\longrightarrow0 .
\]
\end{theorem}
\begin{proof}
Otherwise choose \(n_k\to\infty\) and \(x_k\) for which the absolute
difference is at least a fixed \(\varepsilon>0\). Apply weak
compactness to the empirical measures based at \(x_k\) for time
\(n_k\). Every weak subsequential limit is invariant, because the
telescoping error for any continuous \(h\) is bounded uniformly by
\(2\norm{h}_\infty/n_k\). It must equal \(\nu\), contradicting the
chosen discrepancy for \(g\).
\end{proof}

\section{Rotations and shifts: different forms of recurrence}

\begin{theorem}[Irrational rotations]\label{d13:19:rotation-ergodicity}
Lebesgue probability on \(\mathbb T=\R/\Z\) is preserved by
\(R_\alpha(x)=x+\alpha\). If \(\alpha\) is irrational, this rotation
is ergodic and uniquely ergodic, but not mixing.
For every continuous periodic \(g\), its orbit averages converge
uniformly to \(\int_0^1g(x)\,\dd x\).
\end{theorem}
\begin{proof}
Translation modulo one preserves the lengths of half-open arcs;
the uniqueness part of measure extension gives preservation of
Lebesgue measure. If \(h\in L^2\) is invariant, change variables
in its \(k\)-th Fourier coefficient to get
\[
 \widehat h(k)=e^{2\pi i k\alpha}\widehat h(k).
\]
For nonzero integer \(k\), irrationality makes the factor different
from one; thus all nonconstant coefficients vanish.
Trigonometric polynomials are dense in \(L^2(\mathbb T)\):
continuous functions are \(L^2\)-dense by regularity of Lebesgue
measure, and their Fejer polynomials converge uniformly.
These are the established integration and periodic Fourier
foundations. Hence \(h\) equals its constant mean in \(L^2\).
Apply this to invariant indicators to obtain ergodicity.

For any invariant Borel probability \(\eta\), the same computation
with the continuous characters \(e^{2\pi ikx}\) shows that its
integrals of all trigonometric polynomials agree with Lebesgue
probability. Uniform Fejer approximation extends equality to all
continuous functions, and measure representation uniqueness gives
\(\eta=\mu\). Thus the invariant measure is unique, and the
preceding uniform-averages theorem applies.

Finally take a nontrivial arc \(A\), \(0<\mu(A)<1\), and positive
integers \(q_j\to\infty\) with \(q_j\alpha\to0\) modulo one, as in
Theorem~\ref{d13:15:rigid-rotation}. The measure of the symmetric
difference of an arc and its translate tends to zero. Therefore
\(\mu(A\cap R_\alpha^{-q_j}A)\to\mu(A)\), whereas mixing would
force the strictly smaller value \(\mu(A)^2\).
\end{proof}

\begin{lemma}[A Bernoulli probability on the shift]\label{d13:19:bernoulli-measure}
On either binary symbolic space there is a unique Borel probability
assigning mass \(2^{-m}\) to each cylinder specifying \(m\) distinct
coordinates. Coordinates on disjoint finite index sets are independent.
\end{lemma}
\begin{proof}
On the algebra of finite unions of cylinders the stated values give
a consistent finitely additive probability: refine to a common finite
set of coordinates and count the equally weighted assignments.
This is a premeasure. Indeed, if a member of the algebra is a
disjoint countable union of algebra members, these clopen sets form
an open cover of a compact set. A finite subcover exists, so all
other members of the disjoint union are empty; finite additivity
already proves the required countable equality.
Caratheodory's extension theorem gives a probability on the
generated sigma algebra, which is the Borel sigma algebra because
the countable cylinders form a base. Uniqueness follows from
finiteness of the premeasure. Refining two finite coordinate
conditions to their disjoint union proves independence.
\end{proof}

\begin{theorem}[Bernoulli shifts are mixing]\label{d13:19:bernoulli-mixing}
The binary shifts preserve their Bernoulli probability and are mixing,
hence ergodic. Almost every sequence has frequency \(1/2\) of the
symbol one in its nonnegative coordinates. Periodic sequences are
dense topologically but have total probability zero.
\end{theorem}
\begin{proof}
The inverse image of a cylinder specifies the same number of
coordinates shifted by one, so its measure is unchanged.
Uniqueness of finite measure extension proves preservation on
all Borel sets. For cylinders, or finite unions of them, the
coordinate sets used by \(A\) and \(\sigma^{-n}B\) are disjoint
for all sufficiently large \(n\). Independence gives
\(\mu(A\cap\sigma^{-n}B)=\mu(A)\mu(B)\).

Every measurable set can be approximated in measure by the cylinder
algebra. To see this without an extra compactness assumption, the
class of sets so approximable is a sigma algebra: it is closed
under complements and finite unions, and a countable union is
approximated first by a finite union using continuity from below,
then by algebra members. It contains the generating algebra.
For arbitrary \(A,B\), choose algebra approximants with symmetric
differences of measure less than \(\varepsilon\). Preservation
bounds the difference of intersection probabilities by \(2\varepsilon\)
and that of products by \(2\varepsilon\).
Let \(n\to\infty\), then \(\varepsilon\downarrow0\), to prove mixing.
Ergodicity and the frequency assertion follow from
Corollary~\ref{d13:19:ergodic-criterion} and Birkhoff applied to
the indicator of the first symbol.
Each individual sequence has probability at most \(2^{-m}\) for
every \(m\), hence zero. Periodic sequences form a countable
union of finite sets, so their total probability is zero;
their density was proved topologically in
Theorem~\ref{d13:17:full-shift-chaos}.
\end{proof}

\begin{remark}[What the introduction does and does not claim]
Recurrence is an almost-everywhere conclusion under finite invariant
measure. Ergodicity identifies almost-everywhere time averages with
space averages; mixing is stronger, and uniform convergence for
continuous observables here uses unique ergodicity and compactness.
These statements have different quantifiers and are not substitutes
for one another. Entropy theory, equilibrium states and the general
multiplicative ergodic theorem require additional chapters and are
not adopted as unproved extensions of these results.
\end{remark}

\paragraph{Sources and coverage.}
Walters, Chapters 0--1, 3, 5--6, provides the measure-preserving,
ergodic, rotation and invariant-measure branches. In particular
Section 1.6 is the source branch for the pointwise and maximal
ergodic theorems. The present dense-coboundary proof spells out
the \(L^2\)-to-\(L^1\) and maximal-estimate bridges instead of
identifying mean convergence with pointwise convergence.
Compactness, measure extension, integral convergence, \(L^2\)
completeness, positive Riesz--Markov representation and Fejer
approximation are prior analysis/topology foundations with their
stated hypotheses; the dynamical conclusions are proved above.

% END chapters/M13.19


% BEGIN chapters/M13.20
\chapter{Hyperbolic local theory and bounded Green operators}
\label{ch:d13-20}

This foundation chapter fills the hyperbolic local theory used by the
stability and return-map chapters. A local topological conjugacy is weaker
than a differentiable conjugacy. In particular, the theorem below does not
remove nonlinear terms by a smooth change of coordinates at resonance.

\section{Splitting and controlled extension}
\begin{definition}[Hyperbolicity]\label{d13:20:hyperbolicity}
A real matrix $A$ is hyperbolic if none of its complex eigenvalues has real
part zero. Write $\R^n=E^s\oplus E^u$ for the real invariant sums corresponding
to eigenvalues with negative and positive real parts, and $P^s,P^u$ for the
commuting projections. Either subspace is allowed to be zero.
\end{definition}

\begin{lemma}[Exponential dichotomy]\label{d13:20:dichotomy}
For hyperbolic $A$ there are $K\geq1$ and $\alpha>0$ such that
\[
 \norm{e^{tA}P^s}\leq Ke^{-\alpha t},\qquad
 \norm{e^{-tA}P^u}\leq Ke^{-\alpha t}\quad(t\geq0).
\]
Every bounded solution of $z'=Az$ on all of $\R$ is zero.
\end{lemma}
\begin{proof}
In a complex Jordan basis each entry is a fixed polynomial in $t$ times
$e^{\lambda t}$. Choose $\alpha$ strictly smaller than the least absolute
nonzero real part of an eigenvalue. A polynomial times the remaining
decaying exponential is bounded on $[0,\infty)$, giving the two estimates
after changing basis and restricting to the real invariant spaces.
If $z(t)=e^{tA}z_0$ is bounded by $M$, then
$P^u z_0=e^{-tA}P^u z(t)$ tends to zero as $t\to\infty$.
Similarly $P^s z_0=e^{tA}P^s z(-t)$ tends to zero. Their sum is $z_0$.
\end{proof}

\begin{lemma}[A small bounded nonlinear extension]\label{d13:20:cutoff}
Suppose $f$ is $C^k$, $1\leq k\leq\infty$, near zero, $f(0)=0$, and
$Df(0)=A$. For every $\delta>0$ there is a global $C^k$ map $r:\R^n\to\R^n$
with compact support, $r(0)=Dr(0)=0$, $\Lip(r)\leq\delta$, such that
$Ax+r(x)=f(x)$ in a neighborhood of zero. All derivatives of any fixed
order at most $k$ are bounded and uniformly continuous. The modified
equation has a complete global flow.
\end{lemma}
\begin{proof}
Put $g(x)=f(x)-Ax$. Choose a ball of radius $2\rho$ in the original domain
where $\norm{Dg}\leq\varepsilon$. Since $g(0)=0$, integration along the
radius gives $\norm{g(x)}\leq\varepsilon\norm{x}$ there. Take a smooth cutoff
$\chi(x/\rho)$ equal to one on the radius-$\rho$ ball, zero outside the
radius-$2\rho$ ball, with derivative norm at most $C/\rho$. Extend
$r(x)=\chi(x/\rho)g(x)$ by zero. It is $C^k$ across the outer boundary
because the cutoff is already zero in a neighborhood on the exterior side.
Its derivative is bounded by $(1+2C)\varepsilon$. Set
$\varepsilon\leq\delta/(1+2C)$. The assertions at zero follow from those
for $g$. Compact support and continuity of every derivative give the
boundedness and uniform continuity assertions. Finally $Ax+r(x)$ is
globally Lipschitz and has linear growth; the continuation and growth
results of Chapter~\ref{ch:d13-04} give the complete flow in both directions.
\end{proof}

\section{Two-sided bounded solutions}
Let $C_b(\R,\R^n)$ denote bounded continuous maps with the supremum norm.
Its completeness follows by pointwise completeness and uniform convergence.

\begin{lemma}[The bounded Green inverse]\label{d13:20:green}
For $g\in C_b(\R,\R^n)$ define
\[
 (Gg)(t)=\int_{-\infty}^t e^{(t-s)A}P^s g(s)\,\dd s
          -\int_t^\infty e^{(t-s)A}P^u g(s)\,\dd s.
\]
Then $Gg$ is the unique bounded $C^1$ solution of $z'=Az+g$ and
$\norm{Gg}_\infty\leq C_G\norm{g}_\infty$, where $C_G=2K/\alpha$.
\end{lemma}
\begin{proof}
The dichotomy bounds both improper integrals absolutely and uniformly in
their tails after shifting to $s-t$. They give the stated norm bound and
continuity by dominated convergence. Differentiating truncated integrals
and passing to the uniformly controlled tails yields
$(Gg)'=A(Gg)+(P^s+P^u)g=A(Gg)+g$. Alternatively the same identity follows
first from the integral variation-of-constants formula on each compact
time interval and then from the fundamental theorem of calculus.
The difference of two bounded solutions is a bounded homogeneous solution
on all of $\R$, hence zero by Lemma~\ref{d13:20:dichotomy}.
\end{proof}

\section{Stable and unstable graphs}
Choose $0<\beta<\alpha$, and choose the extension in
Lemma~\ref{d13:20:cutoff} so small that
\[
 q_\beta=K\delta\left(\frac1{\alpha-\beta}
                         +\frac1{\alpha+\beta}\right)<1.
\]
For $t\geq0$ set
\[
 X_\beta=\{w\in C([0,\infty),\R^n):
                 \norm{w}_\beta=\sup_{t\geq0}e^{\beta t}\norm{w(t)}<\infty\}.
\]
Multiplication by $e^{\beta t}$ is an isometry onto a complete bounded
continuous-function space, so $X_\beta$ is a Banach space.

\begin{lemma}[Lyapunov--Perron fixed points and their regularity]
\label{d13:20:perron}
For every $\xi\in E^s$ there is a unique $w_\xi\in X_\beta$ satisfying
\begin{align*}
 w(t)={}&e^{tA}\xi+
       \int_0^t e^{(t-s)A}P^s r(w(s))\,\dd s\\
       &-\int_t^\infty e^{(t-s)A}P^u r(w(s))\,\dd s.
\end{align*}
It solves $w'=Aw+r(w)$, satisfies $P^s w(0)=\xi$, and
\[
 \norm{w_\xi}_\beta\leq\frac{K}{1-q_\beta}\norm{\xi}.
\]
The map $\xi\mapsto w_\xi$ into $X_\beta$ is $C^k$.
\end{lemma}
\begin{proof}
The linear term has weighted norm at most $K\norm{\xi}$. For the stable
integral, multiply by $e^{\beta t}$ and use
\[
 e^{\beta t}\int_0^t Ke^{-\alpha(t-s)}
       \delta e^{-\beta s}\norm{w-v}_\beta\,\dd s
 \leq\frac{K\delta}{\alpha-\beta}\norm{w-v}_\beta.
\]
For the unstable tail the analogous estimate is
$K\delta/(\alpha+\beta)$. Thus the right side is a contraction of the whole
complete space. Its value at zero is $e^{tA}\xi$, proving existence,
uniqueness and the norm bound. Differentiation of the integral equation,
justified as in Lemma~\ref{d13:20:green}, proves the differential equation.
Applying $P^s$ at zero gives the initial projection.

Here is the parameter regularity argument, including its infinite-time
point. Denote the nonlinear integral operator by $N(w)$. Pointwise Taylor
expansion and uniform continuity of $Dr$ give
\[
 \norm{r(w+v)-r(w)-Dr(w)v}_\beta
 \leq\omega(\norm{v}_\infty)\norm{v}_\beta,
\]
where $\omega(a)\to0$ as $a\downarrow0$; $\norm{v}_\infty\leq\norm{v}_\beta$.
The same integral bounds make $N$ Fr\'echet differentiable and
$\norm{DN(w)}\leq q_\beta$. For order $j\leq k$, the pointwise derivative
$D^j r(w)[v_1,\ldots,v_j]$ lies in $X_\beta$, because
$e^{\beta t}\prod_i\norm{v_i(t)}\leq\prod_i\norm{v_i}_\beta$.
Boundedness and uniform continuity of $D^j r$ prove continuity of these
multilinear maps and the Taylor remainder statement at each order. Hence
$N$ is $C^k$ on $X_\beta$.

For completeness, no implicit-function theorem in an unspecified Banach
space is needed. Subtracting the fixed-point equations first gives
$\norm{w_{\xi+h}-w_\xi}_\beta\leq K\norm{h}/(1-q_\beta)$.
The candidate first derivative in direction $h$ is the unique solution
\[
 V=e^{tA}h+DN(w_\xi)V,
 \qquad V=\sum_{j\geq0}DN(w_\xi)^j(e^{tA}h).
\]
The series converges in operator norm, uniformly with bound
$1/(1-q_\beta)$. Subtract this equation from the finite difference equation;
the Taylor remainder divided by $\norm{h}$ tends to zero, and the same
inverse bound controls the difference. Continuity follows from the
resolvent identity and continuity of $DN$. Differentiating again expresses
the next derivative as the same bounded inverse applied to finite sums of
lower derivatives and derivatives of $N$. Induction gives every order up
to $k$, including continuity; each induction step uses the preceding
Taylor remainder estimate. For $k=\infty$ apply this to each finite order.
\end{proof}

\begin{theorem}[Local stable and unstable manifolds]
\label{d13:20:stable-manifold}
Let a $C^k$ real vector field, $1\leq k\leq\infty$, have a hyperbolic zero
at $p$. In coordinates centered at $p$ there are $C^k$ graphs
\[
 W^s_{\rm graph}=\{\xi+h^s(\xi):\xi\in E^s\},\qquad
 W^u_{\rm graph}=\{\eta+h^u(\eta):\eta\in E^u\}
\]
for a small extension agreeing with the original field near $p$.
Here $h^s:E^s\to E^u$, $h^u:E^u\to E^s$ vanish with their derivatives at
zero. The graphs are respectively forward and backward invariant for the
extension. Their sufficiently small portions give local invariant
manifolds for the original field. Their trajectories converge to $p$
exponentially in the appropriate time direction.

More precisely, choose a ball $U$ on which the extension equals the
original field. Every original forward trajectory remaining in a
sufficiently small bounded neighborhood of $p$ for all $t\geq0$ starts
on the stable graph. Conversely all graph points with sufficiently small
$\norm{\xi}$ stay in $U$ for all $t\geq0$ and satisfy
$\norm{\Phi(t,p+\xi+h^s(\xi))-p}\leq C\norm{\xi}e^{-\beta t}$.
The reversed-time statement holds for the unstable graph.
\end{theorem}
\begin{proof}
Use the extension and set $h^s(\xi)=P^u w_\xi(0)$. Evaluation at zero is a
bounded linear map on $X_\beta$, so $h^s$ is $C^k$. The zero fixed point
for $\xi=0$ gives $h^s(0)=0$. Since $Dr(0)=0$, the derivative equation at
zero is $V(t)=e^{tA}h$ for $h\in E^s$. Its unstable projection at zero is
zero, proving $Dh^s(0)=0$. The stable projection of a graph point is its
parameter, making the graph an embedded manifold tangent to $E^s$.

For $\tau\geq0$, the shifted trajectory $w_\xi(t+\tau)$ is in $X_\beta$.
Integrating its stable equation forward and its unstable equation to
infinity gives the same fixed-point formula with initial stable
projection $P^s w_\xi(\tau)$. Uniqueness proves forward invariance.
The weighted bound proves decay and shows that choosing
$K\norm{\xi}/(1-q_\beta)$ smaller than the radius of $U$ keeps the whole
trajectory in the original-field ball.

For the converse, a trajectory remaining in that ball is bounded for all
positive time. Its unstable variation-of-constants equation, solved
backwards from $T$, has a terminal term bounded by
$Ke^{-\alpha(T-t)}\sup\norm{w}$, tending to zero. Its stable equation is
integrated from zero. Thus it satisfies the fixed-point equation in
$X_0=C_b([0,\infty),\R^n)$. The corresponding contraction constant is
$q_0=2K\delta/\alpha\leq q_\beta<1$. Since $w_\xi$ already lies in $X_0$,
uniqueness there makes the bounded trajectory equal to it. In particular
it has the asserted exponential decay; decay was not an extra assumption
used to identify the graph. Replace $t$ by $-t$ and $A,r$ by $-A,-r$ for
the unstable assertions.
\end{proof}

\begin{remark}[Local versus global stable sets]\label{d13:20:local-global}
A full stable set can return from far away and need not coincide with
one small graph throughout an arbitrary neighborhood. The characterization
above uses trajectories remaining in the chosen local neighborhood.
Taking backward flow images of local stable portions gives the global
stable manifold as an immersed manifold, not necessarily a closed
embedded submanifold of the entire ambient space.
Here is the chart argument for that last assertion. Within the modified
stable graph take the open parameter set
$V_s=\{\xi:\norm{w_\xi}_\beta<\eta\}$ for small $\eta$.
The shift estimate
$\norm{w_\xi(\cdot+\tau)}_\beta\leq e^{-\beta\tau}\norm{w_\xi}_\beta$
makes its image forward invariant. The sets obtained by applying
$\Phi(-m,\cdot)$ to this image on its actual backward-existence domain,
for integers $m\geq0$, are nested
manifold charts covering the global stable set: every convergent
trajectory eventually stays in the original small ball, is on the graph
by the theorem, and eventually has arbitrarily small stable parameter.
Finite-time flow maps are $C^k$ local diffeomorphisms. On overlapping
stable-graph parameter charts their derivatives have full stable rank,
so the finite-dimensional inverse theorem gives $C^k$ chart transitions
and open inclusions. The resulting countable atlas is Hausdorff because
its inclusion in $\R^n$ is continuous and injective, and it is second
countable from the countably many Euclidean parameter charts. Its
inclusion has injective differential in every chart, proving immersion.
The atlas topology need not equal the ambient subspace topology.
\end{remark}

\section{Topological linearization without smooth-conjugacy claims}
\begin{theorem}[Hartman--Grobman for flows]\label{d13:20:hartman-grobman}
Let $f\in C^1$ near a hyperbolic zero $p$, and put $A=Df(p)$.
There is a homeomorphism $H$ between neighborhoods of zero and $p$,
$H(0)=p$, which conjugates $x'=Ax$ to $y'=f(y)$ for every time interval
on which the two corresponding trajectories remain in these
neighborhoods. Time is preserved. No differentiability or Lipschitz
regularity of $H$ is asserted.
\end{theorem}
\begin{proof}
Translate $p$ to zero and choose the bounded extension with
$C_G\delta<1$. Write its complete flow as $\phi$. For each $x\in\R^n$
the map
\[
 z\longmapsto G\bigl(r(e^{(\cdot)A}x+z(\cdot))\bigr)
\]
is a contraction of $C_b(\R,\R^n)$. The argument of $r$ need not be bounded;
boundedness of $r$ ensures the forcing is bounded, and its Lipschitz
bound gives contraction constant $C_G\delta$. Let $z_x$ be its unique
fixed point and set $H(x)=x+z_x(0)$. Uniformly in $x$,
$\norm{z_x}_\infty\leq C_G\norm{r}_\infty=:M$.
The Green identity shows
$e^{tA}x+z_x(t)$ solves the modified nonlinear equation. Translating time
in its defining integral equation and using uniqueness gives
\[
 z_{e^{\tau A}x}(t)=z_x(t+\tau),\qquad
 H(e^{\tau A}x)=\phi(\tau,H(x)).
\]

To prove bijectivity, start with any nonlinear orbit $y(t)=\phi(t,y_0)$.
The function $v=G(r(y))$ is bounded and $x(t)=y(t)-v(t)$ solves $x'=Ax$.
Thus $x(t)=e^{tA}x_0$, and $v$ solves the fixed-point equation defining
$z_{x_0}$. It follows that $H(x_0)=y_0$, proving surjectivity.
If $H(x_1)=H(x_2)$, uniqueness of the nonlinear orbit gives
$e^{tA}(x_1-x_2)=z_{x_2}(t)-z_{x_1}(t)$ for all real $t$.
This homogeneous orbit is bounded, so $x_1=x_2$.

Continuity requires care because $e^{tA}x_j$ need not converge uniformly
on all real time when $x_j\to x$. Instead $z_{x_j}$ are uniformly bounded
by $M$, and their equations give the uniform derivative bound
$\norm{z_{x_j}'}\leq\norm{A}M+\norm{r}_\infty$. Arzel\`a--Ascoli on
$[-m,m]$, followed by a diagonal subsequence, gives locally uniform
convergence along a subsequence to a bounded continuous $z$. On every
fixed compact interval, $e^{tA}x_j\to e^{tA}x$ uniformly. In the Green
integral, tails are uniformly bounded by an integrable decaying
exponential, while the remaining compact part converges uniformly.
Therefore $z=G(r(e^{(\cdot)A}x+z))$, and uniqueness gives $z=z_x$.
Every subsequence has a further subsequence with this same limit at zero,
so $H(x_j)\to H(x)$.

The inverse is explicitly
\[
 H^{-1}(y)=y-[G(r(\phi(\cdot,y)))](0).
\]
Continuous dependence of $\phi(t,y)$ is uniform on each compact time
interval locally in $y$, and $r$ is bounded. Splitting the integral into
an arbitrarily small uniformly bounded tail and a compact interval proves
continuity of this inverse. For $x=0$ the zero function is the unique fixed
point, so $H(0)=0$.

We have obtained a global conjugacy for the small bounded extension.
Restrict its homeomorphism to a neighborhood of zero whose image lies in
the ball where the extension equals $f$. Along trajectory segments
remaining in these neighborhoods the modified and original equations
coincide by uniqueness. This is exactly the local assertion.
\end{proof}

\begin{example}[A nonlinear stable graph]\label{d13:20:example-graph}
For $x'=-x+y^2$, $y'=y$, the stable graph at zero is $y=0$.
The unstable graph is $x=y^2/3$: differentiating this identity along a
trajectory gives $x'=2y^2/3=-y^2/3+y^2$. In fact
$y(t)=y_0e^t$ and
$x(t)=e^{-t}(x_0-y_0^2/3)+(y_0^2/3)e^{2t}$.
Decay as $t\to-\infty$ holds precisely when $x_0=y_0^2/3$.
Thus the unstable graph is tangent to the unstable axis but is not that
linear axis. The formula also checks both the nonlinear graph and its
orientation in time directly.
\end{example}

\source{T, Sections 9.1--9.4 (hyperbolic splitting, invariant graphs,
bounded integral equations and Hartman--Grobman); P, Sections 2.7--2.8.
The two-sided Green construction and the explicit parameter remainder
argument give complete bridges needed by the selected local results.}

% END chapters/M13.20


% BEGIN chapters/M13.21
\chapter{Separation by a regular closed plane curve}\label{ch:d13-21}

This is a bounded foundational supplement for the planar dynamics chapter.
The required assertion concerns regular $C^1$ embedded closed curves, not
arbitrary possibly wild continuous Jordan curves. It is proved here rather
than treated as an already available theorem in the reusable library.
The plane in this chapter has its Euclidean norm and its standard orientation;
write $J(x_1,x_2)=(-x_2,x_1)$ and
$\det(u,v)=u_1v_2-u_2v_1$.

\section{Elementary winding facts}
\begin{lemma}[Index properties without a separation theorem]\label{d13:21:index-properties}
For a continuous loop $\alpha:[0,1]\to\R^2\setminus\{p\}$,
with $\alpha(0)=\alpha(1)$, its winding index about $p$ is an integer.
It is unchanged under any continuous loop homotopy avoiding $p$.
For fixed $\alpha$ it is locally constant as a function of $p$ off
the trace, and it is zero for $p$ outside a sufficiently large ball.
For a $C^1$ loop it equals
\[
 \operatorname{ind}(\alpha,p)=\frac{1}{2\pi}
 \int_0^1\frac{\det(\alpha(t)-p,\alpha'(t))}
                    {\norm{\alpha(t)-p}^2}\,\dd t.
\]
\end{lemma}
\begin{proof}
Identify the oriented plane with $\C$. A continuous nonzero path admits
a continuous argument on $[0,1]$: its normalized image is covered by
open semicircles, uniform continuity gives a finite interval subdivision
whose pieces lie in such semicircles, and the argument on successive
pieces is adjusted by multiples of $2\pi$ to agree at the subdivision
points. Two such lifts differ by a constant multiple of $2\pi$, since
their difference is continuous and belongs to $2\pi\Z$. For a loop
the endpoint difference is therefore $2\pi k$ for a uniquely determined
integer $k$; define the index to be $k$.

If two loops $\alpha,\beta$ satisfy
$\sup_t|\beta(t)-\alpha(t)|<\min_t|\alpha(t)-p|$, then
$(\beta-p)/(\alpha-p)$ lies in the disc $|z-1|<1$. This ratio has a
single continuous argument in $(-\pi/2,\pi/2)$ whose endpoint values
agree. Adding that argument to a lift for $\alpha-p$ gives a lift
for $\beta-p$ with the same endpoint difference. Thus sufficiently
close loops have the same index. For a continuous homotopy avoiding
$p$, compactness gives a positive minimum distance from $p$ and uniform
continuity permits a finite subdivision in the homotopy parameter into
such close pairs. This proves homotopy invariance. Varying $p$ slightly
is the same close-loop argument with $\alpha-p$; hence the index is
locally constant in $p$. If $p$ lies outside a ball containing the loop,
contract the loop linearly to the center of that ball. The contraction
stays in the ball and avoids $p$, so the index is zero.

For a $C^1$ path, each of the local argument lifts used above is $C^1$.
Differentiating the local argument gives
$\theta'=\det(\alpha-p,\alpha')/\norm{\alpha-p}^2$.
The expressions agree on overlapping argument charts. Integrating this
derivative gives $\theta(1)-\theta(0)=2\pi k$ and the formula.
Nothing in this proof asserts any classification of the complement of
the loop.
\end{proof}

\section{An embedded transverse collar}
\begin{lemma}[A $C^1$ collar for a regular embedded loop]\label{d13:21:collar}
Let $\gamma:\R/\Z\to\R^2$ be a $C^1$ injective map with
$\gamma'(t)\ne0$ for all $t$. There are a periodic $C^1$ vector
field $v(t)$ and $a>0$ for which
\[
 F:(\R/\Z)\times(-a,a)\longrightarrow\R^2,
 \qquad F(t,r)=\gamma(t)+rv(t)
\]
is a $C^1$ diffeomorphism onto an open neighborhood $N$ of the curve.
Moreover $\det(\gamma'(t),v(t))>0$ for all $t$. The two sets
$N_+=F((\R/\Z)\times(0,a))$ and
$N_-=F((\R/\Z)\times(-a,0))$ are disjoint path-connected subsets
of the curve complement.
\end{lemma}
\begin{proof}
The continuous unit left normal
$n(t)=J\gamma'(t)/\norm{\gamma'(t)}$ is periodic. It has a periodic
$C^1$ approximation $v$ with $\sup\norm{v-n}<1/2$. For completeness,
choose a nonnegative smooth bump of integral one supported in a small
interval about zero, periodize it, and convolve it with $n$ on the
circle. Uniform continuity of $n$ makes the convolution uniformly
close to $n$ as the support shrinks. The convolution is $C^1$ because
its derivative is obtained by differentiating the smooth kernel, not
$n$, under the integral. Such a bump is, after rescaling and normalizing,
$e^{-1/(1-u^2)}$ on $|u|<1$ and zero elsewhere; all its derivatives
vanish at the two endpoints by exponential decay.

Consequently
\[
 \det(\gamma',v)\geq
 \norm{\gamma'}-\norm{\gamma'}\norm{v-n}
 >\tfrac12\norm{\gamma'}.
\]
The minimum of $\norm{\gamma'}$ is positive by compactness. The
Jacobian of $F$ in any circle coordinate is
$\det(\gamma'(t)+rv'(t),v(t))$. It is bounded away from zero for
all $t$ and sufficiently small $|r|$, because $v,v'$ are bounded.
The inverse-function theorem therefore makes $F$ a local $C^1$
diffeomorphism on a thin cylinder.

It is injective on some uniformly thin cylinder. Otherwise there would
be distinct pairs $(t_j,r_j),(u_j,s_j)$ with
$|r_j|,|s_j|\to0$ and $F(t_j,r_j)=F(u_j,s_j)$.
Pass to convergent subsequences of the circle coordinates. Their
limits satisfy $\gamma(t_*)=\gamma(u_*)$, hence $t_*=u_*$ on the
circle by injectivity of $\gamma$. Both pairs then lie, for large
$j$, in a single neighborhood on which the inverse-function theorem
made $F$ injective, a contradiction. Decrease $a$ to secure both
injectivity and the nonzero Jacobian everywhere on the cylinder.
The image of a local diffeomorphism is open, and its local inverses
patch to a global $C^1$ inverse by injectivity. This proves the first
assertion. The two open half-cylinders are path connected, so their
images are path connected and disjoint. Neither image meets the curve,
because points of the curve already have the unique preimages $(t,0)$.
\end{proof}

\begin{lemma}[Local index jump at a regular arc]\label{d13:21:jump}
With the preceding orientation of the collar,
\[
 \operatorname{ind}(\gamma,p_+)
       -\operatorname{ind}(\gamma,p_-)=1
 \qquad(p_+\in N_+,\ p_-\in N_-).
\]
\end{lemma}
\begin{proof}
By local constancy of the index and path connectedness of each collar
side, it suffices to use points sufficiently close to one fixed curve
point. Translate and rotate that point to zero, preserving orientation,
and parametrize a small oriented curve arc by its first coordinate:
\[
 \gamma(u)=(u,h(u)),\quad -d\leq u\leq d,
 \qquad h(0)=h'(0)=0.
\]
This is possible because the tangent has nonzero positive first
coordinate after the rotation and choice of the increasing local
parameter. The rest of the curve is compact and stays a positive
distance from zero, by injectivity and compactness. For small
$\varepsilon>0$, the points $p_\pm=(0,\pm\varepsilon)$ do not lie on
the curve. They belong to $N_\pm$: the transverse coordinate of
$F^{-1}(0,y)$ has positive derivative at $y=0$, since
$\det(\gamma',v)>0$ and the tangent points to the right.

The contribution of the remaining curve to the integral index has the
same limit from $p_+$ and $p_-$ as $\varepsilon\downarrow0$, because
that curve portion stays away from zero. Compute the angular change
of the local graph arc. For $u>0$ use the argument branch in
$(-\pi/2,\pi/2)$ for both vectors
$u+i(h(u)\mp\varepsilon)$. For $u<0$, choose the branch in
$(-3\pi/2,-\pi/2)$ when the reference point is $p_+$, so it joins
the value $-\pi/2$ at $u=0$. For reference point $p_-$, use instead
the left-half-plane branch in $(\pi/2,3\pi/2)$, joining $\pi/2$
at zero. These choices are continuous on the respective entire arc:
the sign of the real part controls the half-plane on each side of
zero, and the vector at zero is respectively $-i\varepsilon$ or
$i\varepsilon$.

At the right endpoint the two argument values converge to the same
value as $\varepsilon\downarrow0$. At the left endpoint the branch
for $p_-$ is $2\pi$ larger than that for $p_+$ in the limit. Hence
the angular change from left to right for $p_+$ minus that for $p_-$
tends to $2\pi$. The two complete integral indices therefore differ
by a quantity tending to one. Each complete index is constant on its
collar side, so their difference is exactly one, as claimed. This
calculation uses no separation theorem.
\end{proof}

\section{The regular-curve separation theorem}
\begin{theorem}[Jordan separation for a regular $C^1$ loop]\label{d13:21:smooth-jordan}
Let $\gamma:\R/\Z\to\R^2$ be $C^1$, injective, and regular
($\gamma'\ne0$), and let $\Gamma$ be its trace. Then
$\R^2\setminus\Gamma$ has exactly two connected components:
a bounded open component $\Omega$ and an unbounded open component $E$.
They are path connected,
\[
 \partial\Omega=\partial E=\Gamma,
 \qquad \overline\Omega=\Omega\cup\Gamma,
\]
and $\Omega$ has $C^1$ boundary. The winding index is zero on $E$
and is identically either $1$ or $-1$ on $\Omega$, according to the
orientation of $\gamma$.
\end{theorem}
\begin{proof}
Take the transverse collar from Lemma~\ref{d13:21:collar}.
Every point $p\notin\Gamma$ can be connected within the complement
to one of its two sides. Indeed, compactness of $\Gamma$ supplies a
nearest point $q=\gamma(t_0)$, at positive distance
$d=\norm{p-q}$. The open segment $q+\theta(p-q)$, $0<\theta\leq1$,
contains no curve point. A point of the curve at some
$0<\theta<1$ would be at distance $(1-\theta)d<d$ from $p$,
contradicting minimality. The endpoint $p$ is not on the curve either.

Differentiating $\norm{p-\gamma(t)}^2$ at its minimizing circle
point gives $\ip{p-q}{\gamma'(t_0)}=0$. This is a nonzero normal
direction. For small $\theta$, write
$q+\theta(p-q)=F(t(\theta),r(\theta))$ using the local inverse
of the collar. Differentiating at zero gives
\[
 r'(0)=\frac{\det(\gamma'(t_0),p-q)}
                   {\det(\gamma'(t_0),v(t_0))}\ne0.
\]
Consequently $r(\theta)$ has a fixed nonzero sign for sufficiently
small positive $\theta$. The terminal portion of the segment is
therefore in $N_+$ or $N_-$. It connects $p$ to that side without
meeting $\Gamma$.

It follows that the complement has at most two path components: every
point joins one of the two path-connected collar sides. An open subset
of the plane has open path components and connected components equal
to its path components. One elementary verification is to take the
set of points reachable from a given point by polygonal paths in the
open set; small open balls show that this set and its relative
complement are open. Thus it is the connected component.
Lemma~\ref{d13:21:jump} shows that the two collar sides have different
indices. Local constancy of the index prevents a path, or a connected
component, from joining them. Hence there are exactly two components.

Choose a large closed ball containing $\Gamma$ in its interior. Its
exterior is connected and belongs to one component, say $E$, on
which the index is zero by Lemma~\ref{d13:21:index-properties}.
Any unbounded component must meet this ball exterior. Therefore the
other component $\Omega$ is bounded, and $E$ is the unique unbounded
one. The jump of one between the two collar sides forces the index
on $\Omega$ to be $1$ or $-1$.

Every point of $\Gamma$ is approached from both collar sides, hence
belongs to the boundary of both components. Conversely a boundary
point of a component cannot lie in the open complement: a small
connected ball in the complement belongs to just one component and
either lies inside that component or misses it. Thus both boundaries
are exactly $\Gamma$, and the stated closure identity follows.
Finally, in each collar chart the curve is $r=0$ and its two sides
are $r>0$ and $r<0$. Equivalently, projection onto the tangent
direction makes the curve a $C^1$ graph by the inverse-function
theorem, with $\Omega$ on exactly one side. This is the usual local
graph definition of a $C^1$ boundary.
\end{proof}

\begin{corollary}[Interior containment in a simply connected domain]\label{d13:21:domain-interior}
Let $U\subset\R^2$ be open and simply connected, meaning path connected
with every continuous loop contractible within $U$. If the trace of a
regular $C^1$ embedded loop $\gamma$ lies in $U$, then its bounded
interior $\Omega$ and its closure both lie in $U$.
\end{corollary}
\begin{proof}
Suppose $p\in\Omega$ but $p\notin U$. A contraction of $\gamma$
to a constant loop within $U$ is a continuous homotopy avoiding $p$.
Its winding index about $p$ is therefore the index of the constant
loop, namely zero, by Lemma~\ref{d13:21:index-properties}.
Theorem~\ref{d13:21:smooth-jordan} makes that same index $1$ or $-1$,
a contradiction. Thus $\Omega\subset U$. Its boundary $\Gamma$
already lies in $U$, and the closure identity of that theorem gives
$\overline\Omega\subset U$ as well.
\end{proof}

\paragraph{Sources and coverage.}
This supplement supplies the regular-curve separation ingredient used
in the planar periodic-orbit and Bendixson arguments. The elementary
argument-lifting and winding formulas agree with the earlier complex
analysis foundation, and the collar uses the multivariable
inverse-function theorem and compactness. The nearest-point and local
index-jump arguments above give the separation proof explicitly. The
general continuous Jordan curve theorem, a Schoenflies theorem, and
classification of wild or nonsmooth embeddings are not invoked.

% END chapters/M13.21

\backmatter

% BEGIN source-register
\chapter*{Source and scope register}
\addcontentsline{toc}{chapter}{Source and scope register}
The source inventory records the exact archived files, versions, page counts
and SHA-256 identities. Chapter source paragraphs identify the actual source
sections used. The accompanying coverage table distinguishes adopted
mathematics from more advanced material retained only as future references.
Proofs here are independently written syntheses, not facsimile chapters.

\begin{description}[style=nextline,leftmargin=12mm]
\item[T.] Gerald Teschl, \emph{Ordinary Differential Equations and Dynamical
Systems}, author's preliminary version of the AMS Graduate Studies in
Mathematics volume 140. Main spine for classical theory, regular spectral
problems and qualitative dynamics.
\item[CL.] Earl A. Coddington and Norman Levinson, \emph{Theory of Ordinary
Differential Equations}, 1955 text, Tata McGraw--Hill reprint. Classical
existence, parameter dependence, linear and boundary-value theory.
\item[P.] Lawrence Perko, \emph{Differential Equations and Dynamical Systems},
third edition, Springer Texts in Applied Mathematics 7. Local, global and
bifurcation theory.
\item[GH.] John Guckenheimer and Philip Holmes, \emph{Nonlinear Oscillations,
Dynamical Systems, and Bifurcations of Vector Fields}, Springer Applied
Mathematical Sciences 42. Geometric nonlinear dynamics and selected
bifurcation constructions.
\item[W.] Peter Walters, \emph{An Introduction to Ergodic Theory},
Springer, GTM 79 (1982; softcover printing 2000).
Introductory measure-preserving and topological dynamics.
\item[HSD.] Morris W. Hirsch, Stephen Smale and Robert L. Devaney,
\emph{Differential Equations, Dynamical Systems, and an Introduction to
Chaos}, third edition, Academic Press, copyright 2013. Undergraduate
examples, geometric explanations and discrete dynamics.
\end{description}

The material here is fixed for this preparation. A subsequent Lean project
must read all adopted mathematical statements independently; neither a
source ledger nor a favorable preparation review is a substitute for the
normal machine checks and independent fidelity reviews. The physical PDF
page count is navigation metadata. Any granularity reference concerns the
remaining mathematical content and does not force a goal count.

% END source-register

\end{document}
